% L’intégrité morale — ECM 1re Tech — Cours signé OnBuch+
% Programme officiel MINESEC. Compiler avec : tectonic N02-integrite-morale.tex
\newcommand{\DOCMATIERE}{ECM}
\newcommand{\DOCNIVEAU}{1re Tech}
\newcommand{\DOCMODULE}{ECM – classe de Première technique}
\newcommand{\DOCLECON}{N02}
\newcommand{\DOCTITRE}{L’intégrité morale}
% OnBuch+ — préambule « cours signés » ENRICHI (compilé avec tectonic / XeLaTeX)
% Système visuel « Pop » d'OnBuch+ : fond crème (jamais blanc), bordures encre
% épaisses, ombres « dures » décalées, titres Archivo Black en capitales.
% Version MATHÉMATIQUES : graphes (pgfplots/tikz), boîtes pédagogiques
% (définition, propriété, méthode, attention, le savais-tu, exercice résolu,
% exercice, TP, bilan), page de garde et signature OnBuch+ ancrée en bas.
%
% Macros attendues dans main.tex AVANT \input{preamble} :
%   \DOCMATIERE  \DOCNIVEAU  \DOCMODULE  \DOCLECON (numéro)  \DOCTITRE
\documentclass[11pt]{article}

\usepackage{fontspec}
\usepackage[french]{babel}
\usepackage[a4paper,margin=2cm,top=3.4cm,bottom=2.8cm,headheight=1.4cm,footskip=1.3cm]{geometry}
\usepackage{amsmath,amssymb,amsfonts}
\usepackage{mathtools} % \xmapsto, \coloneqq... souvent utilisés par les modèles
\usepackage{cancel} % simplifications barrées (bilans de forces)
% \abs{z} (module, valeur absolue) et \norm{u} (norme), utilisés par les modèles
\providecommand{\abs}{}\let\abs\relax\DeclarePairedDelimiter{\abs}{\lvert}{\rvert}
\providecommand{\norm}{}\let\norm\relax\DeclarePairedDelimiter{\norm}{\lVert}{\rVert}
\usepackage[table]{xcolor} % [table] : fournit aussi \rowcolors (alternance de lignes)
\usepackage{pagecolor}
\usepackage{tikz}
\usetikzlibrary{arrows.meta,positioning,calc,shapes.geometric,shapes.misc,shapes.arrows,decorations.pathmorphing,decorations.pathreplacing,decorations.markings,patterns,shadows,mindmap,trees,fit,backgrounds,matrix,intersections,angles,quotes}
\usepackage{pgfplots}
\usepackage[version=4]{mhchem} % équations nucléaires \ce{^{235}_{92}U}
\usepackage[european,siunitx]{circuitikz} % schémas électriques
\pgfplotsset{compat=1.18}
\usepgfplotslibrary{fillbetween,groupplots} % aires entre courbes (calcul intégral)
\usepackage{enumitem}
\usepackage{fancyhdr}
% [most] charge toutes les bibliothèques tcolorbox (listings, minted, raster,
% documentation...) et gonfle inutilement la mémoire TeX sur un document long
% (~300 boîtes) : on ne charge que ce qui est réellement utilisé.
\usepackage[skins,breakable]{tcolorbox}
\usepackage{graphicx}
\usepackage{colortbl}
\usepackage{booktabs}
\usepackage{tabularx}
\usepackage{multirow}
\usepackage{diagbox} % case d'en-tête diagonale des tableaux à double entrée
% Colonnes extensibles centrée / gauche / droite (C, L, R), souvent utilisées par les modèles
\newcolumntype{C}{>{\centering\arraybackslash}X}
\newcolumntype{L}{>{\raggedright\arraybackslash}X}
\newcolumntype{R}{>{\raggedleft\arraybackslash}X}
\usepackage{array}
\usepackage{multicol}
\usepackage{siunitx}
% parse-numbers=false : accepte aussi \SI{2,5 \times 10^{-3}}{...}
\sisetup{locale=FR,output-decimal-marker={,},group-minimum-digits=4,parse-numbers=false}
\DeclareSIUnit\ohm{\ensuremath{\symup{\Omega}}} % U+2126 absent de la police
\DeclareSIUnit\division{div} % réglages d'oscilloscope (ms/div, V/div)
\DeclareSIUnit\tr{tr} % tours (vitesse de rotation, ex. \SI{1500}{\tr\per\minute})
\DeclareSIUnit\molar{M} % concentration molaire (ex. \SI{3}{\molar}, \SI{1}{\micro\molar})
\DeclareSIUnit\an{an} % année (le modèle confond parfois avec \year, un compteur TeX) — ex. \SI{5}{\kilo\meter\per\an}
\DeclareSIUnit\FCFA{FCFA} % franc CFA (ex. \SI{500}{\FCFA\per\kilo\gram})
\DeclareSIUnit\degreeBrix{\textdegree Brix} % teneur en sucre d'un sirop/jus (réfractométrie)
\DeclareSIUnit\month{mois} % le modèle utilise parfois \month, un compteur TeX, comme unité de durée
\DeclareSIUnit\microliter{\micro\liter} % le modèle écrit parfois \microliter en un seul token
\DeclareSIUnit\million{millions} % dénombrement (ex. \SI{15}{\million\per\mL} spermatozoïdes)
\usepackage{qrcode}
% \degree : commande fréquemment utilisée par les modèles pour un angle ou une
% température en dehors de \SI{}{} (non fournie par les paquets chargés).
\providecommand{\degree}{^{\circ}}
% \qed (fin de démonstration), utilisé par les modèles sans amsthm
\providecommand{\qed}{\hfill$\square$}
% \barre : double barre des valeurs interdites dans un tableau de variations (tkz-tab)
\providecommand{\barre}{\ensuremath{\|}}
% environnement proof (amsthm non chargé) utilisé par les modèles
\ifdefined\proof\else\newenvironment{proof}[1][Démonstration]{\par\noindent\textit{#1.}\ }{\qed\par}\fi
\usepackage{lastpage}
\usepackage{hyperref}
\hypersetup{hidelinks}

% ============================================================
% Palette « Pop » OnBuch+ — mêmes hex que la classe P (plus_ui.dart)
% ============================================================
\definecolor{poporange}{HTML}{F59321}
\definecolor{popdark}{HTML}{E07A0C}
\definecolor{popcream}{HTML}{FFF6E8}
\definecolor{popink}{HTML}{1C1714}
\definecolor{poppurple}{HTML}{7A5AE0}
\definecolor{popgreen}{HTML}{2E9E6A}
\definecolor{poppink}{HTML}{E0457B}
\definecolor{popblue}{HTML}{2F7DE1}
\definecolor{popgold}{HTML}{C98A12}
\definecolor{popmuted}{HTML}{8A7F76}
\definecolor{popline}{HTML}{EADFD2}
\definecolor{popgray}{HTML}{8A7F76}
\colorlet{popgrey}{popgray}
% Alias tolérants (le modèle nomme parfois les couleurs différemment)
\colorlet{popwhite}{white}
\colorlet{popblack}{popink}
\colorlet{popred}{poppink}
\colorlet{popyellow}{popgold}
% Teintes claires (fonds de tableaux/encadrés) — le modèle nomme parfois
% les couleurs avec un suffixe L (ex. popblueL) sans les avoir définies.
\colorlet{poporangeL}{poporange!20}
\colorlet{popdarkL}{popdark!20}
\colorlet{poppurpleL}{poppurple!20}
\colorlet{popgreenL}{popgreen!20}
\colorlet{poppinkL}{poppink!20}
\colorlet{popblueL}{popblue!20}
\colorlet{popgoldL}{popgold!20}
\colorlet{popredL}{popred!20}
\colorlet{popgrayL}{popgray!20}
% Teintes claires pour fonds de tableaux / zones
\colorlet{popblueL}{popblue!10!white}
\colorlet{poporangeL}{poporange!14!white}
\colorlet{popgreenL}{popgreen!12!white}
\colorlet{poppurpleL}{poppurple!12!white}
\colorlet{poppinkL}{poppink!12!white}
\colorlet{popgoldL}{popgold!14!white}

\pagecolor{popcream}

% ============================================================
% Typographie
% ============================================================
\defaultfontfeatures{Ligatures=TeX}
\setmainfont{PlusJakartaSans-Medium.ttf}[Path=../../../fonts/,BoldFont=PlusJakartaSans-Bold.ttf]
% Maths en sans-serif (Fira Math) avec chiffres et lettres droites de Plus
% Jakarta, pour que les formules se fondent dans le texte.
\usepackage[math-style=ISO,bold-style=ISO]{unicode-math}
\setmathfont{FiraMath-Regular.otf}[Path=../../../fonts/]
\setmathfont{PlusJakartaSans-Medium.ttf}[Path=../../../fonts/,range={up/num,up/{Latin,latin},\mathup}]
\usepackage{icomma} % virgule décimale sans espace en mode math
% Caractères mathématiques Unicode absents des polices de texte (Plus Jakarta,
% Archivo Black) : rendus par leur équivalent mathématique, en texte comme en
% formule — sinon ils s'affichent en carré vide (ex. « Arithmétique dans ℤ »).
\usepackage{newunicodechar}
\newunicodechar{ℝ}{\ensuremath{\mathbb{R}}}
\newunicodechar{ℤ}{\ensuremath{\mathbb{Z}}}
\newunicodechar{ℂ}{\ensuremath{\mathbb{C}}}
\newunicodechar{ℕ}{\ensuremath{\mathbb{N}}}
\newunicodechar{ℚ}{\ensuremath{\mathbb{Q}}}
\newunicodechar{𝔻}{\ensuremath{\mathbb{D}}}
\newunicodechar{α}{\ensuremath{\alpha}}
\newunicodechar{β}{\ensuremath{\beta}}
\newunicodechar{γ}{\ensuremath{\gamma}}
\newunicodechar{δ}{\ensuremath{\delta}}
\newunicodechar{ε}{\ensuremath{\varepsilon}}
\newunicodechar{θ}{\ensuremath{\theta}}
\newunicodechar{λ}{\ensuremath{\lambda}}
\newunicodechar{μ}{\ensuremath{\mu}}
\newunicodechar{σ}{\ensuremath{\sigma}}
\newunicodechar{τ}{\ensuremath{\tau}}
\newunicodechar{φ}{\ensuremath{\varphi}}
\newunicodechar{ψ}{\ensuremath{\psi}}
\newunicodechar{ω}{\ensuremath{\omega}}
\newunicodechar{Σ}{\ensuremath{\Sigma}}
% majuscules produites par \MakeUppercase dans les titres (α-aminés -> Α)
\newunicodechar{Α}{\ensuremath{\alpha}}
\newunicodechar{Β}{\ensuremath{\beta}}
\newunicodechar{Γ}{\ensuremath{\Gamma}}
\newunicodechar{Δ}{\ensuremath{\Delta}}
\newunicodechar{Ε}{\ensuremath{\varepsilon}}
\newunicodechar{Θ}{\ensuremath{\Theta}}
\newunicodechar{Λ}{\ensuremath{\Lambda}}
\newunicodechar{Μ}{\ensuremath{\mu}}
\newunicodechar{Τ}{\ensuremath{\tau}}
\newunicodechar{Φ}{\ensuremath{\Phi}}
\newunicodechar{Ψ}{\ensuremath{\Psi}}
\newunicodechar{Ω}{\ensuremath{\Omega}}
\newunicodechar{↦}{\ensuremath{\mapsto}}
\newunicodechar{∈}{\ensuremath{\in}}
\newunicodechar{∉}{\ensuremath{\notin}}
\newunicodechar{∧}{\ensuremath{\wedge}}
\newunicodechar{⇔}{\ensuremath{\Leftrightarrow}}
\newunicodechar{⟺}{\ensuremath{\Longleftrightarrow}}
\newunicodechar{⇒}{\ensuremath{\Rightarrow}}
\newunicodechar{⟹}{\ensuremath{\Longrightarrow}}
\newunicodechar{∘}{\ensuremath{\circ}}
\newunicodechar{∪}{\ensuremath{\cup}}
\newunicodechar{∩}{\ensuremath{\cap}}
\newunicodechar{≡}{\ensuremath{\equiv}}
\newunicodechar{⊂}{\ensuremath{\subset}}
\newunicodechar{⊥}{\ensuremath{\perp}}
\newunicodechar{∃}{\ensuremath{\exists}}
\newunicodechar{∀}{\ensuremath{\forall}}
\newunicodechar{∛}{\ensuremath{\sqrt[3]{\,}}}
\newunicodechar{∜}{\ensuremath{\sqrt[4]{\,}}}
\newunicodechar{⃗}{\ensuremath{{}^{\to}}}
\newunicodechar{ⁿ}{\ifmmode^{n}\else\textsuperscript{\ensuremath{n}}\fi}
\newunicodechar{ˣ}{\ifmmode^{x}\else\textsuperscript{\ensuremath{x}}\fi}
\newunicodechar{ᵇ}{\ifmmode^{b}\else\textsuperscript{\ensuremath{b}}\fi}
\newunicodechar{ʳ}{\ifmmode^{r}\else\textsuperscript{\ensuremath{r}}\fi}
\newunicodechar{ᵘ}{\ifmmode^{u}\else\textsuperscript{\ensuremath{u}}\fi}
\newunicodechar{ʸ}{\ifmmode^{y}\else\textsuperscript{\ensuremath{y}}\fi}
\newunicodechar{ᵃ}{\ifmmode^{a}\else\textsuperscript{\ensuremath{a}}\fi}
\newunicodechar{ᵉ}{\ifmmode^{e}\else\textsuperscript{\ensuremath{e}}\fi}
\newunicodechar{ᵏ}{\ifmmode^{k}\else\textsuperscript{\ensuremath{k}}\fi}
\newunicodechar{ᵖ}{\ifmmode^{p}\else\textsuperscript{\ensuremath{p}}\fi}
\newunicodechar{ᴺ}{\ifmmode^{N}\else\textsuperscript{\ensuremath{N}}\fi}
\newunicodechar{ᶜ}{\ifmmode^{c}\else\textsuperscript{\ensuremath{c}}\fi}
\newunicodechar{ʰ}{\ifmmode^{h}\else\textsuperscript{\ensuremath{h}}\fi}
\newunicodechar{ᵠ}{\ifmmode^{q}\else\textsuperscript{\ensuremath{q}}\fi}
\newunicodechar{ₙ}{\ifmmode_{n}\else\textsubscript{\ensuremath{n}}\fi}
\newunicodechar{ₐ}{\ifmmode_{a}\else\textsubscript{\ensuremath{a}}\fi}
\newunicodechar{ᵢ}{\ifmmode_{i}\else\textsubscript{\ensuremath{i}}\fi}
\newunicodechar{ᵣ}{\ifmmode_{r}\else\textsubscript{\ensuremath{r}}\fi}
\newunicodechar{ₖ}{\ifmmode_{k}\else\textsubscript{\ensuremath{k}}\fi}
\newunicodechar{ᵦ}{\ifmmode_{\beta}\else\textsubscript{\ensuremath{\beta}}\fi}
\newunicodechar{ₓ}{\ifmmode_{x}\else\textsubscript{\ensuremath{x}}\fi}
\newfontfamily\popdisplay{ArchivoBlack-Regular.ttf}[Path=../../../fonts/]
\newfontfamily\poplogo{SpaceGrotesk-Bold.ttf}[Path=../../../fonts/]
\newfontfamily\popsemi{PlusJakartaSans-SemiBold.ttf}[Path=../../../fonts/]
\newfontfamily\popxbold{PlusJakartaSans-ExtraBold.ttf}[Path=../../../fonts/]
\newfontfamily\popxboldls{PlusJakartaSans-ExtraBold.ttf}[Path=../../../fonts/,LetterSpace=3.0]

\setlist[enumerate]{leftmargin=1.7em,itemsep=5pt,topsep=4pt}
\setlist[itemize]{leftmargin=1.7em,itemsep=4pt,topsep=4pt,label={\color{poporange}\textbullet}}
\setlength{\parindent}{0pt}
\setlength{\parskip}{5pt}
\renewcommand{\arraystretch}{1.25}

% pgfplots : style Pop par défaut
\pgfplotsset{
  every axis/.append style={axis line style={popink,line width=1pt},tick style={popink},
    grid style={popline,line width=0.5pt},label style={font=\small},tick label style={font=\footnotesize}},
  popcurve/.style={poporange,line width=1.8pt,no markers,smooth},
}
% Même style disponible dans un \draw TikZ simple (hors pgfplots)
\tikzset{popcurve/.style={poporange,line width=1.8pt}}
\tikzset{popgrid/.style={popline,very thin}}
\colorlet{popcurve}{poporange} % le nom du style est parfois utilisé comme couleur

% ============================================================
% Pill / squiggle
% ============================================================
\newcommand{\poppill}[3]{%
  \tikz[baseline=-0.55ex]\node[draw=popink,line width=1.2pt,rounded corners=2.6pt,%
    fill=#2,inner xsep=6.5pt,inner ysep=3pt]{%
    \popxboldls\fontsize{7}{7}\selectfont\color{#3}\MakeUppercase{#1}};%
}
\newcommand{\popsquiggle}[1]{%
  \tikz[baseline=-0.4ex]{
    \draw[#1,line width=2.6pt,line cap=round]
      plot [smooth] coordinates {(0,0.09) (0.34,0.01) (0.68,0.21) (1.02,0.01) (1.36,0.10)};
  }%
}
% Mot-clé surligné (marqueur orange sous le texte)
\newcommand{\cle}[1]{{\popxbold\color{popdark}#1}}

% ============================================================
% En-tête / pied
% ============================================================
\pagestyle{fancy}
\fancyhf{}
\renewcommand{\headrulewidth}{0pt}
\renewcommand{\footrulewidth}{0pt}
\lhead{\raisebox{-0.15ex}{\poplogo\fontsize{13}{13}\selectfont\color{popink}OnBuch\color{poporange}+}}
\rhead{\poppill{\DOCMATIERE\ \textbullet\ \DOCNIVEAU\ \textbullet\ Leçon \DOCLECON}{white}{popink}}
\lfoot{\popxbold\fontsize{7.6}{7.6}\selectfont\color{popmuted}\MakeUppercase{Cours signé OnBuch+}\ \textbullet\ {\popsemi\DOCTITRE}}
\rfoot{\tikz[baseline=-0.6ex]\node[rounded corners=8.5pt,fill=popink,minimum height=17pt,minimum width=17pt,inner xsep=5pt,inner sep=0]{\popxbold\fontsize{8}{8}\selectfont\color{popcream}\thepage\,/\,\pageref*{LastPage}};}

% ============================================================
% Titres
% ============================================================
\newcommand{\chaptitre}[2]{%
  \begin{tcolorbox}[enhanced,colback=poporange,colframe=popink,boxrule=2.6pt,arc=11pt,
      drop shadow=popink,left=15pt,right=15pt,top=12pt,bottom=11pt,before skip=3pt,after skip=18pt]
    \poppill{#2}{popink}{popcream}\par\vspace{9pt}
    {\raggedright\hyphenpenalty=10000\exhyphenpenalty=10000\popdisplay\fontsize{22}{24}\selectfont\color{popcream}\MakeUppercase{#1}\par}
  \end{tcolorbox}
}
% Section numérotée : pastille orange avec le numéro + titre Archivo
\newcounter{popsec}
\newcommand{\coursec}[1]{%
  \stepcounter{popsec}\setcounter{popsub}{0}%
  \par\vspace{16pt}\noindent
  \begin{minipage}{\linewidth}
  \tikz[baseline=-0.7ex]\node[draw=popink,line width=1.6pt,fill=poporange,rounded corners=4pt,minimum size=22pt,inner sep=0]{\popdisplay\fontsize{12}{12}\selectfont\color{popcream}\thepopsec};%
  \hspace{8pt}{\popdisplay\fontsize{15}{17}\selectfont\color{popink}\MakeUppercase{#1}}\par
  \vspace{4pt}\noindent{\color{poporange}\rule{36pt}{3.6pt}}
  \end{minipage}\par\nopagebreak\vspace{8pt}
}
\newcounter{popsub}
\newcommand{\courssub}[1]{%
  \stepcounter{popsub}%
  \par\vspace{9pt}\noindent{\popxbold\fontsize{11.5}{13}\selectfont\color{popink}{\color{poporange}\thepopsec.\thepopsub}\hspace{6pt}#1}\par\nopagebreak\vspace{4pt}
}

% ============================================================
% Boîtes pédagogiques — même recette : fond blanc, bordure épaisse colorée,
% ombre dure encre, badge plein. Titre optionnel [#1].
% ============================================================
\tcbset{popbase/.style={breakable,enhanced,colback=white,boxrule=2.3pt,arc=9pt,
  shadow={3pt}{-3pt}{0pt}{popink},left=14pt,right=14pt,top=11pt,bottom=11pt,before skip=11pt,after skip=15pt,
  fontupper=\popsemi\fontsize{10.6}{14.5}\selectfont\color{popink},
  fontlower=\popsemi\fontsize{10.6}{14.5}\selectfont\color{popink}}}
% Coche dessinée (le glyphe ✓ n'existe pas dans les polices du document)
\newcommand{\popcheck}[1]{\tikz[baseline=-0.5ex]\draw[#1,line width=1.6pt,line cap=round,line join=round](-0.13,0.0)--(-0.04,-0.09)--(0.13,0.1);}
\providecommand{\coche}{\popcheck{popgreen}}
\providecommand{\resultat}[1]{\par\smallskip\noindent\fcolorbox{popgreen}{popgreenL}{\parbox{\dimexpr\linewidth-2\fboxsep-2\fboxrule\relax}{\textbf{Résultat.} #1}}\par\smallskip}
\newcommand{\popboxhead}[3]{\poppill{#1}{#2}{white}\ifx\relax#3\relax\else\hspace{6pt}{\popxbold\fontsize{10.6}{12}\selectfont\color{popink}#3}\fi\par\vspace{7pt}}

\newtcolorbox{aretenir}[1][]{popbase,colframe=popgreen,before upper={\popboxhead{À retenir}{popgreen}{#1}}}
\newtcolorbox{exemplebox}[1][]{popbase,colframe=poporange,before upper={\popboxhead{Exemple}{poporange}{#1}}}
\newtcolorbox{definition}[1][]{popbase,colframe=popblue,before upper={\popboxhead{Définition}{popblue}{#1}}}
\newtcolorbox{propriete}[1][]{popbase,colframe=poppurple,before upper={\popboxhead{Propriété}{poppurple}{#1}}}
\newtcolorbox{methode}[1][]{popbase,colframe=popgold,before upper={\popboxhead{Méthode}{popgold}{#1}}}
\newtcolorbox{attention}[1][]{popbase,colframe=poppink,before upper={\popboxhead{Attention}{poppink}{#1}}}
\newtcolorbox{experience}[1][]{popbase,colframe=popblue,colback=popblueL,before upper={\popboxhead{Expérience}{popblue}{#1}}}
% « Le savais-tu ? » : carte violette pleine (contexte, culture, Cameroun)
\newtcolorbox{savaistu}[1][]{popbase,colback=poppurple,colframe=popink,
  fontupper=\popsemi\fontsize{10.4}{14.2}\selectfont\color{white},
  before upper={\poppill{Le savais-tu\,?}{white}{poppurple}\ifx\relax#1\relax\else\hspace{6pt}{\popxbold\color{white}#1}\fi\par\vspace{7pt}}}
% Exercice résolu : énoncé en haut, \tcblower puis la solution sur fond crème
\newcounter{popexo}\newcounter{popexr}
\newtcolorbox{exoresolu}[1][]{popbase,colframe=poporange,colbacklower=popcream,
  segmentation style={popink,line width=1.2pt,dashed},
  before upper={\stepcounter{popexr}\popboxhead{Exercice résolu \thepopexr}{poporange}{#1}},
  before lower={\poppill{Solution}{popink}{popcream}\par\vspace{6pt}}}
% Exercice (non résolu en place) — la correction est dans la partie Corrigés
\newtcolorbox{exercice}[1][]{popbase,colframe=popink,
  before upper={\stepcounter{popexo}\popboxhead{Exercice \thepopexo}{popink}{#1}}}
% Corrigé
\newtcolorbox{corrige}[1][]{popbase,colframe=popgreen,colback=popgreenL,
  before upper={\popboxhead{Corrigé}{popgreen}{#1}}}
% Objectifs / prérequis / bilan
\newtcolorbox{objectifs}{popbase,colframe=popink,colback=poporangeL,
  before upper={\popboxhead{Objectifs de la leçon}{popink}{}}}
\newtcolorbox{prerequis}{popbase,colframe=popink,colback=popblueL,
  before upper={\popboxhead{Prérequis}{popblue}{}}}
\newtcolorbox{bilan}{popbase,colframe=popink,colback=white,boxrule=2.8pt,
  before upper={\popboxhead{Fiche bilan}{poporange}{L'essentiel à connaître pour l'examen}}}
% Figure encadrée avec légende
\newtcolorbox{popfigure}[1][]{enhanced,breakable=false,colback=white,colframe=popline,boxrule=1.4pt,arc=8pt,
  left=10pt,right=10pt,top=10pt,bottom=8pt,before skip=10pt,after skip=12pt,halign=center,
  lower separated=false,halign lower=center,
  fontlower=\popsemi\fontsize{9}{11.5}\selectfont\color{popmuted},
  #1}
\newcommand{\legende}[1]{\par\vspace{6pt}{\popsemi\fontsize{9}{11.5}\selectfont\color{popmuted}#1}}

% ============================================================
% Signature OnBuch+
% ============================================================
\newcommand{\poplogoinline}[1]{{\poplogo\fontsize{#1}{#1}\selectfont\color{popink}OnBuch\color{poporange}+}}
\newcommand{\signatureonbuch}{%
  \begin{tcolorbox}[enhanced,colback=popink,colframe=popink,boxrule=2.2pt,arc=10pt,
      drop shadow=poporange,left=14pt,right=14pt,top=10pt,bottom=10pt,before skip=0pt,after skip=0pt]
    \tikz[baseline=-0.7ex]\node[circle,fill=popgreen,draw=popcream,line width=1.3pt,minimum size=22pt,inner sep=0]{\popcheck{white}};%
    \hspace{9pt}\begin{minipage}[c]{0.62\linewidth}
      {\popxbold\fontsize{12}{13}\selectfont\color{popcream}Signé\ \textbullet\ {\poplogo OnBuch\color{poporange}+}}\par\vspace{2pt}
      {\raggedright\popsemi\fontsize{8}{10.5}\selectfont\color{popline}\MakeUppercase{Équipe pédagogique OnBuch+}\\\MakeUppercase{Conforme au programme officiel MINESEC}\par}
    \end{minipage}\hfill
    {\poplogo\fontsize{22}{22}\selectfont\color{popcream}OnBuch\color{poporange}+}
  \end{tcolorbox}%
}

% ============================================================
% Page de garde
% #1 titre · #2 matière · #3 niveau · #4 module · #5 n° leçon · #6 durée ·
% #7 description / objectifs (texte ou itemize)
% ============================================================
\newcommand{\pagegarde}[7]{%
  \thispagestyle{empty}
  \newgeometry{a4paper,margin=1.7cm,top=1.6cm,bottom=1.4cm}%
  \begin{tikzpicture}[remember picture,overlay]
    \fill[poporange!18] ([xshift=-3.2cm,yshift=-3.6cm]current page.north east) circle (4.6cm);
    \fill[poppurple!14] ([xshift=2.4cm,yshift=7.6cm]current page.south west) circle (3.8cm);
  \end{tikzpicture}%
  {\poplogo\fontsize{30}{30}\selectfont\color{popink}OnBuch\color{poporange}+}\hfill
  \raisebox{0.6ex}{\poppill{Cours signé \textbullet\ Premium}{popink}{popcream}}\par
  \vspace{1pt}\popsquiggle{poppurple}\par
  \vspace{8pt}
  \begin{tcolorbox}[enhanced,colback=poporange,colframe=popink,boxrule=2.8pt,arc=13pt,
      drop shadow=popink,left=18pt,right=18pt,top=12pt,bottom=13pt,after skip=10pt]
    \poppill{#2 \textbullet\ #3}{popink}{popcream}\hspace{5pt}\poppill{#4}{white}{popink}\par\vspace{8pt}
    {\popdisplay\fontsize{14}{14}\selectfont\color{popink}LEÇON #5}\par\vspace{4pt}
    {\raggedright\hyphenpenalty=10000\exhyphenpenalty=10000\popdisplay\fontsize{25}{27}\selectfont\color{popcream}\MakeUppercase{#1}\par}
    \vspace{8pt}
    {\popxbold\fontsize{9.5}{9.5}\selectfont\color{popink}\MakeUppercase{Durée conseillée : #6}}
  \end{tcolorbox}
  \begin{tcolorbox}[enhanced,colback=white,colframe=popink,boxrule=2.2pt,arc=10pt,
      drop shadow=popink,left=16pt,right=16pt,top=10pt,bottom=10pt,after skip=8pt]
    \poppill{Dans cette leçon}{poppurple}{white}\par\vspace{5pt}
    \popsemi\fontsize{9.3}{12.6}\selectfont\color{popink}#7
  \end{tcolorbox}
  \noindent\begin{minipage}[t]{0.487\linewidth}
    \begin{tcolorbox}[enhanced,colback=popgreen,colframe=popink,boxrule=2.2pt,arc=10pt,
        drop shadow=popink,left=11pt,right=11pt,top=10pt,bottom=10pt,height=3.1cm,valign=center]
      \tcbox[colback=white,colframe=popink,boxrule=1.3pt,arc=5pt,boxsep=0pt,nobeforeafter,
        left=3pt,right=3pt,top=3pt,bottom=3pt,valign=center]{\qrcode[height=1.9cm]{https://play.google.com/store/apps/details?id=com.luvvix.onbuch&hl=fr}}%
      \hspace{8pt}\begin{minipage}[c]{0.5\linewidth}
        {\popdisplay\fontsize{11}{12.5}\selectfont\color{white}\MakeUppercase{Télécharge OnBuch}}\par\vspace{4pt}
        {\raggedright\popsemi\fontsize{8.4}{11}\selectfont\color{white}Gratuit \textbullet\ Google Play\\Tuteur IA, annales, concours, résultats\par}
      \end{minipage}
    \end{tcolorbox}
  \end{minipage}\hfill
  \begin{minipage}[t]{0.487\linewidth}
    \begin{tcolorbox}[enhanced,colback=poppurple,colframe=popink,boxrule=2.2pt,arc=10pt,
        drop shadow=popink,left=11pt,right=11pt,top=10pt,bottom=10pt,height=3.1cm,valign=center]
      \tikz[baseline=-0.7ex]\node[circle,fill=white,draw=popink,line width=1.3pt,minimum size=1.6cm,inner sep=0]{\popdisplay\fontsize{24}{24}\selectfont\color{poppurple}+};%
      \hspace{8pt}\begin{minipage}[c]{0.6\linewidth}
        {\popdisplay\fontsize{11}{12.5}\selectfont\color{white}\MakeUppercase{Passe à OnBuch+}}\par\vspace{4pt}
        {\raggedright\popsemi\fontsize{8.4}{11}\selectfont\color{white}Tous les cours signés, Tuteur IA illimité, fiches PDF — onglet OnBuch+ de l'app\par}
      \end{minipage}
    \end{tcolorbox}
  \end{minipage}
  \vspace{8pt}
  \popcontenu
  \vfill
  \signatureonbuch
  \restoregeometry
  \newpage
}

% Rangée « Ce cours contient » (page de garde)
\newcommand{\poptile}[3]{% #1 couleur #2 gros texte #3 légende
  \begin{tcolorbox}[enhanced,colback=#1,colframe=popink,boxrule=2pt,arc=9pt,shadow={2.5pt}{-2.5pt}{0pt}{popink},
      left=6pt,right=6pt,top=9pt,bottom=9pt,halign=center,valign=center,height=1.75cm,width=\linewidth,nobeforeafter]
    {\popdisplay\fontsize{12.5}{13}\selectfont\color{white}\MakeUppercase{#2}}\par\vspace{3pt}
    {\centering\popsemi\fontsize{7.8}{9.5}\selectfont\color{white}#3\par}
  \end{tcolorbox}}
\newcommand{\popcontenu}{%
  \par\noindent{\popxboldls\fontsize{7.5}{7.5}\selectfont\color{popmuted}CE COURS SIGNÉ CONTIENT}\par\vspace{7pt}
  \noindent\begin{minipage}[t]{0.235\linewidth}\poptile{popblue}{Cours}{complet, illustré, pas à pas}\end{minipage}\hfill
  \begin{minipage}[t]{0.235\linewidth}\poptile{poporange}{Méthodes}{et exercices résolus}\end{minipage}\hfill
  \begin{minipage}[t]{0.235\linewidth}\poptile{poppink}{Exercices}{type examen + corrigés}\end{minipage}\hfill
  \begin{minipage}[t]{0.235\linewidth}\poptile{popgreen}{Fiche bilan}{l'essentiel à retenir}\end{minipage}\par}

% Sommaire
\newtcolorbox{sommaire}{enhanced,colback=white,colframe=popink,boxrule=2.2pt,arc=10pt,
  drop shadow=popink,left=18pt,right=18pt,top=13pt,bottom=13pt,before skip=6pt,after skip=20pt,
  before upper={\poppill{Sommaire}{poppurple}{white}\par\vspace{9pt}\popsemi\fontsize{10.6}{14.5}\selectfont\color{popink}}}

% Signature de fin de cours — poussée tout en bas de la dernière page
\newcommand{\findecours}{%
  \par\vfill\nopagebreak
  \begin{center}\popsquiggle{poporange}\end{center}\vspace{4pt}
  \signatureonbuch
}
\AtBeginDocument{\def\llbracket{\mathopen{[\mkern-3.5mu[}}\def\rrbracket{\mathclose{]\mkern-3.5mu]}}} % intervalles d'entiers (glyphe ⟦ absent de la police)
% Glyphes absents de Fira Math : redéfinis à partir de glyphes disponibles
\AtBeginDocument{%
  \def\setminus{\mathbin{\backslash}}\let\smallsetminus\setminus
  \def\vdots{\mathord{\vbox{\baselineskip3.2pt\lineskiplimit0pt\kern4pt\hbox{.}\hbox{.}\hbox{.}}}}%
  \def\ddots{\mathinner{\mkern1mu\raise6.5pt\hbox{.}\mkern2mu\raise3.6pt\hbox{.}\mkern2mu\raise0.7pt\hbox{.}\mkern1mu}}%
  \def\triangle{\mathord{\scalebox{0.9}{$\symup{\Delta}$}}}%
  \def\nabla{\mathord{\raisebox{\depth}{\scalebox{1}[-1]{$\symup{\Delta}$}}}}%
  \def\checkmark{\popcheck{popdark}}%
  \def\star{\mathbin{\ast}}\def\bigstar{\mathord{\ast}}%
  \def\diamond{\mathbin{\scalebox{0.6}{$\Diamond$}}}%
  \def\bigcup{\mathop{\mathchoice{\raisebox{-0.3em}{\scalebox{1.7}{$\cup$}}}{\scalebox{1.25}{$\cup$}}{\cup}{\cup}}\displaylimits}%
  \def\bigcap{\mathop{\mathchoice{\raisebox{-0.3em}{\scalebox{1.7}{$\cap$}}}{\scalebox{1.25}{$\cap$}}{\cap}{\cap}}\displaylimits}%
  \def\lesseqqgtr{\mathrel{\lessgtr}}\def\bigoplus{\mathop{\oplus}\displaylimits}%
}
\providecommand{\ssi}{si et seulement si} % abréviation fréquente
\AtBeginDocument{\providecommand{\wideparen}{\overparen}} % arc de cercle

%% FIGFIX-BEGIN v4 — correctifs globaux des figures (docs/onbuch-generation/tools/figfix)
\usepackage{adjustbox}\usepackage{etoolbox}
% toute figure TikZ/pgfplots plus large que la ligne est réduite à la largeur disponible
\BeforeBeginEnvironment{tikzpicture}{\begin{adjustbox}{max width=\linewidth}}
\AfterEndEnvironment{tikzpicture}{\end{adjustbox}}
% légendes pgfplots lisibles par-dessus les courbes
\pgfplotsset{every axis/.append style={legend style={fill=white,fill opacity=0.92,text opacity=1,draw=black!15,inner sep=3pt}}}
% étiquettes d'arêtes (syntaxe edge["..."]) avec fond blanc
\tikzset{every edge quotes/.append style={fill=white,inner sep=1.2pt}}
%% FIGFIX-END
\begin{document}

\pagegarde{L’intégrité morale}{ECM}{1re Tech}{ECM – classe de Première technique}{N02}{3 séances (fiche de progression harmonisée)}{Cette leçon analyse le concept d'intégrité morale comme cohérence entre valeurs, paroles et actes. Elle identifie les institutions camerounaises garantes de cette vertu (CONAC, CONSUPE, Justice) et propose un travail dirigé sur les manifestations de l'intégrité (honnêteté, probité, responsabilité) et de ses contraires (triche, corruption, plagiat) en milieu scolaire, pour former des citoyens intègres et des professionnels fiables.
\vspace{4pt}
\begin{multicols}{2}\small
\begin{itemize}
\item À la fin de cette leçon, je sais définir l'intégrité morale et la distinguer de l'honnêteté simple et de la probité.
\item J'identifie les piliers de l'intégrité (cohérence, courage moral, responsabilité, transparence) et je les relie aux valeurs républicaines camerounaises.
\item Je cite et explique le rôle des institutions nationales de promotion et de préservation de l'intégrité (CONAC, CONSUPE, Tribunal des Comptes, Commission Nationale des Droits de l'Homme et des Libertés).
\item J'analyse des cas concrets de corruption, népotisme, tricherie aux examens ou détournement de fonds publics à l'aide de documents (rapports CONAC, articles de presse, statistiques).
\item Je construis un organigramme fonctionnel des institutions de contrôle au Cameroun.
\item J'évalue les conséquences de l'absence d'intégrité sur le développement national (indice de perception de la corruption, fuite des capitaux, dégradation des services publics).
\end{itemize}\end{multicols}}

\begin{sommaire}
\begin{multicols}{2}
\begin{enumerate}
\item 1. L'intégrité morale : concept, piliers et enjeux
\item 2. Le dispositif institutionnel camerounais de promotion et de préservation
\item 3. TD2 : L'intégrité morale en milieu scolaire - Diagnostic et Analyse
\item 4. TP : Construction d'outils de promotion de l'intégrité (Cartes, Croquis, Chartes)
\item 5. Intégrité et avenir professionnel technique : Éthique de l'ingénieur/technicien
\item 6. Synthèse évaluative et Engagement citoyen
\item Activité d'intégration
\item Méthodes et exercices résolus
\item Exercices
\item Corrigés des exercices
\item Fiche bilan
\end{enumerate}
\end{multicols}
\end{sommaire}

\begin{objectifs}
\begin{itemize}
\item À la fin de cette leçon, je sais définir l'intégrité morale et la distinguer de l'honnêteté simple et de la probité.
\item J'identifie les piliers de l'intégrité (cohérence, courage moral, responsabilité, transparence) et je les relie aux valeurs républicaines camerounaises.
\item Je cite et explique le rôle des institutions nationales de promotion et de préservation de l'intégrité (CONAC, CONSUPE, Tribunal des Comptes, Commission Nationale des Droits de l'Homme et des Libertés).
\item J'analyse des cas concrets de corruption, népotisme, tricherie aux examens ou détournement de fonds publics à l'aide de documents (rapports CONAC, articles de presse, statistiques).
\item Je construis un organigramme fonctionnel des institutions de contrôle au Cameroun.
\item J'évalue les conséquences de l'absence d'intégrité sur le développement national (indice de perception de la corruption, fuite des capitaux, dégradation des services publics).
\item Je rédige une charte de l'intégrité pour ma classe/école en justifiant chaque article par des arguments éthiques et juridiques.
\item Je propose des stratégies personnelles de résistance à la corruption passive (petits pots-de-vin, favoritisme) dans mon futur milieu professionnel technique.
\end{itemize}
\end{objectifs}

\begin{prerequis}
\begin{itemize}
\item Notion de valeur morale, de norme et de sanction (programme de 4ème / Seconde).
\item Connaissance basique de l'organisation politique du Cameroun (institutions de la République).
\item Distinction entre droit, morale et religion (Secondes).
\item Notion de bien commun et d'intérêt général (Secondes).
\item Techniques de lecture de documents statistiques (tableaux, graphiques d'évolution).
\end{itemize}
\end{prerequis}

\newpage

\coursec{L'intégrité morale : concept, piliers et enjeux}

\courssub{Situation d'accroche et problématique}

Dans un lycée de Yaoundé, un élève délégué découvre que le trésorier de l'association des élèves a détourné \num{50000} FCFA de la caisse de la kermesse pour s'acheter un smartphone. Face à l'indignation générale, le trésorier plaide : « C'est juste un prêt, je rembourserai ! » Le chef d'établissement doit trancher : sanctionner ou étouffer l'affaire pour préserver l'image de l'école ?

Cette situation, banale en apparence, condense toute la difficulté de la vie morale. Le trésorier n'a pas volé sous la menace d'une arme : il a simplement cédé à une tentation, en se racontant une histoire pour se rassurer. Le délégué, lui, hésite : dénoncer un camarade, est-ce trahir ? Le proviseur, enfin, est tenté par la solution de facilité : cacher le problème pour protéger la réputation de l'établissement. Chacun se trouve face à un choix où ses \cle{valeurs}, ses \cle{intérêts} et ses \cle{actes} risquent de se séparer.

\textbf{Problématique :} Qu'est-ce qu'agir avec intégrité morale ? Sur quels piliers repose cette vertu, quels ennemis la menacent, et pourquoi constitue-t-elle un bien précieux pour l'individu, pour l'école et pour la nation camerounaise ?

\courssub{Définition de l'intégrité morale}

\textbf{Intuition d'abord.} Imagine un mur bien construit : chaque brique est à sa place, aucune fissure ne le traverse. On dit d'un tel mur qu'il est « intact ». L'intégrité morale, c'est exactement cela appliqué à une personne : un être humain dont les convictions, les paroles et les actes forment un tout cohérent, sans fissure ni double fond.

\textbf{Étymologie.} Le mot « intégrité » vient du latin \emph{integer}, qui signifie « entier, intact, non touché ». De cette racine viennent aussi les mots « intégral » et « entier » (comme les nombres entiers, qui ne sont pas fractionnés). Être intègre, c'est donc être \cle{un} : ne pas être divisé entre ce que l'on pense, ce que l'on dit et ce que l'on fait.

\begin{definition}[Intégrité morale]
L'\cle{intégrité morale} est la vertu consistant à agir conformément à ses principes éthiques et aux lois, de manière \textbf{constante}, \textbf{transparente} et \textbf{désintéressée}, \emph{même en l'absence de contrôle externe}.
\end{definition}

Trois mots de cette définition méritent l'attention :
\begin{itemize}
\item \textbf{Constante} : l'intégrité n'est pas un comportement occasionnel. Celui qui est honnête un jour sur deux n'est pas intègre à moitié : il n'est pas intègre.
\item \textbf{Désintéressée} : on n'agit pas bien pour obtenir une récompense ou éviter une sanction, mais parce que c'est bien.
\item \textbf{Même sans contrôle} : c'est le critère décisif. Le vrai test de l'intégrité se passe quand personne ne regarde : pas de caméra, pas de surveillant, pas de gendarme.
\end{itemize}

\begin{exemplebox}[Le caissier du marché Mokolo]
Une caissière d'une coopérative du marché Mokolo à Yaoundé constate une erreur en sa faveur de \num{20000} FCFA dans la caisse, un vendredi soir. Personne ne s'en apercevra avant l'inventaire du mois suivant, et l'erreur pourrait être attribuée à un autre employé. Elle signale immédiatement l'erreur au président de la coopérative. Elle n'a agi ni par peur ni par calcul : ses convictions (l'honnêteté), son discours (elle dénonce la fraude) et ses actes (elle restitue) sont alignés. C'est l'intégrité en acte.
\end{exemplebox}

\courssub{Honnêteté, probité, intégrité : trois notions à distinguer}

Ces trois mots sont souvent confondus. Or ils désignent des réalités de plus en plus exigeantes, comme des cercles imbriqués.

\begin{definition}[Trois degrés de l'exigence morale]
\begin{itemize}
\item L'\cle{honnêteté} est la vertu élémentaire : ne pas mentir, ne pas voler, ne pas tromper autrui dans les relations ordinaires.
\item La \cle{probité} est l'honnêteté rigoureuse appliquée à la gestion des affaires publiques : le fonctionnaire probe ne détourne pas un franc de l'argent public et ne favorise personne dans l'exercice de sa fonction.
\item L'\cle{intégrité} est la vertu globale : une intransigeance éthique qui engage la personne tout entière, dans sa vie publique comme dans sa vie privée, même sans témoin.
\end{itemize}
\end{definition}

\begin{popfigure}
\begin{center}
\begin{tikzpicture}[scale=1.0]
% Cercle extérieur : Intégrité
\draw[fill=popgreen!18,draw=popgreen!70!black,line width=1.2pt] (0,0) circle (4.4);
% Cercle moyen : Probité
\draw[fill=popblue!20,draw=popblue!70!black,line width=1.2pt] (-0.7,-0.3) circle (3.0);
% Cercle intérieur : Honnêteté
\draw[fill=poporange!25,draw=poporange!80!black,line width=1.2pt] (-1.2,-0.6) circle (1.6);
% Étiquettes
\node[font=\bfseries\small,text=popgreen!50!black] at (2.9,3.4) {INTÉGRITÉ};
\node[font=\bfseries\small,text=popblue!60!black] at (1.7,1.9) {PROBITÉ};
\node[font=\bfseries\small,text=poporange!80!black] at (-1.2,-0.6) {HONNÊTETÉ};
% Exemples zone honnêteté
\node[font=\scriptsize,text=popdark,align=center] at (-1.2,-1.9) {Rendre la monnaie exacte\\au client de l'atelier};
% Exemples zone probité seule
\node[font=\scriptsize,text=popdark,align=center] at (1.6,0.6) {Gestionnaire public rigoureux\\mais vie privée trouble};
% Exemples zone intégrité seule
\node[font=\scriptsize,text=popdark,align=center] at (0.6,-3.4) {Refuse le pot-de-vin, déclare ses biens,\\cohérent en public comme en privé};
\end{tikzpicture}
\end{center}
\legende{Ensembles imbriqués : Honnêteté $\subset$ Probité $\subset$ Intégrité. Chaque degré contient le précédent et ajoute une exigence supplémentaire.}
\end{popfigure}

\begin{attention}[Ne pas confondre les degrés]
Un gestionnaire peut être \emph{probe} (jamais un franc détourné) sans être \emph{intègre} : s'il trompe son entourage, triche aux examens de concours ou maltraite ses employés à domicile, sa vertu reste partielle. À l'inverse, l'intégrité inclut nécessairement la probité et l'honnêteté. Dans une copie d'examen, ne jamais employer ces trois termes comme des synonymes : le correcteur attend la hiérarchie précise.
\end{attention}

\begin{propriete}[Indivisibilité de l'intégrité]
L'intégrité est \cle{indivisible} : on ne peut pas être intègre dans sa vie professionnelle et corrompu dans sa vie privée. La « compartmentalisation » éthique (être vertueux au bureau, malhonnête à la maison) est une illusion : tôt ou tard, les habitudes de la sphère corrompue contaminent l'autre, car le caractère se forme par la répétition des actes.
\end{propriete}

\courssub{Les quatre piliers de l'intégrité}

L'intégrité repose sur quatre piliers solidaires : si l'un s'effondre, tout l'édifice vacille.

\begin{enumerate}
\item \textbf{La \cle{cohérence}} : c'est l'alignement entre les valeurs proclamées et les actes réels. Celui qui condamne la tricherie en classe mais recopie le devoir de son voisin manque de cohérence. C'est le pilier central, directement issu de l'étymologie \emph{integer}.
\item \textbf{Le \cle{courage moral}} : c'est la capacité de dire \emph{non} à la facilité, même sous la pression du groupe ou de l'autorité. Refuser de participer à une fraude collective à un examen, alors que « tout le monde le fait », exige plus de courage que bien des exploits physiques.
\item \textbf{La \cle{responsabilité}} : c'est assumer les conséquences de ses actes, sans chercher de bouc émissaire. Le trésorier de notre situation d'accroche qui prétend « c'est juste un prêt » refuse précisément d'assumer.
\item \textbf{La \cle{transparence}} : c'est accepter de rendre des comptes, d'exposer ses décisions et ses gestes au regard des autres. La transparence est la preuve sociale de l'intégrité : qui n'a rien à cacher n'a rien à craindre de la lumière.
\end{enumerate}

\begin{popfigure}
\begin{center}
\begin{tikzpicture}[scale=0.92, every node/.style={transform shape}]
% Sol
\draw[fill=popcream,draw=none] (-6,-3.2) rectangle (6,-2.2);
\draw[popline] (-6,-2.2) -- (6,-2.2);
% Racines
\draw[popgold!80!black,line width=2pt] (0,-2.2) -- (-1.8,-3.0);
\draw[popgold!80!black,line width=2pt] (0,-2.2) -- (0,-3.1);
\draw[popgold!80!black,line width=2pt] (0,-2.2) -- (1.8,-3.0);
\node[font=\scriptsize\itshape,text=popdark] at (0,-3.5) {Racines : valeurs familiales et culturelles};
% Tronc
\draw[fill=popgold!60!black,draw=popgold!50!black] (-0.55,-2.2) rectangle (0.55,0.6);
\node[font=\scriptsize\bfseries,text=white,align=center] at (0,-0.8) {COURAGE\\MORAL};
% Feuillage
\draw[fill=popgreen!35,draw=popgreen!60!black,line width=1pt] (0,2.6) circle (2.9);
% Branches = piliers
\node[font=\scriptsize\bfseries,text=popdark,align=center,fill=white,rounded corners=2pt,inner sep=2pt] at (-2.3,3.4) {Cohérence};
\node[font=\scriptsize\bfseries,text=popdark,align=center,fill=white,rounded corners=2pt,inner sep=2pt] at (2.3,3.4) {Responsabilité};
\node[font=\scriptsize\bfseries,text=popdark,align=center,fill=white,rounded corners=2pt,inner sep=2pt] at (-2.3,1.6) {Transparence};
\node[font=\scriptsize\bfseries,text=popdark,align=center,fill=white,rounded corners=2pt,inner sep=2pt] at (2.3,1.6) {Courage};
% Fruits
\foreach \p/\t in {(-1.1,4.6)/Confiance,(0.4,5.0)/Paix,(1.5,4.3)/Développement}{
  \draw[fill=poporange!80,draw=poporange!60!black] \p circle (0.28);
  \node[font=\scriptsize,text=popdark] at \p [yshift=-16pt] {\t};
}
% Parasites
\foreach \p/\t in {(-4.6,0.4)/Corruption,(-4.9,2.2)/Népotisme,(4.7,0.6)/Tricherie,(4.9,2.4)/Trafic d'influence}{
  \node[font=\scriptsize\bfseries,text=poppink!70!black,fill=poppink!15,rounded corners=2pt,inner sep=2pt] at \p {\t};
  \draw[-{Stealth},poppink!70!black,line width=1pt] \p -- ++(0,0.6);
}
\end{tikzpicture}
\end{center}
\legende{L'arbre de l'intégrité : les racines (valeurs familiales et culturelles) nourrissent le tronc (courage moral), qui porte les branches (les quatre piliers) et les fruits (confiance, paix, développement). Les parasites (corruption, népotisme, tricherie, trafic d'influence) l'attaquent de l'extérieur.}
\end{popfigure}

\begin{methode}[Le test de l'intégrité en trois questions]
Face à un choix délicat, pose-toi successivement trois questions :
\begin{enumerate}
\item \textbf{Est-ce légal ?} L'acte respecte-t-il la loi et le règlement ?
\item \textbf{Est-ce équitable pour tous ?} Accepterais-je que n'importe qui agisse ainsi à ma place ?
\item \textbf{Le test de la Une} : accepterais-je que cet acte soit publié en première page de \emph{Cameroon Tribune}, sous mon nom, et lu par ma famille, mes enseignants, mon futur employeur ?
\end{enumerate}
Si la réponse est « non » à l'une seule de ces questions, il y a atteinte à l'intégrité : renonce à l'acte. Appliqué à notre trésorier : le détournement est illégal (question 1), inéquitable pour les autres élèves (question 2), et impubliable (question 3). Triple échec.
\end{methode}

\courssub{Les ennemis de l'intégrité}

L'intégrité a des adversaires identifiés, que le droit camerounais sanctionne et que la morale condamne.

\begin{itemize}
\item \textbf{La \cle{corruption}} : elle est \emph{active} quand on propose ou donne un avantage indu (le pot-de-vin offert), \emph{passive} quand on le sollicite ou l'accepte. Les deux sont condamnables : il n'y a pas de corrompu sans corrupteur.
\item \textbf{Le \cle{népotisme} et le \cle{clientélisme}} : favoriser un parent, un ami ou un « frère de village » dans une nomination ou un marché, au détriment du mérite. Le clientélisme échange des faveurs contre des soutiens politiques ou personnels.
\item \textbf{Le \cle{trafic d'influence}} : monnayer sa position ou ses relations pour obtenir un avantage (ex. : être payé pour « intervenir » auprès d'un jury de concours).
\item \textbf{L'\cle{enrichissement illicite}} : accroissement du patrimoine d'un agent public sans justification par ses revenus légaux.
\item \textbf{La \cle{tricherie académique}} : fraude aux examens, plagiat, achat de diplômes. C'est la forme scolaire de la corruption : celui qui triche à 16 ans s'entraîne à corrompre à 30 ans.
\end{itemize}

\begin{exemplebox}[Le coût chiffré de la corruption au Cameroun]
Selon Transparency International (indice de perception de la corruption, CPI 2023), le Cameroun obtient le score de \textbf{26/100} et se classe \textbf{140\textsuperscript{e} sur 180 pays}. Le coût de la corruption est estimé entre \textbf{5 \% et 10 \% du PIB par an} (rapports de la CONAC), soit des centaines de milliards de FCFA qui manquent chaque année pour construire des hôpitaux, des routes et des salles de classe. La corruption n'est donc pas un « petit arrangement » : c'est un prélèvement sur l'avenir de chaque Camerounais.
\end{exemplebox}

\courssub{Le cadre juridique camerounais}

L'intégrité n'est pas seulement une affaire de conscience : le Cameroun l'a inscrite dans son droit.

\begin{itemize}
\item \textbf{La Constitution du 18 janvier 1996} (Préambule) : le peuple camerounais y proclame son attachement aux valeurs de justice, d'honnêteté et de travail, et affirme que « tout individu a droit à la vie, à l'intégrité physique et morale ».
\item \textbf{La loi n\textsuperscript{o} 2006/015 du 29 décembre 2006} portant statut général des fonctionnaires de l'État : elle impose aux agents publics des devoirs de probité, de neutralité et de discrétion, et interdit le cumul d'intérêts privés avec la fonction.
\item \textbf{La loi n\textsuperscript{o} 2011/002 du 14 juillet 2011} relative à la lutte contre la corruption : elle définit et réprime pénalement la corruption active et passive, le trafic d'influence, l'enrichissement illicite et les infractions assimilées.
\end{itemize}

\begin{aretenir}[L'essentiel du cadre juridique]
La Constitution de 1996 consacre l'intégrité comme valeur républicaine ; le statut de la fonction publique (loi de 2006) en fait un devoir professionnel des agents de l'État ; la loi de 2011 en fait une obligation pénale dont la violation est punie par les tribunaux. L'intégrité est ainsi à la fois une \textbf{vertu morale}, une \textbf{règle professionnelle} et une \textbf{exigence légale}.
\end{aretenir}

\courssub{L'intégrité, un capital social et professionnel}

Pourquoi l'intégrité est-elle un \emph{investissement} et non une perte ? Parce qu'elle produit trois biens précieux.

\begin{itemize}
\item \textbf{La confiance} : la confiance se gagne lentement par des actes intègres répétés, et se perd en une seule faute. Un technicien dont la parole est sûre devient indispensable : on lui confie les clés, les comptes, les machines coûteuses.
\item \textbf{L'employabilité} : dans les métiers techniques (électricité, mécanique, comptabilité, bâtiment), l'employeur redoute deux choses : l'incompétence et la malhonnêteté. Or l'incompétence se corrige par la formation ; la malhonnêteté, presque jamais. Un jeune diplômé réputé intègre est embauché avant un brillant tricheur.
\item \textbf{La paix sociale} : là où chacun respecte les règles sans qu'on le surveille, on économise les contrôles, les procès et les conflits. L'intégrité collective est le ciment invisible du vivre-ensemble.
\end{itemize}

\begin{attention}[Deux confusions fréquentes]
\begin{itemize}
\item \textbf{Intégrité $\neq$ intégrisme} : l'intégrisme est une rigidité dogmatique qui refuse toute nuance et s'impose aux autres. L'intègre, lui, tient ses principes avec humilité et discernement.
\item \textbf{Intégrité $\neq$ naïveté} : être intègre n'est pas ignorer les codes sociaux ni se laisser exploiter. L'intégrité demande de l'intelligence : savoir refuser fermement, mais avec tact et au bon moment.
\end{itemize}
\end{attention}

\begin{exoresolu}[Appliquer le test de l'intégrité]
\textbf{Énoncé.} Un élève de Première technique, en stage dans une entreprise de Bafoussam, reçoit de son maître de stage l'ordre de facturer à un client des pièces neuves alors que des pièces d'occasion ont été montées. Le maître de stage ajoute : « Tout le monde fait ça, et ton attestation de stage en dépend. » En utilisant le test de l'intégrité en trois questions, analyse la situation et propose une conduite à tenir.
\tcblower
\textbf{Solution détaillée.}
\begin{enumerate}
\item \emph{Est-ce légal ?} Non. Facturer des pièces neuves non fournies constitue une escroquerie, punie par le Code pénal, et une violation de la loi n\textsuperscript{o} 2011/002 du 14 juillet 2011 si l'entreprise travaille sur des marchés publics.
\item \emph{Est-ce équitable pour tous ?} Non. Le client paie pour une qualité qu'il ne reçoit pas ; les concurrents honnêtes sont lésés ; l'argument « tout le monde fait ça » est faux et, même s'il était vrai, une injustice collective ne devient pas juste.
\item \emph{Le test de la Une} : l'élève accepterait-il de lire dans \emph{Cameroon Tribune} : « Un stagiaire facture des pièces d'occasion au prix du neuf » ? Évidemment non : sa réputation et son avenir professionnel seraient détruits.
\end{enumerate}
\textbf{Conduite à tenir} : refuser poliment mais fermement d'établir la fausse facture, en invoquant le règlement et la loi ; en informer si nécessaire la hiérarchie ou le chef d'établissement scolaire qui suit le stage ; documenter les faits. Le courage moral (pilier 2) et la responsabilité (pilier 3) priment sur la menace concernant l'attestation : un diplôme obtenu par la compromission est un fruit empoisonné.
\end{exoresolu}

\begin{exercice}[À toi de juger]
Reprends la situation d'accroche (le trésorier et les \num{50000} FCFA). 1) Identifie, pour chacun des trois personnages (trésorier, délégué, proviseur), le pilier de l'intégrité mis en jeu. 2) Applique le test des trois questions à la décision du proviseur d'étouffer l'affaire. 3) Rédige en cinq lignes la décision qu'un proviseur intègre devrait prendre, en la justifiant.
\end{exercice}

\begin{corrige}[Exercice : À toi de juger]
1) Le trésorier manque de \textbf{cohérence} (il a accepté une charge de confiance qu'il trahit) et de \textbf{responsabilité} (il nie la faute en parlant de « prêt »). Le délégué fait face à un problème de \textbf{courage moral} (dénoncer malgré la pression amicale). Le proviseur est tenté de sacrifier la \textbf{transparence} à l'image de l'école.
2) Étouffer l'affaire est légal à la limite (le proviseur peut choisir une sanction interne), mais inéquitable (les élèves lésés ne sont pas dédommagés, le fautif est encouragé) et impubliable (un journal titrerait : « Un lycée cache un détournement »). Deux réponses négatives : atteinte à l'intégrité.
3) Un proviseur intègre constate les faits, entend le trésorier, exige le remboursement intégral, applique la sanction prévue par le règlement intérieur et transforme l'incident en moment éducatif pour toute la communauté scolaire. Ainsi, la sanction répare, dissuade et forme, au lieu de cacher.
\end{corrige}

\begin{savaistu}[L'intégrité dans les cultures camerounaises]
Bien avant les textes modernes, les sociétés camerounaises valorisaient la parole donnée et la droiture. Chez de nombreux peuples, la palabre et les serments traditionnels engageaient la personne tout entière devant la communauté et les ancêtres : trahir sa parole, c'était se « briser » soi-même. La notion latine d'\emph{integer} rejoint ainsi des sagesses locales très anciennes : l'intégrité est un héritage à faire fructifier, pas une importation.
\end{savaistu}

\coursec{Le dispositif institutionnel camerounais de promotion et de préservation de l'intégrité morale}

Après avoir clarifié le concept d'intégrité morale, ses piliers (honnêteté, probité, transparence, responsabilité, exemplarité) et ses enjeux pour la cohésion sociale et le développement, il convient d'examiner l'arsenal institutionnel mis en place par l'État du Cameroun pour transformer ces idéaux en réalités tangibles. L'intégrité ne se décrète pas ; elle se construit, se surveille et se sanctionne à travers un dispositif où s'articulent des organes de prévention, de contrôle, de répression et de veille citoyenne.

\courssub{Les cinq piliers institutionnels : une architecture à niveaux multiples}

La lutte pour l'intégrité morale au Cameroun repose sur cinq institutions majeures, chacune dotée d'un mandat spécifique, mais appelées à agir en synergie. On peut les classer selon leur fonction dominante : prévention et répression de la corruption, contrôle supérieur de l'action de l'État, contrôle des finances publiques, promotion des droits humains et de l'éthique, et vérification patrimoniale.

\begin{popfigure}
\begin{tikzpicture}[
    node distance=1.2cm and 1.5cm,
    pilier/.style={rectangle, rounded corners, draw=popblue, thick, fill=popblueL, text width=3.2cm, align=center, minimum height=1.8cm, font=\small\bfseries},
    mission/.style={rectangle, draw=popgreen, fill=popgreenL, text width=3.0cm, align=center, minimum height=1.2cm, font=\small},
    arrow/.style={-Stealth, thick, popdark},
    level1/.style={rectangle, rounded corners, draw=poppurple, thick, fill=poppurpleL, text width=4cm, align=center, minimum height=1.5cm, font=\small\bfseries},
    level3/.style={rectangle, rounded corners, draw=poppink, thick, fill=poppinkL, text width=3.5cm, align=center, minimum height=1.3cm, font=\small\bfseries},
    level4/.style={rectangle, rounded corners, draw=popgold, thick, fill=popgoldL, text width=3.5cm, align=center, minimum height=1.3cm, font=\small\bfseries},
]
% Niveau 1 : Impulsion
\node[level1, align=center] (presidence) at (0, 7){Présidence de la République\\Impulsion politique \& Constitution};

% Niveau 2 : Organes de contrôle / Promotion
\node[pilier, align=center] (conac) at (-7, 4){CONAC\\Commission Nationale Anti-Corruption};
\node[pilier, align=center] (consupe) at (-3.5, 4){CONSUPE\\Contrôle Supérieur de l'État};
\node[pilier, align=center] (cdc) at (0, 4){Chambre des Comptes\\(Cour Suprême)};
\node[pilier, align=center] (cndhl) at (3.5, 4){CNDHL\\Droits de l'Homme \& Libertés};
\node[pilier] (cdb) at (7, 4) {Commission Déclaration des Biens et Avoirs};

% Missions
\node[mission, align=center] (m1) at (-7, 1.5){Éducation, Prévention\\Enquête, Transmission};
\node[mission, align=center] (m2) at (-3.5, 1.5){Audit des services publics\\Contrôle de gestion};
\node[mission, align=center] (m3) at (0, 1.5){Contrôle des comptes\\Jugement des comptables};
\node[mission, align=center] (m4) at (3.5, 1.5){Promotion DDH\\Éthique publique};
\node[mission, align=center] (m5) at (7, 1.5){Déclaration entrée/sortie\\de fonction, vérification};

% Niveau 3 : Répression
\node[level3, align=center] (justice) at (0, -1.5){Appareil Judiciaire\\Parquets, TGI, Cour Suprême};

% Niveau 4 : Société Civile
\node[level4, align=center] (sc) at (0, -4){Société Civile \& Médias\\Transparency Int. Cameroun, lanceurs d'alerte};

% Flèches verticales
\draw[arrow] (presidence.south) -- ++(0,-0.5) -| (conac.north);
\draw[arrow] (presidence.south) -- ++(0,-0.5) -| (consupe.north);
\draw[arrow] (presidence.south) -- ++(0,-0.5) -| (cdc.north);
\draw[arrow] (presidence.south) -- ++(0,-0.5) -| (cndhl.north);
\draw[arrow] (presidence.south) -- ++(0,-0.5) -| (cdb.north);

\draw[arrow] (conac.south) -- (m1.north);
\draw[arrow] (consupe.south) -- (m2.north);
\draw[arrow] (cdc.south) -- (m3.north);
\draw[arrow] (cndhl.south) -- (m4.north);
\draw[arrow] (cdb.south) -- (m5.north);

% Flux vers la justice
\draw[arrow, poporange] (m1.south) -- ++(0,-0.3) -| node[near start, left, font=\tiny, text width=2cm, align=center] {Signalement / Dénonciation} (justice.west);
\draw[arrow, poporange] (m2.south) -- ++(0,-0.3) -| node[near start, left, font=\tiny, text width=2cm, align=center] {Irrégularités de gestion} (justice.west);
\draw[arrow, poporange] (m3.south) -- ++(0,-0.3) -| node[near start, right, font=\tiny, text width=2cm, align=center] {Fautes de gestion} (justice.east);
\draw[arrow, poporange] (m5.south) -- ++(0,-0.3) -| node[near start, right, font=\tiny, text width=2cm, align=center] {Enrichissement illicite} (justice.east);

\draw[arrow] (justice.south) -- node[left, font=\small] {Sanctions pénales / administratives} (sc.north);

% Boucle rétroaction
\draw[arrow, dashed, popblue] (sc.east) .. controls +(2,0) and +(2,0) .. node[right, font=\small, align=center] {Plaidoyer\\Rapports\\Veille} (presidence.east);
\draw[arrow, dashed, popblue] (justice.east) .. controls +(1.5,0) and +(1.5,0) .. node[right, font=\small, align=center] {Jurisprudence\\Dissuasion} (conac.east);

\end{tikzpicture}
\legende{Architecture institutionnelle de l'intégrité au Cameroun : de l'impulsion politique au contrôle citoyen, en passant par les organes spécialisés et la répression judiciaire.}
\end{popfigure}

\textbf{1. La CONAC (Commission Nationale Anti-Corruption) : le chef de file.}\par
Créée par le Décret n°2006/088 du 31 mars 2006 et renforcée par la Loi n°2011/020 du 14 décembre 2011, la CONAC a pour mission de contribuer à la lutte contre la corruption sur l'ensemble du territoire national. Son champ d'action couvre la prévention (éducation, sensibilisation), la réception des dénonciations, l'enquête administrative et la transmission des dossiers à la justice.

\textbf{2. Le CONSUPE (Contrôle Supérieur de l'État) : le contrôleur général des services publics.}\par
Placé sous l'autorité du Président de la République, le CONSUPE effectue des missions de contrôle, d'audit et d'évaluation de la gestion des services publics, des établissements et entreprises publics et parapublics. Il traque les irrégularités de gestion (détournements, surfacturations, mauvaise exécution des marchés publics) et transmet ses rapports aux autorités compétentes, y compris à la justice.

\textbf{3. La Chambre des Comptes de la Cour Suprême : la gardienne de la probité financière.}\par
Juridiction financière prévue par la Constitution et organisée par la Loi n°2003/005 du 21 avril 2003, elle juge les comptes des comptables publics, contrôle la régularité des dépenses et des recettes publiques, et sanctionne les fautes de gestion (prononcé de débets, amendes). Elle est le pendant juridictionnel du contrôle administratif des finances.

\textbf{4. La CNDHL (Commission Nationale des Droits de l'Homme et des Libertés) : le socle éthique.}\par
Institution indépendante créée par la Loi n°2004/016 du 22 juillet 2004, elle promeut et protège les droits humains. L'intégrité morale est indissociable du respect de la dignité humaine, de la non-discrimination et de l'accès à la justice. Elle enquête sur les violations (mauvais traitements, arrestations arbitraires, corruption dans la police ou la justice) et émet des avis et recommandations.

\textbf{5. La Commission de la Déclaration des Biens et Avoirs : le miroir patrimonial.}\par
Prévue par l'article 66 de la Constitution, elle reçoit les déclarations de patrimoine des hautes personnalités de l'État (Président de la République, Premier ministre, ministres, députés, sénateurs, magistrats, directeurs généraux, etc.) à l'entrée et à la sortie de fonction. Elle vérifie la sincérité de ces déclarations et peut saisir le Procureur Général près la Cour Suprême en cas d'enrichissement illicite présumé.

\begin{definition}[CONAC -- Commission Nationale Anti-Corruption]
Organe national chargé de contribuer à la lutte contre la corruption au Cameroun. Créée par le Décret n°2006/088 du 31 mars 2006, renforcée par la Loi n°2011/020 du 14 décembre 2011, elle dispose de la personnalité juridique et de l'autonomie financière. Ses missions : éducation civique, prévention, réception des dénonciations, enquête administrative, transmission des dossiers à la justice.
\end{definition}

\courssub{Focus approfondi : la CONAC, organisation et pouvoirs}

La CONAC n'est pas une simple structure de sensibilisation ; c'est une institution hiérarchisée dotée de pouvoirs d'investigation réels, bien que limités à la phase administrative : elle enquête mais ne juge pas.

\begin{popfigure}
\begin{tikzpicture}[
    node distance=1cm and 1.5cm,
    org/.style={rectangle, rounded corners, draw=popblue, thick, fill=popblueL, text width=4.5cm, align=center, minimum height=1.2cm, font=\small\bfseries},
    dept/.style={rectangle, draw=popgreen, fill=popgreenL, text width=3cm, align=center, minimum height=1cm, font=\small},
    arrow/.style={-Stealth, thick, popdark},
]
% Top
\node[org, align=center] (commission){Commission\\(Président + Commissaires)\\Décision / Orientation};
% Middle
\node[org, below=1.5cm of commission, align=center] (sg){Secrétariat Permanent\\Exécution / Administration};
% Departments
\node[dept] (dept1) at (-6, 1.5) {Prévention \& Éducation};
\node[dept] (dept2) at (-2, 1.5) {Enquêtes \& Investigations};
\node[dept] (dept3) at (2, 1.5) {Affaires Juridiques};
\node[dept] (dept4) at (6, 1.5) {Coopération \& Communication};
% Regional
\node[org, align=center] (reg) at (0, -1.5){Antennes régionales\\Proximité / Réception des plaintes};

\draw[arrow] (commission) -- (sg);
\draw[arrow] (sg.west) -| (dept1.south);
\draw[arrow] (sg.west) -| (dept2.south);
\draw[arrow] (sg.east) -| (dept3.south);
\draw[arrow] (sg.east) -| (dept4.south);
\draw[arrow] (sg) -- (reg);

% Pouvoirs
\node[draw=poporange, thick, fill=poporangeL, rounded corners, text width=8cm, align=left, font=\small, anchor=north] at (0, -3.5) {
\textbf{Pouvoirs clés :} auto-saisine ou saisine sur dénonciation ; audition de tout agent public ou privé ; réquisition de tout document utile ; demande de gel des avoirs suspects auprès du procureur ; transmission obligatoire du dossier au parquet si des indices d'infraction pénale sont réunis.
};
\end{tikzpicture}
\legende{Organigramme simplifié de la CONAC et ses pouvoirs d'investigation.}
\end{popfigure}

\begin{methode}[Comment saisir la CONAC ? -- Démarche citoyenne]
Tout citoyen, témoin ou victime d'actes de corruption (demande de pot-de-vin, détournement, népotisme, favoritisme) peut saisir la CONAC. La procédure est gratuite et l'anonymat du dénonciateur est protégé.
\begin{enumerate}
\item \textbf{Physiquement :} au siège à Yaoundé ou dans les antennes régionales (chefs-lieux de région).
\item \textbf{Par téléphone :} numéro vert \textbf{8200} (appel gratuit).
\item \textbf{En ligne :} site web officiel de la CONAC, rubrique « Dénonciation », formulaire sécurisé.
\item \textbf{Par écrit :} courrier adressé au Président de la CONAC, avec accusé de réception.
\end{enumerate}
\textbf{Pièces à fournir :} tout élément de preuve (procès-verbal, facture, enregistrement, capture d'écran, témoignage écrit). La CONAC accuse réception et informe le dénonciateur de la suite donnée.
\end{methode}

\begin{exoresolu}[Analyse d'un signalement à la CONAC]
\textbf{Énoncé :} Un élève de Première technique observe que le surveillant général de son établissement exige 2 000 FCFA par élève pour « faciliter » l'obtention de la fiche d'inscription au Probatoire, pourtant encadrée par les textes ministériels. L'élève souhaite dénoncer sans subir de représailles. Qualifier juridiquement les faits et proposer une démarche.
\tcblower
\textbf{Solution détaillée :}
\textbf{Analyse juridique :} il s'agit d'un acte de corruption passive (le fonctionnaire sollicite ou reçoit des avantages pour accomplir un acte de sa fonction) et d'une concussion (exaction d'une somme non due). L'établissement est public, le surveillant est un agent public : la CONAC est donc compétente.
\textbf{Démarche conseillée :}
\begin{enumerate}
\item Collecter discrètement la preuve : photo de la liste des payeurs, capture d'écran de messages, témoignages écrits de camarades.
\item Saisir la CONAC par le numéro vert 8200 ou par courrier, en demandant la protection de l'anonymat.
\item Signaler parallèlement les faits au chef d'établissement (voie hiérarchique) et, en cas d'inaction, à la Délégation régionale des Enseignements secondaires.
\end{enumerate}
\textbf{Résultat attendu :} enquête administrative de la CONAC $\rightarrow$ transmission au procureur de la République si les indices sont suffisants $\rightarrow$ poursuites pénales et sanctions administratives (suspension, révocation). L'élève dénonciateur est protégé contre les représailles.
\end{exoresolu}

\courssub{Focus : le CONSUPE, contrôleur de la gestion publique, y compris dans l'éducation}

Le Contrôle Supérieur de l'État (CONSUPE) est l'organe de contrôle \emph{a posteriori} de la gestion des deniers publics. Ses missions d'audit couvrent tous les services publics, y compris les établissements scolaires publics et les ministères de l'Éducation.

\begin{propriete}[Rôle du CONSUPE en matière d'intégrité]
\begin{itemize}
\item \textbf{Audit de gestion :} contrôle de la gestion des budgets des établissements publics, des subventions et des fonds de fonctionnement ; détection des « frais fantômes » et des surfacturations.
\item \textbf{Contrôle des marchés publics :} vérification de la régularité des passations de marchés (fournitures, travaux) et de leur exécution.
\item \textbf{Évaluation des politiques publiques :} il vérifie que les moyens alloués (par exemple aux examens officiels comme le Probatoire et le Baccalauréat) sont utilisés conformément à leur destination.
\item \textbf{Transmission :} ses rapports sont adressés au Président de la République et peuvent donner lieu à des sanctions administratives ou à des saisines de la justice.
\end{itemize}
\end{propriete}

\begin{exemplebox}[Cas concret : la lutte contre la fraude aux examens officiels]
Lorsque des fuites de sujets ou des fraudes massives sont découvertes dans un centre d'examen, la chaîne d'intégrité se met en marche :
\begin{enumerate}
\item Le ministère des Enseignements secondaires (MINSEC) et l'Office du Baccalauréat du Cameroun (OBC) constatent les faits et prennent des mesures conservatoires (annulation d'épreuves, changement de centre).
\item Les sanctions administratives frappent les personnels fautifs (chefs de centre, surveillants).
\item En cas de détournement des fonds d'examen, le CONSUPE ou la Chambre des Comptes peut être saisi pour contrôle.
\item Les infractions pénales (escroquerie, faux, corruption) sont transmises au parquet.
\end{enumerate}
Cet exemple montre l'articulation entre contrôle administratif, décision ministérielle et répression judiciaire.
\end{exemplebox}

\courssub{L'articulation avec la justice : du contrôle à la sanction}

Les organes de contrôle (CONAC, CONSUPE, Chambre des Comptes, Commission de la Déclaration des Biens) sont des \textbf{organes administratifs ou juridictions financières}. À l'exception de la Chambre des Comptes qui juge les comptes, ils n'ont pas le pouvoir de condamner pénalement (prison, amendes pénales). Leur rôle est d'instruire, de constater, de proposer ou de prononcer des sanctions administratives (révocation, débets) et surtout de \textbf{transmettre} au parquet.

\begin{center}
\begin{tabularx}{\linewidth}{|l|X|X|}\hline
\rowcolor{popblueL}
\textbf{Organe de contrôle} & \textbf{Infraction-type détectée} & \textbf{Juridiction saisie ensuite} \\ \hline
CONAC & Corruption, concussion, trafic d'influence, enrichissement illicite & Procureur de la République / TGI du lieu de l'infraction \\ \hline
CONSUPE & Détournement de deniers publics, irrégularités de marchés & TGI / juridictions répressives compétentes \\ \hline
Chambre des Comptes & Fautes de gestion, comptes irréguliers & Elle juge les comptes ; renvoie les faits pénaux au parquet \\ \hline
Commission Déclaration des Biens & Enrichissement illicite, fausse déclaration & Procureur Général près la Cour Suprême \\ \hline
CNDHL & Violations des droits humains, corruption judiciaire & TGI / Cour d'Appel / Cour Suprême \\ \hline
\end{tabularx}
\end{center}

\begin{attention}[Erreur fréquente : confondre contrôle et jugement]
Ne pas écrire : « La CONAC a condamné X à 5 ans de prison. »
\textbf{Correct :} « La CONAC a enquêté sur X et transmis le dossier au procureur de la République. Le tribunal compétent a ensuite condamné X à 5 ans de prison ferme. »
La séparation des pouvoirs impose que l'autorité qui contrôle (administratif) soit distincte de celle qui juge (judiciaire). Seule la Chambre des Comptes a une double nature : juridictionnelle pour le jugement des comptes, administrative pour le contrôle.
\end{attention}

\courssub{Le contrôle citoyen : société civile et médias, le « cinquième pouvoir »}

L'État ne peut pas tout voir. La société civile organisée et les médias indépendants constituent un contre-pouvoir indispensable, protégé par la liberté de communication sociale (Loi n°90/052 du 19 décembre 1990).

\begin{itemize}
\item \textbf{Transparency International Cameroun (TI-C) :} section nationale de Transparency International. Elle publie des analyses sur la perception de la corruption et mène des enquêtes de terrain (marchés publics, santé, éducation). Elle plaide pour une loi de protection des lanceurs d'alerte.
\item \textbf{Les ONG d'appui aux victimes :} elles accompagnent les citoyens dans la saisine de la CONAC et forment des « veilleurs citoyens » dans les arrondissements.
\item \textbf{Les lanceurs d'alerte :} agents publics ou privés qui révèlent des faits de corruption ou de détournement dont ils ont connaissance dans l'exercice de leurs fonctions.
\end{itemize}

\begin{aretenir}[Protection des lanceurs d'alerte : un chantier inachevé]
Le Cameroun a ratifié la Convention de l'Union Africaine (2003) et la Convention des Nations Unies contre la corruption (2003), qui recommandent la protection des témoins et des dénonciateurs. La saisine de la CONAC garantit l'anonymat du dénonciateur. \textbf{Mais} il n'existe pas encore de loi-cadre générale protégeant le lanceur d'alerte dans tous les secteurs (entreprise privée, médias, administration) contre le licenciement, le harcèlement ou les poursuites abusives. C'est un chantier majeur pour rendre l'intégrité effective.
\end{aretenir}

\begin{savaistu}[Les engagements internationaux du Cameroun]
Le Cameroun a ratifié :
\begin{enumerate}
\item \textbf{La Convention de l'Union Africaine sur la prévention et la lutte contre la corruption} (Maputo, 2003). Obligations : criminaliser la corruption active et passive, l'enrichissement illicite, le blanchiment ; mettre en place des autorités nationales de lutte (comme la CONAC).
\item \textbf{La Convention des Nations Unies contre la corruption (CNUCC)} (Mérida, 2003), instrument universel doté d'un mécanisme d'examen par les pairs. Les recommandations récurrentes adressées au Cameroun portent sur l'indépendance effective des organes de lutte, la protection des lanceurs d'alerte et la transparence de la déclaration des biens.
\end{enumerate}
Ces engagements contraignent l'État à renforcer régulièrement son arsenal juridique (lois anti-corruption, Code pénal, lois contre le blanchiment).
\end{savaistu}

\courssub{Limites structurelles et défis persistants : un diagnostic lucide}

Malgré un arsenal juridique et institutionnel complet sur le papier, l'efficacité réelle bute sur des obstacles que tout citoyen lucide -- et tout futur technicien -- doit connaître pour ne tomber ni dans le cynisme ni dans la complaisance.

\begin{center}
\begin{tabularx}{\linewidth}{|l|X|X|}\hline
\rowcolor{popblueL}
\textbf{Défi} & \textbf{Manifestation concrète} & \textbf{Impact sur l'intégrité} \\ \hline
\textbf{Indépendance relative} & Les dirigeants de plusieurs organes sont nommés par décret présidentiel ; les budgets dépendent de l'exécutif. & Risque d'auto-saisine limitée sur les dossiers sensibles impliquant de hauts responsables. \\ \hline
\textbf{Moyens matériels et humains insuffisants} & Effectifs réduits, manque de véhicules et d'outils d'analyse financière modernes. & Traitement lent des dossiers ; difficulté à traiter la criminalité financière complexe. \\ \hline
\textbf{Lenteur de la procédure judiciaire} & Enquête, instruction, jugement, appel, cassation : la chaîne peut durer plusieurs années. & Impunité de fait : les preuves disparaissent, les témoins se rétractent ; la sanction tardive perd son effet dissuasif. \\ \hline
\textbf{Impunité perçue des « gros poissons »} & Les condamnations concernent surtout des agents subalternes ; peu de hauts responsables sont condamnés définitivement. & Perte de confiance des citoyens : « la loi est pour les faibles » ; démobilisation des dénonciateurs. \\ \hline
\textbf{Chevauchement des compétences} & CONAC, police judiciaire, gendarmerie, inspections générales d'État et ministérielles se concurrencent parfois. & Non-partage des informations, dossiers « enterrés », guerres de leadership. \\ \hline
\end{tabularx}
\end{center}

\begin{exoresolu}[Lecture critique de données : le rendement de la chaîne anti-corruption]
\textbf{Énoncé :} Un rapport annuel fictif mais réaliste d'un organe anti-corruption donne les chiffres suivants :
\begin{itemize}
\item Dossiers reçus et traités : \textbf{2 500} (dont 60\,\% de dénonciations de citoyens, 20\,\% d'auto-saisines, 20\,\% issues de la presse et des ONG).
\item Dossiers transmis à la justice : \textbf{150}.
\item Condamnations prononcées : \textbf{50}.
\item Montants récupérés ou ordonnancés : \textbf{2 milliards de FCFA}.
\end{itemize}
\begin{enumerate}
\item Calculer le taux de transmission à la justice par rapport aux dossiers traités.
\item Calculer le taux de condamnation final par rapport aux signalements initiaux.
\item Interpréter l'écart entre les dossiers traités (2 500) et les dossiers transmis (150).
\end{enumerate}
\tcblower
\textbf{Solution détaillée :}
\begin{enumerate}
\item Taux de transmission $= \dfrac{150}{2500} \times 100 = \mathbf{6\,\%}$.
\item Taux de condamnation final $= \dfrac{50}{2500} \times 100 = \mathbf{2\,\%}$.
\item \textbf{Interprétation :} 94\,\% des dossiers s'arrêtent au stade administratif (classement sans suite, insuffisance de preuves, prescription, incompétence de l'organe). Cela révèle la difficulté de transformer une dénonciation en preuve judiciaire recevable. Les 2 milliards récupérés sont réels mais restent marginaux face au coût estimé de la corruption pour l'économie nationale, évalué par les institutions internationales à plusieurs points du PIB, soit des centaines de milliards de FCFA par an.
\end{enumerate}
\end{exoresolu}

\courssub{Vers une intégrité systémique : la responsabilité du technicien}

En tant qu'élèves de Première technique, futurs techniciens supérieurs, gestionnaires de chantiers, de réseaux, de systèmes d'information ou de production, vous serez les \textbf{premiers opérateurs} de l'intégrité sur le terrain. Le dispositif institutionnel n'est efficace que si les acteurs de terrain refusent le « système D » et utilisent les leviers à leur disposition.

\begin{methode}[Réflexe « Intégrité technique » face à une situation à risque]
\begin{enumerate}
\item \textbf{Identifier :} est-ce un détournement de matériel ? une fausse facture ? un rapport de chantier falsifié ? un logiciel piraté ? un pot-de-vin pour valider un procès-verbal ?
\item \textbf{Qualifier :} quelle infraction ? (corruption, concussion, détournement, faux, escroquerie ?).
\item \textbf{Se protéger :} ne pas affronter seul le corrupteur ; conserver les preuves (photos, courriels, PV signés, captures d'écran) hors de sa portée.
\item \textbf{Signaler :} hiérarchie directe (par écrit, avec accusé de réception) $\rightarrow$ en cas d'inaction ou de complicité : CONAC (8200) $\rightarrow$ médias et ONG en dernier recours.
\item \textbf{Documenter :} tenir un journal de bord personnel (daté, signé) des événements, des refus et des pressions subies : ce sera une preuve devant le Conseil de Prud'hommes ou le Tribunal administratif.
\end{enumerate}
\end{methode}

\begin{exercice}[Application : le cas du stage en entreprise]
Lors de ton stage de fin d'année de Première technique (Génie civil), tu découvres que le chef de chantier utilise un ciment de qualité inférieure à celle prescrite par le marché, pour « économiser » et partager la différence de prix avec le fournisseur. Le béton coulé pour les poteaux du bâtiment ne respecte pas les normes de résistance prévues par l'étude de béton armé.
\begin{enumerate}
\item Quels sont les risques techniques immédiats et différés (résistance, durabilité, sécurité des usagers) ?
\item Quelles infractions pénales et administratives sont caractérisées ?
\item Quelle démarche concrète suis-tu (hiérarchie, CONAC, ordre professionnel, médias) pour stopper le chantier sans te mettre en danger ?
\item Rédige le courrier type de signalement à la CONAC (objet, faits, preuves, demande).
\end{enumerate}
\end{exercice}

\begin{corrige}[Exercice : le cas du stage en entreprise]
\begin{enumerate}
\item \textbf{Risques :} résistance des poteaux inférieure aux hypothèses de calcul, fissuration précoce, corrosion accélérée des armatures (béton plus poreux), risque d'effondrement partiel ou total sous les charges d'exploitation ou en cas de séisme. Le stagiaire qui se tait alors qu'il a constaté les faits engage sa responsabilité morale et peut voir sa responsabilité juridique recherchée.
\item \textbf{Infractions :} mise en danger de la vie d'autrui ; escroquerie (livraison non conforme au marché) ; corruption active et passive (entente chef de chantier / fournisseur) ; faux et usage de faux (PV de réception falsifié) ; violation des normes techniques et des règles de l'art.
\item \textbf{Démarche :} 1. Refus ferme et oral de signer le PV de réception. 2. Écrit au chef de chantier, au maître de stage (entreprise) et au chef de mission de stage (école) : « constat de non-conformité, refus de valider, mise en demeure de corriger ». 3. En cas de refus ou de pression : signalement à la CONAC (8200, photos des sacs de ciment, bons de livraison, notes de calcul). 4. Signalement à l'ordre professionnel compétent si le chef de chantier est un ingénieur inscrit. 5. Alerte aux médias uniquement en cas de danger imminent et d'inaction totale des autorités.
\item \textbf{Courrier type :} \emph{Objet :} signalement de corruption et de non-conformité sur le chantier [nom], localisation [adresse]. \emph{Faits :} dates, description technique (ciment livré vs ciment prescrit), noms des auteurs, montant estimé du détournement. \emph{Preuves jointes :} photos, bons de commande et de livraison, notes de calcul. \emph{Demande :} enquête urgente, suspension conservatoire des travaux, protection de l'identité du signataire. \emph{Signature} (ou dénonciation anonyme).
\end{enumerate}
\end{corrige}

\begin{aretenir}[Synthèse : le dispositif est un écosystème, pas une liste]
L'intégrité au Cameroun ne repose pas sur UNE institution miracle, mais sur l'\textbf{interaction vertueuse} de cinq piliers spécialisés (CONAC, CONSUPE, Chambre des Comptes, CNDHL, Commission de la Déclaration des Biens et Avoirs), relayés par une justice indépendante, stimulés par une société civile vigilante, et \textbf{incarnés au quotidien par chaque citoyen-technicien}. Connaître cet organigramme, ses forces (légalité, pouvoirs d'enquête, ancrage dans les textes) et ses faiblesses (moyens, indépendance relative, lenteur judiciaire), c'est déjà exercer sa citoyenneté technique. La section suivante (TD2) nous plongera dans un laboratoire concret : ton établissement scolaire.
\end{aretenir}

\coursec{TD2 : L'intégrité morale en milieu scolaire -- Diagnostic et Analyse}

Après avoir découvert, dans la section précédente, les institutions camerounaises chargées de promouvoir et de préserver l'intégrité morale (CONAC, Commission nationale anti-corruption, tribunaux, cellules d'éthique des ministères), une question naturelle se pose : \emph{que se passe-t-il concrètement dans notre propre cadre de vie, l'école ?} Car l'intégrité ne se décrète pas seulement dans les textes officiels : elle se vit ou se trahit chaque jour dans les salles de classe, les salles d'examen et les bureaux d'administration. Ce travail dirigé (TD2) nous invite à devenir de véritables \cle{enquêteurs citoyens} : à partir d'un dossier documentaire, nous allons diagnostiquer les atteintes à l'intégrité morale en milieu scolaire, en analyser les causes et les conséquences, puis construire collectivement un outil d'analyse.

\courssub{La méthode d'analyse de documents}

Avant de plonger dans le dossier, il faut se doter d'une démarche rigoureuse. En ECM comme dans les matières techniques, on ne juge pas un problème « au feeling » : on l'examine méthodiquement, comme un technicien diagnostique une panne avant de réparer.

\begin{methode}[Analyser un document en TD]
Pour exploiter correctement un document (texte, statistique, image, texte de loi), suivre cinq étapes :
\begin{enumerate}
\item \textbf{Identifier la nature du document} : s'agit-il d'un texte argumentatif, d'un tableau statistique, d'une photographie, d'un extrait de loi ou de règlement ?
\item \textbf{Extraire l'information clé} : relever le chiffre marquant, le fait rapporté, l'argument central ou la sanction prévue.
\item \textbf{Contextualiser} : préciser la date, la source et l'auteur du document. Un rapport de la CONAC n'a pas la même portée qu'un message anonyme sur WhatsApp.
\item \textbf{Croiser avec d'autres documents} : confronter les informations entre elles pour vérifier leur cohérence et dégager une tendance.
\item \textbf{Conclure et proposer une solution} : formuler un diagnostic argumenté et suggérer une piste d'action réaliste.
\end{enumerate}
\end{methode}

\courssub{Présentation et analyse du dossier documentaire}

Le dossier TD comprend quatre documents que la classe analyse en groupes de travail.

\textbf{Document 1 : Extrait d'un rapport de la CONAC sur la corruption en milieu scolaire.} Ce document statistique recense les pratiques corruptives les plus fréquentes dans le système éducatif camerounais : \cle{achat de notes} (un élève ou ses parents versent de l'argent à un enseignant pour obtenir une meilleure moyenne), \cle{fuites de sujets d'examens} (des épreuves officielles circulent avant la date de composition) et \cle{faux diplômes} (fabrication ou achat d'attestations de réussite sans avoir composé). La source est institutionnelle et donc fiable ; les chiffres montrent que ces pratiques touchent tous les niveaux, du primaire à l'enseignement supérieur.

\textbf{Document 2 : Témoignages d'élèves.} Plusieurs élèves racontent des pressions subies : un enseignant qui conditionne une note à un « cadeau », une élève victime de \cle{harcèlement sexuel} à qui l'on promet des points en échange de faveurs. Ces témoignages, bien qu'anonymes, croisent le document 1 : ils montrent que la corruption scolaire n'est pas qu'une question d'argent, mais aussi d'\emph{abus de pouvoir}.

\textbf{Document 3 : Extrait d'un règlement intérieur type.} Ce texte réglementaire prévoit les sanctions en cas de tricherie : zéro à l'évaluation, retenue, exclusion temporaire, voire renvoi définitif en cas de récidive. Il prouve que l'école dispose déjà d'un arsenal disciplinaire -- encore faut-il l'appliquer.

\textbf{Document 4 : Photographies d'affichages « Zéro Tolérance » et de boîtes à idées/signalement.} Ces images montrent que certains établissements mettent en place des dispositifs concrets de prévention et de dénonciation. Elles illustrent la possibilité d'agir localement, sans attendre les grandes institutions.

\begin{exemplebox}[Croisement des documents]
En croisant les documents 1 et 2, on constate que les chiffres officiels de la CONAC correspondent aux faits vécus par les élèves : la corruption scolaire est une réalité mesurée \emph{et} ressentie. En croisant les documents 3 et 4, on observe un écart : les règlements existent, mais leur application dépend de la volonté des chefs d'établissement et de la confiance des élèves dans les mécanismes de signalement.
\end{exemplebox}

\courssub{Identifier les formes de corruption scolaire}

L'analyse du dossier permet de classer les atteintes à l'intégrité morale en milieu scolaire en quatre grandes catégories.

\begin{definition}[Les formes de corruption scolaire]
\begin{itemize}
\item La \cle{corruption active} : c'est celui qui \emph{propose} qui faute. Un élève ou un parent offre de l'argent, un cadeau ou un avantage à un enseignant pour obtenir une note, un sujet ou une faveur.
\item La \cle{corruption passive} : c'est celui qui \emph{exige ou accepte} qui faute. Un enseignant demande ou reçoit un avantage en échange d'un acte de sa fonction (bonne note, indulgence, admission).
\item Le \cle{trafic d'influence} : une personne abuse de ses relations (connaissance du directeur, parenté avec un responsable) pour obtenir un avantage indu, comme une admission ou un classement favorable.
\item La \cle{fraude aux examens} : toute manœuvre visant à fausser les résultats d'une évaluation : usage du téléphone portable, « cahiers de brouillon » pré-remplis, substitution de candidat (quelqu'un compose à la place d'un autre).
\end{itemize}
\end{definition}

\begin{definition}[Fraude aux examens et sanctions légales]
La \cle{fraude aux examens} est toute manœuvre frauduleuse visant à fausser les résultats d'un examen ou d'un concours officiel. Elle est réprimée par le droit camerounais (article 141 du Code pénal ; loi n\textsuperscript{o} 98/004 du 14 avril 1998 d'orientation de l'éducation au Cameroun). Les sanctions prévues sont :
\begin{itemize}
\item le zéro à l'épreuve et l'annulation des résultats ;
\item l'exclusion de tout examen officiel pour une durée de un à cinq ans ;
\item une peine de prison de trois mois à trois ans, assortie d'une amende, pour les cas les plus graves (fuites de sujets, faux diplômes, substitution de candidat).
\end{itemize}
\end{definition}

\begin{attention}[Le piège de la « petite corruption »]
Le cadeau au professeur « pour son anniversaire », la « motivation » glissée au surveillant, le « petit quelque chose » au censeur : ces gestes semblent anodins, mais ils \emph{banalisent l'acte corruptif}. L'école de la petite corruption est l'antichambre de la grande corruption. Le proverbe le dit clairement : \emph{« Qui vole un \oe uf vole un b\oe uf »}. Celui qui achète une note en classe de Première achètera demain un marché public ou un jugement.
\end{attention}

\begin{savaistu}[Le « piston », une corruption déguisée]
Le \cle{piston} (ou népotisme) utilisé pour l'admission dans une grande école ou dans la fonction publique est une forme de trafic d'influence : il combine corruption active (celui qui sollicite la relation) et corruption passive (celui qui use de son pouvoir pour favoriser un proche). Conséquence directe : un candidat méritant, qui a travaillé et réussi honnêtement, est privé de sa place. Le piston n'est donc pas une « solidarité » : c'est un vol d'opportunité.
\end{savaistu}

\courssub{Analyser les causes de la corruption scolaire}

Pourquoi des élèves, des parents et même des enseignants en arrivent-ils à corrompre ou à tricher ? L'analyse du dossier et la discussion en classe dégagent cinq causes principales.

\begin{enumerate}
\item \textbf{La pression de la réussite} : la famille et la société valorisent le diplôme plus que la compétence. « Réussis coûte que coûte » devient un mot d'ordre qui justifie tous les moyens.
\item \textbf{La précarité des enseignants} : des salaires modestes, parfois payés avec retard, peuvent fragiliser certains personnels et les exposer à la tentation de la corruption passive. Cela n'excuse rien, mais cela explique une partie du phénomène.
\item \textbf{L'impunité antérieure} : lorsque des fraudeurs ont été pris sans être sanctionnés, le message envoyé est que « ça passe ». L'absence de sanction nourrit la récidive.
\item \textbf{La défaillance de la surveillance} : salles mal organisées, surveillants insuffisants ou complaisants, téléphones non confisqués : la fraude prospère là où le contrôle faiblit.
\item \textbf{La culture du « piston » et de la « connaissance »} : dans certains milieux, on estime qu'on ne réussit rien sans relation. Cette mentalité décourage le mérite et légitime le trafic d'influence.
\end{enumerate}

\courssub{Analyser les conséquences : un danger pour l'individu et pour la nation}

La corruption scolaire n'est pas un « crime sans victime ». Ses conséquences sont lourdes et durables.

\begin{itemize}
\item \textbf{L'incompétence des diplômés} : un diplôme acheté certifie des savoir-faire que le titulaire ne possède pas. Dans les filières techniques, c'est un \emph{danger réel} : un électricien incompétent provoque des incendies, un maçon incompétent construit des bâtiments qui s'effondrent, un infirmier incompétent met des vies en péril.
\item \textbf{L'injustice sociale} : la corruption brise la \cle{méritocratie}, c'est-à-dire le principe selon lequel chacun obtient une place selon son travail et son talent. L'élève pauvre mais travailleur est évincé par l'élève riche mais paresseux.
\item \textbf{La dévalorisation du diplôme camerounais} : lorsque les fraudes se multiplient, les employeurs et les universités étrangères doutent de la valeur des diplômes du pays. La reconnaissance internationale de nos certifications s'affaiblit, et ce sont \emph{tous} les diplômés, y compris les honnêtes, qui en paient le prix.
\item \textbf{La perte de confiance dans l'école} : une école corrompue ne forme plus, elle décourage. Les familles, les élèves et les enseignants intègres perdent foi dans l'institution censée construire l'avenir.
\end{itemize}

\begin{exemplebox}[Le diplômé dangereux]
Un élève de Première technique F4 (Électrotechnique) achète le sujet de l'épreuve pratique du baccalauréat blanc pour \num{5000} FCFA. Il « réussit » son Bac, décroche son diplôme et ouvre un atelier d'installation électrique. Mais il ne sait pas câbler correctement un tableau de distribution. Quelques mois plus tard, une installation mal protégée provoque un incendie chez un client : dégâts matériels, blessés, procès, et finalement une condamnation à la prison. Les \num{5000} FCFA de la fraude se transforment en millions de francs de dommages et en années de détention. La note achetée n'a jamais remplacé la compétence.
\end{exemplebox}

\courssub{Étude de cas : « L'affaire des sujets du Probatoire 202X »}

\begin{exoresolu}[Étude de cas : la fuite des sujets du Probatoire 202X]
\textbf{Énoncé.} À la veille de l'épreuve de Mathématiques du Probatoire technique, le sujet circule sur des groupes WhatsApp d'élèves de plusieurs lycées de la région. L'enquête administrative révèle les faits suivants : un imprimeur chargé de reproduire les sujets a photographié l'épreuve et l'a vendue à un surveillant ; celui-ci l'a revendue à un délégué de classe ; le délégué l'a partagée aux élèves contre \num{2000} FCFA par copie ; certains parents, informés, ont payé pour leurs enfants. Les imprimeurs et surveillants impliqués sont arrêtés et déférés devant les tribunaux.

\textbf{Questions.} 1) Identifier les maillons de la chaîne de complicité. 2) Pour chaque maillon, préciser la forme de corruption commise. 3) Quelles sanctions encourent les responsables ? 4) Proposer deux mesures pour empêcher une telle fuite.
\tcblower
\textbf{Solution détaillée.}
\begin{enumerate}
\item \textbf{Chaîne de complicité} : imprimeur $\rightarrow$ surveillant $\rightarrow$ délégué de classe $\rightarrow$ élèves $\rightarrow$ parents complices.
\item \textbf{Qualification des fautes} : l'imprimeur commet une \emph{corruption active} (il vend un document confidentiel) et une violation du secret professionnel ; le surveillant commet une \emph{corruption passive} puis active (il achète puis revend) ; le délégué joue le rôle d'\emph{intermédiaire corrupteur} ; les élèves qui achètent commettent une \emph{corruption active} et une fraude aux examens ; les parents qui paient sont \emph{complices} de la fraude.
\item \textbf{Sanctions encourues} : annulation des épreuves et zéro pour les candidats fraudeurs ; exclusion des examens officiels de un à cinq ans ; pour les organisateurs de la fuite (imprimeur, surveillant), une peine de prison de trois mois à trois ans et une amende, conformément à l'article 141 du Code pénal, sans compter les sanctions disciplinaires (révocation pour les personnels).
\item \textbf{Mesures préventives} : renforcer la sécurité de la chaîne d'impression et de transport des sujets (scellés numérotés, convoyage sous escorte, responsabilisation nominative) ; instaurer des sanctions systématiquement appliquées et médiatisées pour briser le sentiment d'impunité, tout en sensibilisant les élèves et les parents avant chaque session d'examen.
\end{enumerate}
\end{exoresolu}

Pour visualiser cette affaire, on peut représenter la mécanique générale de toute corruption scolaire sous forme de chaîne.

\begin{popfigure}
\begin{center}
\begin{tikzpicture}[
box/.style={draw=popblue, thick, fill=popblueL, rounded corners, text width=2.6cm, align=center, minimum height=1.4cm, font=\small},
badbox/.style={draw=poporange, thick, fill=poporangeL, rounded corners, text width=2.6cm, align=center, minimum height=1.4cm, font=\small},
arr/.style={-{Stealth[length=3mm]}, very thick, popdark}]
\node[badbox, align=center] (dem) at (0,0){\textbf{Demandeur}\\ Parent / Élève};
\node[box, align=center] (int) at (3.6,0){\textbf{Intermédiaire}\\ Délégué / Surveillant};
\node[badbox, align=center] (pouv) at (7.2,0){\textbf{Détenteur du pouvoir}\\ Prof / Imprimeur / Officier};
\node[box, align=center] (ben) at (10.8,0){\textbf{Bénéfice indu}\\ Note / Diplôme};
\node[badbox, fill=poporange!40, align=center] (perte) at (7.2,-2.6){\textbf{Perte collective}\\ Qualité de la formation détruite};
\draw[arr] (dem) -- (int) node[midway, above, font=\scriptsize]{argent, cadeau};
\draw[arr] (int) -- (pouv) node[midway, above, font=\scriptsize]{commission};
\draw[arr] (pouv) -- (ben) node[midway, above, font=\scriptsize]{faveur};
\draw[arr, poporange] (ben.south) |- (perte.east);
\draw[arr, poporange] (dem.south) |- (perte.west);
\end{tikzpicture}
\end{center}
\legende{La chaîne de la corruption scolaire : chaque maillon profite individuellement, mais la collectivité entière subit la perte finale.}
\end{popfigure}

\courssub{Lire les statistiques : évolution des fraudes signalées}

Les rapports des organismes chargés des examens officiels (Office du Baccalauréat du Cameroun, MINESEC) permettent de suivre l'évolution des cas de fraude signalés. Le graphique ci-dessous présente une évolution indicative sur plusieurs sessions.

\begin{popfigure}
\begin{center}
\begin{tikzpicture}
\begin{axis}[
ybar, bar width=16pt,
width=0.85\linewidth, height=6.5cm,
xlabel={Session d'examen}, ylabel={Cas de fraude signalés},
symbolic x coords={2018,2019,2020,2021,2022,2023},
xtick=data,
ymin=0, ymax=1100,
nodes near coords, nodes near coords style={font=\scriptsize},
ymajorgrids=true, grid style={popline!40},
every axis plot/.append style={fill=popblue, draw=popdark}]
\addplot[fill=popblue, draw=popdark] coordinates {(2018,620) (2019,580) (2020,340) (2021,950) (2022,710) (2023,480)};
\end{axis}
\end{tikzpicture}
\end{center}
\legende{Évolution indicative des cas de fraude signalés aux examens officiels (sessions 2018-2023). Source : OBC / MINESEC (données reconstituées à des fins pédagogiques). On observe un pic post-COVID en 2021 (reprise des examens après les perturbations sanitaires), puis une baisse liée au renforcement des mesures de sécurisation.}
\end{popfigure}

La lecture du graphique appelle deux remarques. D'une part, le pic de 2021 montre que les périodes de désorganisation (crise sanitaire, reports de sessions) créent des failles exploitées par les fraudeurs : entre 2020 et 2021, les cas signalés sont passés de 340 à 950, soit une hausse d'environ $180\,\%$ en un an. D'autre part, la baisse des années suivantes (de 950 cas en 2021 à 480 cas en 2023, soit une diminution de près de $50\,\%$ en deux ans) prouve que les mesures de prévention (sécurisation des sujets, sanctions médiatisées, sensibilisation) \emph{fonctionnent} lorsqu'elles sont appliquées avec constance : la corruption n'est pas une fatalité.

\courssub{Construction collective : l'arbre des causes et des effets de la tricherie}

Pour clore le diagnostic, la classe construit au tableau un \cle{arbre des causes et des effets}, inspiré de la méthode d'Ishikawa (diagramme en arêtes de poisson utilisé en gestion de qualité dans les entreprises). Le principe est simple : le tronc représente le problème central (la tricherie), les racines représentent ses causes profondes et les branches ses conséquences.

\begin{popfigure}
\begin{center}
\begin{tikzpicture}[
root/.style={draw=popgreen, thick, fill=popgreenL, rounded corners, text width=3.1cm, align=center, font=\scriptsize},
trunk/.style={draw=popdark, very thick, fill=popgoldL, rounded corners, text width=3.2cm, align=center, font=\small\bfseries},
branch/.style={draw=poporange, thick, fill=poporangeL, rounded corners, text width=3.1cm, align=center, font=\scriptsize},
link/.style={thick, popmuted}]
\node[trunk] (t) at (0,0) {LA TRICHERIE en milieu scolaire};
\node[root] (c1) at (-4.5,-2.6) {Pression de la réussite (famille, société)};
\node[root] (c2) at (-1.5,-3.2) {Précarité de certains enseignants};
\node[root] (c3) at (1.5,-3.2) {Impunité et sanctions non appliquées};
\node[root] (c4) at (4.5,-2.6) {Culture du « piston » et défaillance de surveillance};
\node[branch] (e1) at (-4.5,2.6) {Diplômés incompétents (danger technique)};
\node[branch] (e2) at (-1.5,3.2) {Injustice sociale, mérite bafoué};
\node[branch] (e3) at (1.5,3.2) {Dévalorisation du diplôme camerounais};
\node[branch] (e4) at (4.5,2.6) {Perte de confiance dans l'école};
\draw[link] (c1.north) -- (t.south west);
\draw[link] (c2.north) -- (t.south);
\draw[link] (c3.north) -- (t.south);
\draw[link] (c4.north) -- (t.south east);
\draw[link] (e1.south) -- (t.north west);
\draw[link] (e2.south) -- (t.north);
\draw[link] (e3.south) -- (t.north);
\draw[link] (e4.south) -- (t.north east);
\end{tikzpicture}
\end{center}
\legende{Arbre des causes (racines) et des effets (branches) de la tricherie, construit collectivement selon la méthode d'Ishikawa. Couper les racines, c'est assécher le problème.}
\end{popfigure}

\begin{methode}[Construire un arbre causes-effets en classe]
\begin{enumerate}
\item Écrire le problème central au milieu du tableau (le tronc).
\item En dessous, faire émerger les \textbf{causes} par brainstorming : chaque groupe propose une cause issue de l'analyse du dossier ; on ne retient que celles justifiées par un document.
\item Au-dessus, inscrire les \textbf{conséquences} en distinguant celles qui touchent l'individu, la communauté scolaire et la nation entière.
\item Relier chaque cause et chaque effet au tronc, puis discuter : quelle racine peut-on couper en priorité, et par quelle action concrète ?
\item Conclure en dégageant une ou deux solutions réalisables à l'échelle de l'établissement (charte d'honnêteté, boîte de signalement, sensibilisation avant les examens).
\end{enumerate}
\end{methode}

\begin{aretenir}[L'essentiel du TD2]
\begin{itemize}
\item La corruption scolaire prend quatre formes : corruption active, corruption passive, trafic d'influence et fraude aux examens.
\item Elle s'explique par la pression de la réussite, la précarité, l'impunité, les défaillances de surveillance et la culture du piston.
\item Ses conséquences dépassent l'école : diplômés incompétents et dangereux, injustice sociale, dévalorisation du diplôme national, perte de confiance dans l'école.
\item La fraude aux examens est un délit puni par la loi (article 141 du Code pénal ; loi n\textsuperscript{o} 98/004 du 14 avril 1998) : zéro à l'épreuve, exclusion de un à cinq ans, prison et amende pour les cas graves.
\item L'analyse de documents et l'arbre causes-effets sont des outils de diagnostic citoyen : comprendre les racines du mal pour agir efficacement.
\end{itemize}
\end{aretenir}

\begin{exercice}[À faire en classe ou à la maison]
\begin{enumerate}
\item À partir du document 1 (rapport CONAC), rédige en dix lignes un paragraphe argumenté répondant à la question : « Pourquoi la corruption scolaire est-elle une menace pour l'avenir professionnel des élèves de l'enseignement technique ? »
\item Dans l'affaire des sujets du Probatoire 202X, rédige la lettre qu'un délégué de classe intègre pourrait adresser au proviseur pour signaler la fraude sans s'exposer.
\item Complète l'arbre causes-effets construit en classe en ajoutant une cause et une conséquence non mentionnées, en les justifiant.
\end{enumerate}
\end{exercice}

\begin{corrige}[Exercice n\textsuperscript{o} 1]
Éléments de réponse attendus. \textbf{Question 1} : le paragraphe doit mobiliser au moins deux conséquences (incompétence dangereuse dans les métiers techniques, dévalorisation du diplôme auprès des employeurs) et s'appuyer sur un chiffre ou un fait du document. \textbf{Question 2} : la lettre doit respecter la structure d'un courrier administratif (objet, faits précis, demande de confidentialité, formules de politesse) et montrer que le signalement est un acte de courage civique protégé, non une délation. \textbf{Question 3} : toute cause ou conséquence est acceptée si elle est justifiée par un fait observable (exemple de cause : la banalisation des cadeaux aux enseignants ; exemple de conséquence : la fuite des élèves doués vers les filières ou les pays où le mérite est reconnu).
\end{corrige}

\coursec{TP : Construction d'outils de promotion de l'intégrité (Cartes, Croquis, Chartes)}

\noindent Après avoir diagnostiqué les atteintes à l'intégrité dans notre établissement lors du TD2 (analyse de cas, questionnaire, débat), nous passons à l'action. Comprendre les mécanismes de la corruption ou du harcèlement ne suffit plus ; il faut \emph{outiller} la communauté éducative pour les prévenir, les signaler et les sanctionner. Ce TP transforme l'élève en \cle{acteur de l'intégrité} : cartographe des risques, concepteur de procédures, législateur de sa classe et analyste de données. Chaque outil produit ici a vocation à être affiché, diffusé ou intégré au règlement intérieur.

\courssub{TP1 -- Cartographie des risques : la « Carte de vigilance intégrité »}

\begin{methode}[Réaliser un croquis géographique scolaire]
\begin{enumerate}
    \item \textbf{Cadre et échelle :} délimiter le périmètre (enceinte du lycée). Choisir une échelle approximative (ex. : 1 cm pour 10 m) pour tenir sur une feuille A4 ou A3.
    \item \textbf{Structures principales :} tracer au crayon les bâtiments (administration, salles de classe, ateliers, réfectoire, toilettes, portail, mur d'enceinte). Indiquer l'orientation (flèche du Nord).
    \item \textbf{Thématisation (risques) :} superposer les symboles de risque selon une légende codée, avec des couleurs (rouge / orange / vert) ou des pastilles.
    \item \textbf{Légende organisée :} classer par nature de risque (financier, sexuel, académique, disciplinaire) et par niveau de criticité (critique, surveillé, sain).
    \item \textbf{Métadonnées :} titre précis, date de réalisation, auteur(s) (groupe/classe), source (enquête de terrain).
    \item \textbf{Propreté :} règle, équerre, écriture lisible, pas de ratures. Numérisation possible (photo/scanner) pour affichage.
\end{enumerate}
\end{methode}

\begin{experience}[Enquête de terrain pour la cartographie]
\textbf{Objectif :} identifier les « zones chaudes » où l'intégrité est le plus menacée.

\textbf{Matériel :} plan simplifié du lycée (fourni par l'intendance ou dessiné à main levée), grille d'observation (4 colonnes : Lieu, Risque principal, Fréquence estimée, Témoins), crayons de couleur, règle.

\textbf{Protocole :}
\begin{enumerate}
    \item En binômes, parcourir l'établissement aux heures creuses (récréation, intercours, sortie).
    \item Interroger 3 à 5 élèves ou agents par zone (anonymat garanti) : « As-tu déjà vu ou subi une situation anormale ici ? » (argent contre note, harcèlement, racket, fuite de sujets).
    \item Reporter sur le plan les symboles : \textbf{F} (caisse/intendance, risque financier), \textbf{A} (secrétariat des examens, risque académique), \textbf{D} (toilettes/recoins, risque disciplinaire ou harcèlement), \textbf{S} (portail, risque sécurité).
    \item Qualifier la criticité : \textbf{Rouge} (faits avérés et récurrents), \textbf{Orange} (rumeurs, configuration propice), \textbf{Vert} (rien à signaler, zone bien surveillée).
\end{enumerate}

\textbf{Observations types :} bureau de la caisse : paiements « accélérés » sans reçu ; secrétariat des examens : photocopies de sujets circulant avant l'heure ; toilettes : racket par de grands élèves, consommation de stupéfiants ; portail : faux badges, entrées non autorisées.

\textbf{Interprétation :} la carte révèle que le risque \emph{financier} et \emph{académique} se concentre dans l'enceinte administrative (accès restreint), tandis que le risque de \emph{harcèlement} et de \emph{racket} prospère dans les zones mal surveillées (toilettes, arrière-bâtiments).
\end{experience}

\begin{popfigure}
\begin{tikzpicture}[scale=0.72, every node/.style={font=\small}]
    \tikzset{
        bat/.style={draw=popdark, fill=popblueL, thick, rectangle, rounded corners=3pt, minimum width=2.5cm, minimum height=1.5cm, align=center},
        zoneR/.style={circle, fill=poppink, minimum size=9pt, inner sep=0},
        zoneO/.style={circle, fill=poporange, minimum size=9pt, inner sep=0},
        zoneV/.style={circle, fill=popgreen, minimum size=9pt, inner sep=0},
        leg/.style={rectangle, draw=popline, fill=popcream, rounded corners=2pt, inner sep=4pt, align=left, font=\footnotesize}
    }
    % Fond quadrillé léger
    \draw[popline!50, step=1cm] (0,0) grid (14,10);
    % Murs d'enceinte
    \draw[very thick, popdark] (0,0) rectangle (14,10);
    % Bâtiments
    \node[bat, align=center] (Admin) at (2,8){ADMINISTRATION\\Proviseur / Caisse / Secrétariat};
    \node[bat, align=center] (SallesA) at (7,8){BÂTIMENT A\\Salles 1--10};
    \node[bat, align=center] (SallesB) at (12,8){BÂTIMENT B\\Salles 11--20};
    \node[bat, align=center] (Ateliers) at (2,4){ATELIERS\\Techniques / Labos};
    \node[bat, align=center] (Refect) at (7,4){RÉFECTOIRE\\Salle de détente};
    \node[bat, align=center] (Toilettes) at (12,4){TOILETTES\\Bloc sanitaire};
    \node[bat, minimum width=1.8cm, minimum height=1cm, align=center] (Gardien) at (1.5,1){LOGEMENT\\GARDIEN};
    % Portail
    \draw[very thick, popcream] (6.3,0) -- (7.7,0);
    \draw[very thick, popdark] (6.3,0) -- (6.3,-0.5) -- (7.7,-0.5) -- (7.7,0);
    \node[below, font=\footnotesize] at (7,-0.5) {PORTAIL PRINCIPAL};
    % Zones à risque
    \node[zoneR] at (1.2,8.6) {};
    \node[above, font=\tiny\bfseries, text=poppink] at (1.2,8.7) {F};
    \node[zoneO] at (2.9,8.6) {};
    \node[above, font=\tiny\bfseries, text=poporange] at (2.9,8.7) {A};
    \node[zoneR] at (11.4,3.4) {};
    \node[below, font=\tiny\bfseries, text=poppink] at (11.4,3.3) {D};
    \node[zoneO] at (7,0.6) {};
    \node[above, font=\tiny\bfseries, text=poporange] at (7,0.7) {S};
    \node[zoneO] at (0.8,2.6) {};
    \node[right, font=\tiny\bfseries, text=poporange] at (0.9,2.6) {D};
    % Zones vertes
    \node[zoneV] at (4.5,6.3) {};
    \node[zoneV] at (10,6.3) {};
    % Légende
    \node[leg, anchor=north east, align=center] at (13.8,9.8){
        \textbf{LÉGENDE -- CARTE DE VIGILANCE}\\
        \textbf{Criticité :}\\
        \tikz\draw[fill=poppink] (0,0) circle(3pt); \textbf{Rouge} : risque avéré, récurrent\\
        \tikz\draw[fill=poporange] (0,0) circle(3pt); \textbf{Orange} : risque potentiel\\
        \tikz\draw[fill=popgreen] (0,0) circle(3pt); \textbf{Vert} : zone saine, encadrée\\
        \textbf{Nature du risque :}\\
        F : financier (détournement, faux reçus)\\
        A : académique (fuite de sujets, notes)\\
        D : disciplinaire (harcèlement, racket)\\
        S : sécurité (faux badges, intrusion)\\
        \textit{Lycée X -- 1re Tech -- Octobre 2025}\\
        \textit{Échelle indicative : 1 cm $\approx$ 10 m}
    };
    % Flèche Nord
    \draw[->, thick, popdark] (13,8.6) -- (13,9.4) node[above] {N};
\end{tikzpicture}
\legende{Croquis simplifié « Carte de vigilance intégrité -- Lycée X » (orientation : flèche du Nord ; échelle indicative). Les zones rouges (caisse, toilettes) nécessitent une action immédiate (éclairage, présence adulte, boîte à idées). Les zones oranges (secrétariat, portail, recoins) demandent un renforcement des procédures (contrôle d'accès, rondes).}
\end{popfigure}

\begin{aretenir}[Intérêt pédagogique de la carte]
Le croquis spatialise l'abstrait : il transforme le « on dit que... » en « ici, précisément, tel risque existe ». C'est un outil de dialogue avec l'administration (demande d'éclairage, de rondes, de caméras) et un support de prévention pour les nouveaux élèves (visite guidée « sécurité et intégrité » à la rentrée).
\end{aretenir}

\courssub{TP2 -- Diagramme de flux : procédure de signalement « Du témoin à la sanction »}

\begin{definition}[Organigramme / diagramme de flux]
Représentation graphique d'un processus séquentiel. \textbf{Ovale} = début / fin ; \textbf{rectangle} = action, tâche ou acteur ; \textbf{losange} = décision ou test (Oui/Non) ; \textbf{flèche} = flux (étape suivante). Il permet de visualiser les goulots d'étranglement, les boucles de retour et les points de vigilance (délais, anonymat).
\end{definition}

\begin{savaistu}[Logiciels libres pour réaliser des diagrammes]
\begin{itemize}
    \item \textbf{Draw.io (diagrams.net) :} en ligne, gratuit, export PNG/SVG/PDF, formes normalisées, travail collaboratif.
    \item \textbf{Dia :} installable (Windows/Linux), léger, export EPS/SVG.
    \item \textbf{LibreOffice Draw :} intégré à la suite bureautique, suffisant pour des organigrammes simples.
\end{itemize}
Ces outils permettent de numériser la carte du TP1 et l'organigramme du TP2 pour un affichage propre et durable.
\end{savaistu}

\begin{popfigure}
\begin{tikzpicture}[
    node distance=1.1cm and 2.2cm,
    startstop/.style={ellipse, draw=popblue, thick, fill=popblueL, minimum width=2.5cm, minimum height=1cm, align=center},
    process/.style={rectangle, draw=popdark, thick, fill=popcream, minimum width=3cm, minimum height=1cm, align=center, rounded corners=3pt},
    decision/.style={diamond, draw=poporange, thick, fill=poporangeL, minimum width=2.4cm, minimum height=1.1cm, align=center, aspect=1.6},
    arrow/.style={->, thick, popdark},
    vigilance/.style={rectangle, draw=poppink, thick, fill=poppinkL, minimum width=3.2cm, minimum height=0.8cm, align=center, font=\footnotesize\itshape, rounded corners=2pt}
]
% Nœuds
\node[startstop, align=center] (start){DÉBUT\\Témoin / Victime};
\node[process, below of=start, align=center] (signal){SIGNALEMENT\\Boîte à idées / Délégué / CPE\\Proviseur / Appli « Alerte Lycée »};
\node[decision, below of=signal, align=center] (anonymat){Anonymat\\garanti ?};
\node[process, below of=anonymat, align=center] (enquete){ENQUÊTE DISCRÈTE\\CPE / Vie scolaire /\\Proviseur adjoint};
\node[decision, below of=enquete, align=center] (preuve){Preuves\\suffisantes ?};
\node[process, left of=preuve, xshift=-3.6cm, align=center] (conseil){CONSEIL DE DISCIPLINE\\Sanction éducative /\\Exclusion};
\node[process, right of=preuve, xshift=3.8cm, align=center] (conac){TRANSMISSION EXTERNE\\CONAC / Police / Justice\\(détournement, violences graves)};
\node[process, below of=preuve, yshift=-1.6cm, align=center] (sanction){SANCTION / RÉPARATION\\Exclusion temporaire, travaux\\d'intérêt général, médiation};
\node[process, below of=sanction, align=center] (protection){PROTECTION DU LANCEUR D'ALERTE\\Suivi, anonymat maintenu,\\non-représailles};
\node[startstop, below of=protection, align=center] (end){FIN\\Clôture du dossier /\\Retour au signalant};

% Points de vigilance
\node[vigilance, right of=signal, xshift=4.2cm, align=center] (vig1){POINT VIGILANCE 1\\Anonymat physique / numérique ?\\(boîte scellée, formulaire sans e-mail)};
\node[vigilance, right of=enquete, xshift=4.2cm, align=center] (vig2){POINT VIGILANCE 2\\Délai maximal de 48 h pour\\l'accusé de réception ?};
\node[vigilance, right of=sanction, xshift=4.2cm, align=center] (vig3){POINT VIGILANCE 3\\Retour au signalant ?\\(information sur la suite donnée)};

% Flèches principales
\draw[arrow] (start) -- (signal);
\draw[arrow] (signal) -- (anonymat);
\draw[arrow] (anonymat) -- node[left, font=\footnotesize] {OUI} (enquete);
\draw[arrow] (anonymat.east) -- ++(1,0) |- (vig1.south west) node[pos=0.3, above, font=\footnotesize] {NON $\to$ corriger};
\draw[arrow] (enquete) -- (preuve);
\draw[arrow] (preuve) -- node[left, font=\footnotesize] {OUI} (sanction);
\draw[arrow] (preuve) -- node[above, font=\footnotesize] {Grave} (conac);
\draw[arrow] (preuve.west) -- node[above, font=\footnotesize] {OUI} (conseil);
\draw[arrow] (conseil) |- (sanction);
\draw[arrow] (conac) |- (sanction);
\draw[arrow] (sanction) -- (protection);
\draw[arrow] (protection) -- (end);

% Flèches vigilance (pointillés)
\draw[arrow, dashed, poppink] (signal.east) -- (vig1.west);
\draw[arrow, dashed, poppink] (enquete.east) -- (vig2.west);
\draw[arrow, dashed, poppink] (sanction.east) -- (vig3.west);

% Boucle retour
\draw[arrow, popgreen, bend right=45] (end.west) to node[left, font=\footnotesize, text=popgreen] {Bilan / statistiques annuelles} (start.west);
\end{tikzpicture}
\legende{Organigramme « Parcours d'un signalement d'atteinte à l'intégrité au lycée ». Les encadrés roses sont des points de contrôle qualité de la procédure. L'absence de retour au signalant (vigilance 3) détruit la confiance et tarit les signalements.}
\end{popfigure}

\begin{exoresolu}[Analyse critique d'une procédure]
\textbf{Énoncé :} un élève dépose un mot anonyme dans la boîte à idées : « M. X (enseignant) demande 5 000 FCFA pour valider le TP ». La boîte est relevée par le délégué de classe, qui remet le mot au CPE. Le CPE convoque l'élève (anonymat levé de fait) et l'enseignant. L'enseignant nie. Il n'y a pas de preuve écrite. Le CPE classe sans suite.

\tcblower

\textbf{Analyse à l'aide de l'organigramme :}
\begin{enumerate}
    \item \textbf{Signalement :} correct (boîte à idées $\to$ délégué $\to$ CPE).
    \item \textbf{Vigilance 1 (anonymat) :} \textbf{échec.} Le délégué a lu le mot ; le CPE a convoqué l'élève sur la base du mot, levant l'anonymat de fait. Risque de représailles contre l'élève.
    \item \textbf{Enquête :} faible (confrontation immédiate sans recherche de preuves).
    \item \textbf{Décision « preuves suffisantes ? » :} NON (parole contre parole).
    \item \textbf{Issue :} classement sans suite = impunité perçue.
    \item \textbf{Vigilance 3 (retour) :} aucun. L'élève ne sait pas ce qu'il est advenu de son signalement.
\end{enumerate}
\textbf{Procédure corrigée :} la boîte est relevée \textbf{uniquement} par le CPE ou le Proviseur (clé unique) ; enquête discrète (vérifications, recherche d'autres témoignages) avant toute confrontation ; réponse écrite au signalant (via la boîte ou une application anonyme) : « signalement reçu, enquête en cours / clôturée faute de preuve / transmise à la hiérarchie ».
\end{exoresolu}

\courssub{TP3 -- Rédaction collaborative : la « Charte de l'intégrité de la 1re Tech »}

\begin{attention}[Pièges à éviter dans la rédaction d'une charte]
\begin{itemize}
    \item \textbf{Vœux pieux :} « Je serai honnête », « Je respecterai tout le monde » : inapplicable, non vérifiable, subjectif.
    \item \textbf{Négations floues :} « Je ne tricherai pas » : qu'est-ce que tricher ? Copier ? Avoir son téléphone ? Partager un fichier ?
    \item \textbf{Sanctions non définies :} « Les contrevenants seront punis » : par qui ? comment ? selon quelle échelle ?
    \item \textbf{Oubli des « zones grises » :} prêt de cahier pendant un contrôle, « coup de main » sur une machine en TP, gestion de la caisse de classe par le délégué.
\end{itemize}
\textbf{Règle d'or :} un article = un engagement \textbf{observable, mesurable, contextualisé}.
\end{attention}

\begin{methode}[Élaborer la charte de classe (10 articles maximum)]
\begin{enumerate}
    \item \textbf{Remue-méninges (15 min) :} par groupes de 4, lister 5 situations concrètes vécues ou redoutées dans la spécialité (électrotechnique, mécanique, génie civil, commerce, etc.).
    \item \textbf{Mise en commun et fusion (15 min) :} au tableau, regrouper par thème (évaluation, matériel/atelier, vie de classe/argent, relations/autorité). Voter les 10 situations prioritaires.
    \item \textbf{Rédaction formelle (20 min) :} chaque groupe rédige 2 à 3 articles. Vérifier : sujet (je/nous), verbe d'action, contexte, critère de réussite observable.
    \item \textbf{Relecture croisée (10 min) :} échange des feuilles. Test : « Puis-je \emph{prouver} que je l'ai respecté aujourd'hui ? »
    \item \textbf{Signature solennelle (5 min) :} impression A3, lecture collective, signature de chaque élève, du professeur principal et des délégués. Affichage en salle de classe et copie dans le cahier de textes.
\end{enumerate}
\end{methode}

\begin{exemplebox}[Modèle d'articles « spécialisés » (exemple : 1re Électrotechnique)]
\begin{itemize}
    \item \textbf{Art. 1 (Évaluation) :} « Je compose mes devoirs et TP sans communication ni document non autorisé ; je signale tout échange suspect au surveillant. »
    \item \textbf{Art. 2 (Atelier / matériel) :} « Je rends le matériel complet et en bon état ; je ne prête pas mon badge d'accès à l'atelier à un camarade non habilité. »
    \item \textbf{Art. 3 (Conflit d'intérêts -- délégué trésorier) :} « Si je suis délégué trésorier, je présente le budget mensuel (recettes/dépenses) à la classe avant le 5 du mois suivant ; je refuse tout achat sans facture. »
    \item \textbf{Art. 4 (Stage / entreprise) :} « Je refuse tout pourboire ou cadeau de l'entreprise d'accueil pour une tâche relevant de ma formation ; je signale toute demande anormale au tuteur de l'école. »
    \item \textbf{Art. 5 (Signalement) :} « Je dépose tout signalement (corruption, harcèlement, vol) dans la boîte "Intégrité" (verrouillée, clé détenue par le CPE) ou via l'application "Alerte Lycée" ; je sais que mon anonymat est protégé. »
    \item \textbf{Art. 6 (Réparation) :} « En cas de manquement avéré, j'accepte la médiation par les pairs (délégués + CPE) et la sanction éducative (réparation, excuses publiques, travaux d'intérêt général) avant tout conseil de discipline. »
    \item \textbf{Art. 7 (Numérique) :} « Je ne partage pas mes codes d'accès ; je ne modifie pas les notes ou fichiers d'autrui ; je cite mes sources dans mes rapports. »
    \item \textbf{Art. 8 (Climat scolaire) :} « Je m'oppose au racket et au harcèlement dans les toilettes et recoins : j'interviens si je ne me mets pas en danger, sinon j'alerte immédiatement un adulte. »
    \item \textbf{Art. 9 (Exemplarité) :} « Je rappelle la charte à tout nouvel arrivant ; je la fais respecter lors des travaux de groupe (pas de "passager clandestin"). »
    \item \textbf{Art. 10 (Révision) :} « Cette charte est révisée chaque semestre sur proposition de 5 élèves ou du professeur principal ; toute version validée fait l'objet d'une nouvelle signature. »
\end{itemize}
\end{exemplebox}

\begin{savaistu}[Transparence budgétaire : de l'État à la classe]
Au Cameroun, l'article 66 de la Constitution du 18 février 1996 prévoit la \textbf{déclaration des biens et avoirs} par les hauts responsables publics, et la CONAC (Commission nationale anti-corruption, créée par le décret n° 2006/088 du 11 mars 2006) lutte contre la corruption dans l'administration. Transposé en classe : le \textbf{délégué trésorier} (élu) présente un \textbf{rapport mensuel} (recettes : cotisations ; dépenses : fournitures, activités) avec pièces justificatives (factures, tickets), que la classe valide. C'est l'apprentissage concret de la \cle{reddition des comptes} et de la \cle{transparence}, piliers de l'intégrité publique.
\end{savaistu}

\courssub{TP4 -- Analyse de données : sanctions disciplinaires 2021-2023}

\begin{experience}[Exploitation de données disciplinaires (tableur $\to$ graphique $\to$ interprétation)]
\textbf{Données fournies (extrait simulé du registre de vie scolaire) :}
\begin{center}
\begin{tabularx}{\linewidth}{|l|X|X|X|X|X|}\hline
\rowcolor{popblueL}\textbf{Année} & \textbf{Tricherie} & \textbf{Vol} & \textbf{Violence} & \textbf{Absentéisme/Autre} & \textbf{Total} \\ \hline
2021 & 45 & 20 & 15 & 20 & 100 \\ \hline
2022 & 38 & 22 & 18 & 22 & 100 \\ \hline
2023 & 30 & 25 & 25 & 20 & 100 \\ \hline
\textbf{Total} & \textbf{113} & \textbf{67} & \textbf{58} & \textbf{62} & \textbf{300} \\ \hline
\end{tabularx}
\end{center}
\textit{Note : effectifs constants, environ 1 200 élèves ; les pourcentages sont calculés sur le total des sanctions de chaque année (100).}

\textbf{Consignes :}
\begin{enumerate}
    \item Saisir le tableau dans LibreOffice Calc ou Excel. Vérifier le type des cellules (nombre).
    \item Calculer les \textbf{parts en \%} par type et par année (ex. : tricherie 2021 $= 45/100 = 45\,\%$).
    \item Construire \textbf{deux graphiques} :
        \begin{itemize}
            \item un \textbf{diagramme en barres groupées} : années en abscisse, nombre de sanctions en ordonnée, 4 séries colorées $\to$ visualise l'évolution par type ;
            \item un \textbf{diagramme circulaire} pour 2023 : part de chaque faute dans le total $\to$ visualise la structure actuelle.
        \end{itemize}
    \item Rédiger l'interprétation (15 lignes maximum) : la tricherie est-elle la faute la plus sanctionnée ? Quelles hypothèses sur les causes des évolutions ?
\end{enumerate}
\end{experience}

\begin{popfigure}
\begin{tikzpicture}
\begin{axis}[
    ybar,
    width=\linewidth,
    height=8.5cm,
    ylabel={Nombre de sanctions},
    xlabel={Année},
    symbolic x coords={2021,2022,2023},
    xtick=data,
    legend style={at={(0.5,-0.18)}, anchor=north, legend columns=-1, font=\footnotesize, draw=popline},
    ymin=0, ymax=55,
    bar width=0.55cm,
    enlarge x limits=0.25,
    nodes near coords, every node near coord/.style={font=\tiny},
]
\addplot[fill=popblue, draw=popdark] coordinates {(2021,45) (2022,38) (2023,30)}; \addlegendentry{Tricherie}
\addplot[fill=poporange, draw=popdark] coordinates {(2021,20) (2022,22) (2023,25)}; \addlegendentry{Vol}
\addplot[fill=popgreen, draw=popdark] coordinates {(2021,15) (2022,18) (2023,25)}; \addlegendentry{Violence}
\addplot[fill=poppink, draw=popdark] coordinates {(2021,20) (2022,22) (2023,20)}; \addlegendentry{Absentéisme/Autre}
\end{axis}
\end{tikzpicture}
\legende{Évolution des sanctions par type de faute (2021-2023, données simulées). Baisse nette de la tricherie (de 45 à 30, soit $-33\,\%$ en trois ans). Hausse continue du vol (de 20 à 25) et surtout de la violence (de 15 à 25, soit $+67\,\%$). L'absentéisme reste stable autour de 20.}
\end{popfigure}

\begin{popfigure}
\begin{tikzpicture}
    % Diagramme circulaire 2023 : Tricherie 30% (108°), Vol 25% (90°), Violence 25% (90°), Absentéisme 20% (72°)
    \fill[popblue]   (0,0) -- (90:3) arc (90:198:3) -- cycle;
    \fill[poporange] (0,0) -- (198:3) arc (198:288:3) -- cycle;
    \fill[popgreen]  (0,0) -- (288:3) arc (288:378:3) -- cycle;
    \fill[poppink]   (0,0) -- (378:3) arc (378:450:3) -- cycle;
    \draw[popdark, thick] (0,0) circle (3);
    % Étiquettes
    \node[font=\footnotesize\bfseries, text=white] at (144:1.8) {30\,\%};
    \node[font=\footnotesize\bfseries, text=white] at (243:1.8) {25\,\%};
    \node[font=\footnotesize\bfseries, text=white] at (333:1.8) {25\,\%};
    \node[font=\footnotesize\bfseries, text=white] at (54:1.8) {20\,\%};
    % Légende
    \node[anchor=west, font=\footnotesize] at (3.6,1.5) {\tikz\fill[popblue] (0,0) rectangle (0.3,0.3); \ Tricherie (30\,\%)};
    \node[anchor=west, font=\footnotesize] at (3.6,0.8) {\tikz\fill[poporange] (0,0) rectangle (0.3,0.3); \ Vol (25\,\%)};
    \node[anchor=west, font=\footnotesize] at (3.6,0.1) {\tikz\fill[popgreen] (0,0) rectangle (0.3,0.3); \ Violence (25\,\%)};
    \node[anchor=west, font=\footnotesize] at (3.6,-0.6) {\tikz\fill[poppink] (0,0) rectangle (0.3,0.3); \ Absentéisme/Autre (20\,\%)};
\end{tikzpicture}
\legende{Répartition des sanctions en 2023 (total : 100). La tricherie reste la première cause unique (30\,\%), mais la somme vol + violence (50\,\%) la dépasse : le climat scolaire se dégrade malgré la baisse de la fraude académique.}
\end{popfigure}

\begin{exoresolu}[Interprétation type attendue]
\textbf{Énoncé :} à partir des deux graphiques, répondre aux questions : la tricherie est-elle la faute la plus sanctionnée ? Comment expliquer les évolutions observées ? Quelles conséquences pour nos outils ?

\tcblower

\textbf{1. La tricherie est-elle la plus sanctionnée ?} \textbf{Oui}, en valeur absolue et relative chaque année (de 45\,\% à 30\,\% du total annuel). C'est la faute la plus traitée par l'institution.

\textbf{2. Est-elle la plus fréquente en réalité ?} Probablement, car le taux de détection est élevé (surveillance des examens, vérification des copies). Mais il existe un « chiffre noir » : les fraudes non détectées n'apparaissent pas dans les statistiques.

\textbf{3. Analyse dynamique :}
\begin{itemize}
    \item \textbf{Baisse de la tricherie} (45 $\to$ 30, soit $-33\,\%$) : corrélation avec l'installation de caméras dans les salles d'examen (2022) et le renforcement de la surveillance (2023). Les mesures techniques de contrôle fonctionnent sur la fraude académique.
    \item \textbf{Hausse de la violence} (15 $\to$ 25, soit $+67\,\%$) \textbf{et du vol} (20 $\to$ 25) : la répression technique déplace peut-être la déviance vers des actes relationnels (racket, bagarres, vols de téléphones) moins détectables par caméra ; elle peut aussi révéler un mal-être social non traité (tensions entre groupes, manque d'encadrement de la vie scolaire).
    \item \textbf{Stabilité de l'absentéisme} : facteur structurel (éloignement, travail, démotivation), peu sensible aux mesures de sécurité.
\end{itemize}

\textbf{Conclusion pour l'action :} la carte de vigilance (TP1) doit intégrer ces nouvelles zones de violence (cour de récréation, abords du portail) ; la charte (TP3) doit inclure des articles sur la non-violence et le respect du bien d'autrui ; la procédure (TP2) doit traiter le signalement des violences avec la même rigueur que celui de la corruption.
\end{exoresolu}

\courssub{Restitution orale, débat et engagement}

\begin{methode}[Restitution orale structurée (5 min par groupe)]
\begin{enumerate}
    \item \textbf{Présentation de l'outil} (1 min) : nom, format (affiche A3, fichier PDF, tableau de bord), auteurs.
    \item \textbf{Démarche} (1 min) : comment les informations ont-elles été collectées ? (enquête, données administratives, remue-méninges).
    \item \textbf{Résultat clé} (1 min) : la zone rouge n°1, l'article 3 de la charte, la courbe de la violence en 2023.
    \item \textbf{Applicabilité et protection} (2 min) : comment cela fonctionne-t-il demain matin ? Qui gère la boîte ? Qui valide la charte ? Comment l'anonymat est-il \emph{techniquement} garanti (boîte scellée, formulaire sans e-mail, tiers de confiance) ? Quelle sanction si un adulte viole l'anonymat ?
\end{enumerate}
\end{methode}

\begin{aretenir}[Débat : « Ces outils sont-ils applicables ? Comment garantir l'anonymat et la protection ? »]
\textbf{Arguments en faveur :} outils concrets, issus du terrain, appropriés par les élèves (légitimité), faible coût (papier, logiciels libres), création d'une culture de la preuve et de la procédure.

\textbf{Risques identifiés :} boîte à idées vidée par une personne non autorisée ? délégué corrompu ? charte oubliée après un mois ? signalement calomnieux (dénonciation mensongère) ?

\textbf{Garanties techniques et institutionnelles :}
\begin{enumerate}
    \item \textbf{Tiers de confiance :} la boîte physique possède \textbf{deux clés} (Proviseur + représentant des parents d'élèves). Le formulaire numérique ne collecte \textbf{aucune adresse e-mail} (vérifié par le professeur d'ECM).
    \item \textbf{Traçabilité du traitement :} registre numéroté des signalements (date, nature, suite donnée), accessible aux délégués élus (contrôle simulé).
    \item \textbf{Protection du lanceur d'alerte :} clause explicite dans le règlement intérieur : « Toute représaille contre un signalant de bonne foi est une faute grave ».
    \item \textbf{Rituel :} « Point intégrité » mensuel de 10 minutes en vie de classe (bilan des signalements, avancement de la charte, mise à jour de la carte).
\end{enumerate}
\end{aretenir}

\begin{exercice}[Mise en situation : « Le cas du délégué trésorier »]
Le délégué trésorier de la 1re Tech Génie Civil gère la caisse de classe (cotisation de 1 000 FCFA par mois). En décembre, il annonce : « Il manque 15 000 FCFA, je ne sais pas où c'est passé. » Il n'a tenu aucun cahier de caisse et n'a gardé aucune facture pour les achats de peinture de l'atelier.
\begin{enumerate}
    \item Quel article de votre charte (TP3) s'applique ? Est-il assez précis ?
    \item Sur la carte (TP1), où se situe le risque ? (symbole, couleur)
    \item Dans la procédure (TP2), qui saisit qui ? En quel délai ?
    \item Proposez une \textbf{mesure de réparation éducative} (et pas seulement punitive) pour le délégué et pour la classe.
\end{enumerate}
\end{exercice}

\begin{corrige}[Exercice : Le cas du délégué trésorier]
\begin{enumerate}
    \item \textbf{Charte :} article 3 (conflit d'intérêts / transparence budgétaire) et article 6 (réparation). Si l'article dit seulement « je gère bien la caisse », il est insuffisant. Il doit imposer : « cahier de caisse tenu à jour, présenté mensuellement, pièces justificatives obligatoires pour toute dépense supérieure à 2 000 FCFA ».
    \item \textbf{Carte :} salle de classe (caisse conservée en classe) $\to$ symbole \textbf{F} (risque financier) en \textbf{rouge}, car le risque est avéré et le contrôle absent.
    \item \textbf{Procédure :} un élève cotisant (témoin) $\to$ boîte à idées ou CPE $\to$ enquête (vérification des comptes, audition du délégué) $\to$ médiation de vie scolaire $\to$ remboursement et/ou travaux d'intérêt général $\to$ élection d'un nouveau trésorier. Délai : une semaine maximum.
    \item \textbf{Réparation éducative :} le délégué démissionne ou est révoqué. Il établit un \textbf{relevé de comptes rétroactif} (recherche des tickets, témoignages) aidé par le professeur de gestion-comptabilité, puis le présente à la classe. Il effectue 5 heures de travaux d'intérêt général (tenue du nouveau cahier de caisse pour son successeur). La classe vote éventuellement une cotisation exceptionnelle de solidarité (500 FCFA) pour combler le déficit, validant ainsi la responsabilité collective.
\end{enumerate}
\end{corrige}

\coursec{Intégrité et avenir professionnel technique : éthique de l'ingénieur et du technicien}

Après avoir construit collectivement des outils de promotion de l'intégrité au sein même de l'établissement (chartes, affiches, croquis d'espaces de vie scolaire), il est temps de projeter cette exigence morale vers votre avenir professionnel immédiat. En tant qu'élèves de l'enseignement technique — futurs techniciens supérieurs, ingénieurs, géomètres, informaticiens, comptables ou chefs de chantier — vous ne serez pas de simples exécutants. Vous serez des \textbf{garants de la sécurité, de la fiabilité et de la conformité} d'ouvrages, de systèmes ou de données dont dépendent des vies humaines et des intérêts économiques majeurs. Cette section ancre l'intégrité morale dans la réalité technique et juridique de vos futures professions.

\courssub{La spécificité de l'intégrité technique : une responsabilité vitale}

L'intégrité morale, dans un contexte technique, dépasse largement l'honnêteté intellectuelle ou le respect des règles de vie scolaire. Elle engage la \textbf{responsabilité vitale} : l'erreur, la négligence ou la fraude technique ne sont pas de simples fautes de gestion, elles sont potentiellement mortelles.

\begin{definition}[Intégrité technique]
L'intégrité technique est l'obligation morale, déontologique et légale pour un professionnel technique d'agir avec compétence, honnêteté, rigueur et impartialité dans la conception, la réalisation, le contrôle et la maintenance d'ouvrages, de systèmes ou de services, en privilégiant impérativement la sécurité des personnes, la protection de l'environnement et l'intérêt général sur tout intérêt particulier, financier ou hiérarchique.
\end{definition}

\begin{propriete}[Portée juridique de la négligence technique]
Au Cameroun, la négligence caractérisée dans l'exercice d'une profession technique est pénalement sanctionnée. Le Code pénal réprime la \textbf{mise en danger de la vie d'autrui} ainsi que les \textbf{blessures et homicides involontaires} résultant d'une maladresse, d'une imprudence, d'une inattention, d'une négligence ou de l'inobservation des règlements (art. 275 et s.). La loi n°94/018 portant organisation de la profession d'ingénieur et la loi n°2011/022 portant organisation de la profession d'architecte prévoient des sanctions disciplinaires (jusqu'à la radiation) et pénales pour manquements graves. Les normes ISO 9001 (management de la qualité), ISO 14001 (management environnemental) et ISO 45001 (santé et sécurité au travail) imposent une traçabilité des décisions techniques qui engage la responsabilité personnelle du signataire.
\end{propriete}

\emph{Exemples concrets par filière :}
\begin{itemize}
    \item \textbf{BTP / Génie civil :} utilisation de fer à béton de diamètre inférieur à celui prescrit (par exemple HA10 au lieu de HA16), ciment frelaté (taux de clinker insuffisant), sable non lavé (contenant des sels ou de l'argile), dosage eau/ciment non respecté. Conséquences : effondrement de dalles, fissuration structurelle, ruine de l'ouvrage, morts.
    \item \textbf{Électricité / Électrotechnique :} câblage sous-dimensionné, absence de protection différentielle, mise à la terre défaillante, utilisation de matériel non certifié (contrefaçon). Conséquences : électrocution, incendie d'origine électrique, destruction d'équipements.
    \item \textbf{Informatique / Réseaux :} laisser une faille de sécurité connue non corrigée, stocker des mots de passe en clair, vendre des données clients, installer des logiciels espions. Conséquences : vol d'identité, paralysie d'un hôpital ou d'une administration, ruine financière, atteinte à la souveraineté numérique.
    \item \textbf{Industrie / Maintenance :} remise en service d'une machine sans vérification des dispositifs de sécurité (arrêt d'urgence, capteurs), falsification de carnets de maintenance. Conséquences : accident du travail grave, pollution industrielle.
\end{itemize}

\begin{attention}[Le piège du « petit arrangement » technique]
Valider un câblage non conforme « pour dépanner le client », signer un procès-verbal de réception sans avoir vérifié le béton frais, accepter un planning irréaliste en sachant que les essais seront bâclés : ce ne sont pas des « services rendus » ni de la « flexibilité ». Ce sont des \textbf{fautes graves engageant votre responsabilité pénale}. L'écrit (votre signature sur un PV, un rapport, un bon de commande, un schéma validé) constitue une preuve devant les tribunaux. « J'ai obéi à mon chef » n'est \emph{jamais} une défense suffisante en matière de sécurité des personnes.
\end{attention}

\courssub{Le cadre déontologique et institutionnel camerounais}

L'exercice des professions techniques réglementées au Cameroun est encadré par des Ordres professionnels qui détiennent le pouvoir disciplinaire et veillent au respect de la déontologie.

\begin{center}
\begin{tabularx}{\linewidth}{|l|X|X|}\hline
\rowcolor{popblueL}\textbf{Profession / Ordre} & \textbf{Textes fondateurs (repères)} & \textbf{Engagements clés du code de déontologie} \\ \hline
\textbf{Ingénieurs (ONIGC)} & Loi n°94/018 du 21 décembre 1994 portant organisation de la profession d'ingénieur & Exercer avec conscience, probité et honneur ; ne rien faire qui puisse porter atteinte à la sécurité, à la santé et aux intérêts légitimes des usagers. \\ \hline
\textbf{Architectes (ONAC)} & Loi n°2011/022 du 14 décembre 2011 & Respect des règles de l'art, des normes d'urbanisme et de l'environnement ; indépendance vis-à-vis des entreprises ; secret professionnel. \\ \hline
\textbf{Géomètres (ONIGC, section géomètres)} & Textes organisant la profession de géomètre & Exactitude des bornages, impartialité dans les expertises foncières, respect du secret professionnel. \\ \hline
\textbf{Experts-comptables (ONECCA)} & Loi n°92/006 du 14 août 1992 ; Actes uniformes OHADA & Indépendance, compétence, secret professionnel, lutte contre le blanchiment, régularité et sincérité des comptes. \\ \hline
\textbf{Techniciens supérieurs / Agents de maîtrise} & Code du travail (Loi n°92/007) ; règlements intérieurs d'entreprise & Obéissance aux consignes de sécurité, signalement des dangers (droit d'alerte), loyauté envers l'employeur, respect de la propriété intellectuelle. \\ \hline
\end{tabularx}
\end{center}

\begin{savaistu}[La radiation par l'Ordre : la fin d'une carrière]
L'Ordre National des Ingénieurs de Génie Civil du Cameroun (ONIGC) peut prononcer la \textbf{radiation} du tableau de l'Ordre pour manquement grave à l'honneur, à la probité ou à la dignité professionnelle. La radiation entraîne l'\textbf{interdiction d'exercer la profession} sur tout le territoire national : plus de signature de plans, plus de direction de travaux, plus de consultation. C'est la « peine de mort » professionnelle. La même logique disciplinaire s'applique aux autres Ordres (architectes, géomètres, experts-comptables).
\end{savaistu}

\courssub{Conflits d'intérêts types en milieu technique : identifier le piège}

Le conflit d'intérêts est la situation la plus fréquente où l'intégrité technique bascule. Il ne s'agit pas seulement de corruption avérée, mais de toute situation où le jugement technique est — ou paraît — altéré.

\begin{definition}[Conflit d'intérêts en milieu technique]
Situation où l'intérêt personnel (financier, familial, amical, de carrière) d'un professionnel technique interfère ou \emph{paraît} interférer avec son obligation professionnelle d'agir exclusivement dans l'intérêt de son client, de son employeur ou du public, conformément aux règles de l'art et aux normes en vigueur.
\end{definition}

\begin{exemplebox}[Conflits d'intérêts « classiques » en entreprise et sur chantier]
\begin{enumerate}
    \item \textbf{Choix des fournisseurs / sous-traitants (rétro-commissions) :} le chef de projet oriente le marché vers l'entreprise d'un proche ou celle qui lui verse une commission occulte, au détriment du meilleur rapport qualité/prix/délai. Résultat : matériaux bas de gamme livrés, factures majorées.
    \item \textbf{Certification de complaisance :} un laboratoire de contrôle valide des résultats d'essais (compression du béton, granulométrie du sable) non conformes aux normes, sous pression du chef de chantier ou du fournisseur. « On ferme les yeux, ça passera. »
    \item \textbf{Faux rapports d'avancement / PV de réception :} signature de procès-verbaux de réception d'ouvrages (électricité, plomberie, VRD) sans visite sur site, ou en masquant des réserves majeures pour déclencher le paiement.
    \item \textbf{Vol de propriété intellectuelle :} copie de plans d'exécution, de codes source, de schémas de câblage ou de méthodes de calcul développés par l'employeur ou un concurrent, pour créer sa propre structure ou les revendre.
    \item \textbf{Double emploi / conseil concurrent :} un ingénieur salarié d'un bureau d'études contrôle les travaux d'une entreprise pour laquelle il fait secrètement du conseil technique le week-end.
\end{enumerate}
\end{exemplebox}

\begin{methode}[Gestion d'un conflit d'intérêts : la procédure en 4 étapes]
\begin{enumerate}
    \item \textbf{Identifier (auto-évaluation) :} suis-je personnellement concerné ? Un proche (conjoint, enfant, parent, ami intime) l'est-il ? Ai-je un intérêt financier direct ou indirect (parts, commissions, promesse d'embauche) dans la décision à prendre ?
    \item \textbf{Déclarer (formalisation écrite) :} rédiger une \textbf{déclaration de conflit d'intérêts} adressée à la hiérarchie directe et/ou au service conformité/éthique de l'entreprise. Préciser la nature du lien, l'opération concernée, le risque identifié. Conserver une copie horodatée (mail avec accusé de réception, lettre recommandée, registre interne).
    \item \textbf{S'abstenir (récusation effective) :} ne \emph{pas} participer à la délibération, au vote, à la négociation, à la signature ou au contrôle technique de l'opération concernée. Se faire remplacer par un collègue compétent et non concerné.
    \item \textbf{Assurer la traçabilité (archivage) :} conserver la déclaration, la décision de remplacement et les pièces du dossier. Cette traçabilité est votre \textbf{protection juridique} en cas de contentieux ultérieur.
\end{enumerate}
\end{methode}

\courssub{Étude de cas technique : l'affaire de l'immeuble « Les Jardins du Lac » (Yaoundé)}

\begin{exoresolu}[Analyse d'un effondrement structurel imputable à une chaîne de manquements à l'intégrité]
\textbf{Énoncé (cas fictif d'école) :} Un immeuble R+4 en construction à Yaoundé s'effondre partiellement lors du coulage de la dalle du 3\textsuperscript{e} niveau. Bilan : 3 ouvriers tués, 7 blessés graves. L'enquête judiciaire et l'expertise technique révèlent :
\begin{itemize}
    \item \textbf{Ferraillage :} diamètre HA10 utilisé au lieu du HA16 prévu au plan d'exécution (économie d'environ 15 \% sur l'acier).
    \item \textbf{Béton :} ciment de classe inférieure à celle prescrite ; dosage ciment de 280 kg/m\textsuperscript{3} au lieu de 350 kg/m\textsuperscript{3} ; ajout d'eau en excessif pour « fluidifier » (rapport eau/ciment supérieur à 0,65).
    \item \textbf{Granulats :} sable de carrière non lavé (teneur en fines supérieure à 8 \%, présence de chlorures), gravier mélangé à de la latérite.
    \item \textbf{Contrôle :} le laboratoire privé missionné pour le contrôle qualité a délivré des PV de conformité sur la base d'éprouvettes « témoins » fabriquées avec un béton de référence, et non prélevées sur le béton réellement coulé. Le chef de chantier a « validé » les livraisons sans vérifier les bons et certificats.
    \item \textbf{Flux financiers :} virements suspects du fournisseur vers les comptes personnels du chef de chantier et du responsable du laboratoire (découverts par la police judiciaire).
\end{itemize}

\textbf{Analyse par le modèle de décision éthique technique en 5 étapes :}

\begin{popfigure}
\centering
\begin{tikzpicture}[
    node distance=1.2cm and 2.2cm,
    box/.style={draw=popblue, thick, rounded corners, fill=popblue!5, minimum width=3.4cm, minimum height=1.6cm, align=center, font=\small},
    arrow/.style={-{Stealth}, thick, popdark}
]
\node[box, align=center] (facts){\textbf{1. FAITS TECHNIQUES}\\\emph{Qu'observe-t-on ?}\\Fer HA10 au lieu de HA16\\Ciment sous-classé\\Sable non lavé\\PV de labo falsifiés};
\node[box, right=of facts, align=center] (rules){\textbf{2. RÈGLES / NORMES}\\\emph{Qu'exige-t-on ?}\\Règles BAEL / Eurocode 2\\Norme béton (NF EN 206)\\Code du travail (sécurité)\\Loi 94/018 (déontologie)};
\node[box, below=of facts, align=center] (conseq){\textbf{3. CONSÉQUENCES}\\\emph{Si l'on ne fait rien ?}\\\textbf{Sécurité :} morts, blessés\\\textbf{Juridique :} poursuites pénales\\\textbf{Financier :} ruine, indemnisations\\\textbf{Réputation :} radiation, exclusion};
\node[box, right=of conseq, align=center] (alts){\textbf{4. ALTERNATIVES}\\\emph{Que faire ?}\\\textbf{STOP chantier immédiat}\\Alerter hiérarchie / maître d'ouvrage\\Saisir l'Ordre professionnel\\Signalement au Procureur\\Refus de signer le PV};
\node[box, below=of conseq, xshift=3.4cm, align=center] (decision){\textbf{5. DÉCISION / ACTION}\\\emph{Choix éthique et légal}\\\textbf{Résistance active}\\Démission si ordre maintenu\\Témoignage devant la justice\\Reconstruction conforme};

\draw[arrow] (facts) -- (rules);
\draw[arrow] (facts) -- (conseq);
\draw[arrow] (rules) -- (alts);
\draw[arrow] (conseq) -- (alts);
\draw[arrow] (alts) -- (decision);
\draw[arrow, poppink, dashed] (decision) to[out=180, in=-90, looseness=1.4] node[left, font=\tiny, poppink, align=center] {Retour\\d'expérience} (facts);
\end{tikzpicture}
\legende{Modèle de décision éthique technique en 5 étapes appliqué à l'affaire « Les Jardins du Lac ». La boucle en pointillés symbolise le retour d'expérience et l'amélioration continue du système qualité-intégrité.}
\end{popfigure}

\textbf{Application au cas :}
\begin{enumerate}
    \item \textbf{Faits :} non-conformités majeures des matériaux et de l'exécution, falsification des PV de contrôle.
    \item \textbf{Règles :} règles de calcul des structures (BAEL / Eurocode 2), norme béton NF EN 206, Code pénal (mise en danger d'autrui), loi 94/018 (déontologie de l'ingénieur).
    \item \textbf{Conséquences :} \textbf{réalisées} (3 morts, 7 blessés graves). \textbf{Encourues} par les signataires (chef de chantier, ingénieur structure, responsable laboratoire) : lourdes peines de prison pour homicide involontaire et corruption, radiation des Ordres, exclusion des marchés publics, ruine assurantielle.
    \item \textbf{Alternatives :} arrêt immédiat (droit d'alerte et droit de retrait), alerte écrite au maître d'ouvrage et à l'Ordre, refus de signer tout document validant l'existant.
    \item \textbf{Décision :} le technicien ou l'ingénieur intègre choisit l'alternative « Arrêt + Alerte + Refus ». Si la direction impose la poursuite des travaux : démission immédiate et signalement au Procureur de la République, car tout citoyen ayant connaissance d'un crime ou d'un délit peut et doit le dénoncer.
\end{enumerate}

\textbf{Enseignement :} la chaîne de complicité (fournisseur $\rightarrow$ chef de chantier $\rightarrow$ laboratoire $\rightarrow$ hiérarchie) est brisée par \textbf{un seul acte d'intégrité} : le refus de valider le non-conforme. Le coût du « non » (pression, menace de licenciement) est infiniment inférieur au coût du « oui » (prison, morts, fin de carrière).
\end{exoresolu}

\courssub{Construire son « CV Intégrité » : la preuve par le parcours}

Face à un recruteur (entreprise, bureau d'études, administration, ONG) ou pour une inscription à un Ordre professionnel, le CV classique (diplômes, stages, logiciels) ne suffit plus. La \textbf{fiabilité éthique} devient un critère de sélection majeur.

\begin{center}
\begin{tabularx}{\linewidth}{|l|X|X|}\hline
\rowcolor{popblueL}\textbf{Rubrique CV classique} & \textbf{Rubrique « CV Intégrité »} & \textbf{Preuves tangibles à archiver} \\ \hline
Stages effectués & \textbf{Stages validés sans réserve éthique} & Attestations mentionnant « respect des consignes de sécurité », « rigueur dans les rapports », « signalement d'anomalie » ; lettre de recommandation du tuteur citant un acte d'intégrité. \\ \hline
Projets académiques & \textbf{Projets menés avec traçabilité complète} & Cahier des charges respecté, PV de réception signés, gestion des non-conformités documentée (ex. : « pièce non conforme détectée, rebutée, remplacée, tracée »). \\ \hline
Compétences techniques & \textbf{Compétences \emph{normatives} maîtrisées} & Liste des normes maîtrisées (NF C 15-100, Eurocode 2, ISO 9001, référentiel comptable OHADA) + attestations de formation continue (normes, sécurité, éthique). \\ \hline
Références & \textbf{Références morales et professionnelles} & Contacts de 2 ou 3 personnes (maîtres de stage, professeurs, chefs de chantier) acceptant de témoigner : « Il/elle a refusé de signer un faux PV, a signalé un danger, a corrigé une erreur à ses frais ». \\ \hline
Sanctions & \textbf{Casier vierge / gestion transparente} & Extrait de casier judiciaire ; attestation employeur « aucune sanction disciplinaire » ; si sanction passée : explication écrite de la leçon tirée et mesures correctrices mises en place. \\ \hline
Engagement associatif & \textbf{Engagement citoyen / Ordre professionnel} & Adhésion étudiante à un Ordre ; participation à des commissions « éthique / développement durable » ; bénévolat technique (diagnostic gratuit pour une association). \\ \hline
\end{tabularx}
\end{center}

\begin{methode}[Construire son dossier « Intégrité » dès la Première technique]
\begin{enumerate}
    \item \textbf{Créez un dossier numérique sécurisé :} « Mon\_Integrite\_Pro ».
    \item \textbf{Archivez systématiquement :} chaque attestation de stage, chaque rapport de projet noté, chaque mail de signalement d'anomalie (même mineure), chaque formation suivie (secourisme, hygiène-sécurité-environnement, normes).
    \item \textbf{Demandez des recommandations ciblées :} à la fin de chaque stage ou projet, demandez au tuteur : « Pourriez-vous mentionner dans mon attestation mon respect des procédures de sécurité et de qualité ? ».
    \item \textbf{Documentez vos « non » :} notez (date, contexte, décision, issue) chaque refus d'une pratique douteuse (ex. : « 12/03 : refus de valider un bon de livraison de ciment sans certificat de conformité ; certificat fourni le 14/03 »). C'est votre capital confiance.
    \item \textbf{Rapprochez-vous des Ordres :} les Ordres professionnels proposent des statuts étudiants ou stagiaires. Cela signale votre engagement précoce.
\end{enumerate}
\end{methode}

\courssub{Stratégies de résistance opérationnelle : dire « non » techniquement}

Dire non à un supérieur, un client ou un fournisseur puissant demande du courage, mais surtout de la \textbf{technique} et de la \textbf{procédure}. L'improvisation émotionnelle affaiblit ; la rigueur procédurale protège.

\begin{propriete}[Arsenal juridique du résistant technique au Cameroun]
\begin{itemize}
    \item \textbf{Code du travail (Loi n°92/007) :} droit d'alerte et droit de retrait en cas de danger grave et imminent pour la vie ou la santé ; l'employeur ne peut pas sanctionner le travailleur qui use de ce droit de bonne foi.
    \item \textbf{Code de procédure pénale :} possibilité pour tout citoyen de dénoncer les crimes et délits (corruption, mise en danger, faux en écriture) au Procureur de la République.
    \item \textbf{Dispositif national de lutte contre la corruption :} la CONAC (Commission Nationale Anti-Corruption) reçoit les dénonciations ; les textes protègent les dénonciateurs de bonne foi contre les représailles.
    \item \textbf{Normes ISO 37001 (management anti-corruption) et ISO 37002 (systèmes d'alerte) :} elles exigent des canaux de signalement internes confidentiels au sein des organisations certifiées.
\end{itemize}
\end{propriete}

\begin{methode}[La « boîte à outils » de résistance sur le terrain]
\begin{enumerate}
    \item \textbf{Technique du « disque rayé » (assertivité calme) :} répéter inlassablement la règle technique, la norme ou le refus, sans s'énerver et sans sur-justifier.
        \begin{quote}
        \emph{— « Signe ce PV, le client attend. »}\\
        \emph{— « Je ne peux pas signer : le ferraillage n'est pas conforme au plan d'exécution (HA16 requis, HA10 posé). Je signerai quand ce sera conforme. »}\\
        \emph{— « On n'a pas le temps, fais-le pour cette fois. »}\\
        \emph{— « Je ne peux pas signer un PV non conforme, c'est ma responsabilité pénale. Je signerai quand ce sera conforme. »}
        \end{quote}
    \item \textbf{Renvoi à la procédure, à la hiérarchie, à la norme :} dépersonnaliser le conflit. Ce n'est pas « moi contre toi », c'est « la norme, le contrat, la loi contre cette action ».
        \begin{quote}
        \emph{« La procédure qualité impose un essai de compression à 7 jours avant décoffrage. Je ne peux pas déroger. Voici la procédure. »}
        \end{quote}
    \item \textbf{Écriture protectrice (traçabilité) :} \textbf{tout par écrit} : mail, note de service, courrier recommandé, registre de chantier.
        \begin{quote}
        \emph{Objet : Réserves techniques — Coulage dalle niveau 3 — Non-conformité ferraillage.\\
        Monsieur le Directeur, par la présente, je vous signale que le ferraillage posé (HA10) ne correspond pas aux plans d'exécution (HA16) validés par le bureau d'études. Conformément aux règles de l'art et au Code du travail, je ne peux valider le coulage. Je me tiens à disposition pour la mise en conformité. Fait à Yaoundé, le... Signature.}
        \end{quote}
    \item \textbf{Alerte éthique interne :} utiliser les canaux officiels de l'entreprise (boîte mail éthique, numéro vert, application dédiée), de manière anonyme ou identifiée. Conserver la preuve de l'envoi (capture d'écran, accusé de réception).
    \item \textbf{Démission éthique (dernier recours) :} si l'entreprise persiste dans une pratique illégale et dangereuse après alerte écrite et refus de signer, démissionner en motivant la lettre : « Refus de cautionner des pratiques mettant en danger la vie d'autrui et contraires à la déontologie professionnelle ». Cette démarche préserve l'honneur et la responsabilité du professionnel.
\end{enumerate}
\end{methode}

\courssub{Le coût de la non-qualité (CNQ) face à l'investissement intégrité : l'argument économique}

L'argument « l'intégrité coûte cher, fait perdre du temps et des marchés » est un mythe. L'analyse économique rigoureuse (modèle PAF : Prévention, Appréciation, Défaillance) démontre l'inverse.

\begin{popfigure}
\centering
\begin{tikzpicture}
\begin{axis}[
    ybar stacked,
    bar width=35pt,
    width=\linewidth,
    height=9cm,
    enlarge x limits=0.35,
    ylabel={Coût relatif (échelle indicative)},
    ylabel style={font=\small},
    xlabel={Stratégie de gestion de la qualité et de l'intégrité},
    xlabel style={font=\small},
    symbolic x coords={Integrite-Prevention, Non-integrite-Reparation},
    xtick=data,
    xticklabels={Intégrité et prévention, Non-intégrité et réparation},
    xticklabel style={font=\small, align=center},
    legend style={at={(0.5,-0.28)}, anchor=north, legend columns=1, font=\small, cells={anchor=west}},
    ymin=0, ymax=1050,
    axis lines=left,
    tick style={draw=none},
]
\addplot[fill=popgreen!70] coordinates {(Integrite-Prevention, 50) (Non-integrite-Reparation, 10)};
\addplot[fill=poporange!70] coordinates {(Integrite-Prevention, 30) (Non-integrite-Reparation, 80)};
\addplot[fill=poppink!70] coordinates {(Integrite-Prevention, 5) (Non-integrite-Reparation, 900)};
\legend{
    Coûts de prévention (formation, contrôle qualité, audits, assurance),
    Coûts de défaillance interne (rebuts, reprises, retards avant livraison),
    Coûts de défaillance externe (garanties, procès, ruine, radiation, vies humaines)
}
\end{axis}
\end{tikzpicture}
\legende{Le coût de la non-qualité selon le modèle PAF (Prévention -- Appréciation -- Défaillance). L'intégrité technique (colonne de gauche) investit en prévention (vert) pour réduire drastiquement les coûts de défaillance externe (rose), qui sont exponentiels. La non-intégrité (colonne de droite) économise sur la prévention mais paie au prix fort, souvent de la vie de l'entreprise et de la liberté des responsables. Valeurs indicatives, adaptées des travaux de Juran et Crosby.}
\end{popfigure}

\begin{aretenir}[L'équation économique de l'intégrité]
\[
\text{Coût total qualité} = \underbrace{C_{\text{prévention}} + C_{\text{appréciation}}}_{\text{investissement intégrité (maîtrisé, linéaire)}} + \underbrace{C_{\text{défaillance interne}} + C_{\text{défaillance externe}}}_{\text{coût du manquement (exponentiel, imprévisible)}}
\]
\begin{itemize}
    \item \textbf{Intégrité :} quand $C_{\text{prévention}}$ augmente, $C_{\text{défaillance externe}}$ s'effondre. Un franc investi en formation et en contrôle en économise des centaines en contentieux et en réparations.
    \item \textbf{Non-intégrité :} quand $C_{\text{prévention}}$ diminue, $C_{\text{défaillance externe}}$ explose : effondrement, procès, prison, fin de carrière.
\end{itemize}
L'intégrité n'est pas un luxe moral : c'est la \textbf{stratégie de gestion des risques la plus rentable} pour l'entreprise et la \textbf{seule véritable assurance-vie} pour le professionnel technique.
\end{aretenir}

\courssub{Exercice d'application : le dilemme du câblage « provisoire »}

\begin{exercice}[Mise en situation : pression hiérarchique contre norme électrique]
\textbf{Contexte :} vous êtes technicien électricien fraîchement embauché dans une PME d'installation électrique à Douala. Le chef d'équipe (votre supérieur direct) vous ordonne de câbler un tableau de distribution tertiaire avec du câble de section 2,5 mm\textsuperscript{2} pour des circuits de prises, alors que la protection amont est un disjoncteur 20 A et que la chute de tension dépasse la limite admise par la norme NF C 15-100 sur la longueur de la ligne. Le chef déclare : « C'est provisoire, le client n'a pas le budget pour la section supérieure, on repassera plus tard. Signe le PV de réception, on a besoin du chèque final. »

\textbf{Questions :}
\begin{enumerate}
    \item Identifiez les non-conformités techniques précises (protection, section, chute de tension).
    \item Qualifiez le conflit d'intérêts : qui gagne quoi ? Qui risque quoi ?
    \item Appliquez la méthode de gestion d'un conflit d'intérêts (4 étapes) à votre situation.
    \item Rédigez la note écrite (mail ou courrier) que vous adressez à votre hiérarchie pour formaliser votre refus de signer le PV, en citant les normes et textes concernés.
    \item Si la direction maintient son ordre, quelle est votre décision finale (droit de retrait, alerte interne, démission) ? Justifiez.
\end{enumerate}
\end{exercice}

\begin{corrige}[Exercice — Le dilemme du câblage « provisoire »]
\begin{enumerate}
    \item \textbf{Non-conformités :} (1) protection contre les surintensités inadaptée : un disjoncteur 20 A ne protège pas correctement un câble de 2,5 mm\textsuperscript{2} dont le courant admissible est limité ; risque d'échauffement et d'incendie. (2) Chute de tension supérieure à la limite fixée par la norme NF C 15-100 : dysfonctionnement des appareils, échauffement des moteurs. (3) Le « provisoire » non documenté, non daté et non tracé devient de fait une installation définitive non conforme.
    \item \textbf{Conflit d'intérêts :} l'entreprise et le chef d'équipe gagnent immédiatement (trésorerie, marge, rapidité). Le technicien risque sa responsabilité pénale (mise en danger d'autrui, faux PV), sa qualification future et sa conscience professionnelle. Le client risque l'incendie, le refus d'indemnisation par l'assurance (installation non conforme) et le surcoût d'une reprise ultérieure.
    \item \textbf{Méthode en 4 étapes :} (1) \emph{Identifier :} je suis le signataire du PV, donc le responsable technique ; mon intérêt personnel (garder mon emploi) entre en conflit avec mon obligation de sécurité. (2) \emph{Déclarer :} mail formel au conducteur de travaux ou au directeur : « Refus de signer un PV non conforme à la NF C 15-100 ». (3) \emph{S'abstenir :} ne pas signer ; quitter le poste si l'ordre persiste (droit de retrait, danger grave et imminent). (4) \emph{Tracer :} conserver le mail, photographier le disjoncteur 20 A et le câble de 2,5 mm\textsuperscript{2}, rédiger une note de chantier datée.
    \item \textbf{Note type :} reprendre le modèle « Écriture protectrice » de la méthode précédente, en citant la norme NF C 15-100 (protection des circuits et chute de tension), le droit d'alerte et de retrait du Code du travail, et la mise en danger d'autrui du Code pénal.
    \item \textbf{Décision finale :} \textbf{droit de retrait immédiat} (danger grave et imminent : incendie, électrocution), puis \textbf{signalement écrit} à l'inspection du travail et, si la pratique persiste, dénonciation aux autorités compétentes. En cas de licenciement abusif consécutif à l'alerte, contester devant l'inspection du travail. L'intégrité technique prime sur le confort d'un emploi : un emploi se retrouve, une vie perdue et une condamnation pénale ne s'effacent pas.
\end{enumerate}
\end{corrige}

\begin{attention}[Dernier rappel avant la synthèse]
L'intégrité technique n'est pas une option réservée « aux héros ». C'est la \textbf{condition sine qua non} de l'exercice légal et durable de votre métier. Chaque signature, chaque validation, chaque silence face à une non-conformité construit votre avenir : soit une carrière respectée, assurée et fière ; soit une carrière brisée, judiciaire et tragique. Le choix se fait \textbf{aujourd'hui}, sur les bancs de l'atelier, en stage, sur votre premier chantier. Choisissez le camp de la norme, de la vie et de l'honneur.
\end{attention}

\coursec{Synthèse évaluative et Engagement citoyen}

Nous avons découvert, dans les sections précédentes, que l'intégrité morale n'est pas une affaire de discours : elle s'incarne dans les gestes quotidiens du technicien et de l'ingénieur — refus de falsifier un rapport de chantier, respect des normes, honnêteté dans les mesures et les devis. Mais une leçon n'est acquise que si elle est \cle{révisée}, \cle{évaluée} et \cle{transformée en engagement personnel}. Cette dernière section a donc un triple objectif : vérifier que les savoirs sont bien installés, transformer les productions de la classe (les chartes) en outils vivants, et faire de chaque élève un acteur engagé de l'intégrité, pour son Probatoire ou son Baccalauréat technique, puis pour sa vie active. C'est le passage du « je sais » au « je m'engage ».

\courssub{Révision structurée : le grand QCM de la leçon}

La révision prend la forme d'un QCM collectif (sur papier ou sous forme de quiz type Kahoot si la salle est équipée). L'intérêt pédagogique est double : l'élève réactive ses connaissances dans un format proche de l'examen, et l'enseignant repère immédiatement les notions encore fragiles. Le QCM balaie les quatre champs de la leçon : les \cle{définitions} (intégrité, corruption, conflit d'intérêts), les \cle{institutions} (CONAC, CONSUPE, Audit, DGSN, tribunaux), les \cle{lois et procédures de signalement}, et les \cle{cas techniques} propres aux spécialités de la classe.

\begin{exoresolu}[Exemple de question de QCM corrigée]
Énoncé — Un contrôleur technique d'une entreprise de BTP constate que le ferraillage d'une dalle ne respecte pas les normes. Le chef de chantier lui propose 50\,000 FCFA pour « valider quand même ». Que doit faire le contrôleur ?\par
A. Accepter, car le chef de chantier est son supérieur hiérarchique.\par
B. Refuser l'argent, se taire et laisser faire pour garder son emploi.\par
C. Refuser l'argent, consigner l'anomalie par écrit et alerter sa hiérarchie ou les organes de contrôle.\par
D. Exiger une somme plus importante pour couvrir le risque.
\tcblower
Solution détaillée — La bonne réponse est la \textbf{C}. L'option A constitue une \cle{corruption passive} (recevoir un avantage pour accomplir un acte contraire à ses devoirs), réprimée par le Code pénal camerounais (Loi n°2011/002, art. 134). L'option B est une non-dénonciation qui engage la responsabilité morale et parfois pénale du technicien en cas d'accident. L'option D aggrave la faute (corruption active). La réponse C applique la démarche d'intégrité vue en section 5 : \textbf{refuser} (point de non-retour), \textbf{tracer} (consigner par écrit, preuve de bonne foi), \textbf{signaler} (hiérarchie, puis organes de contrôle si la hiérarchie est complice). L'intégrité technique protège ici des vies humaines : une dalle mal ferraillée peut s'effondrer.
\end{exoresolu}

\begin{attention}[Pièges fréquents au QCM]
Trois confusions reviennent chaque année : (1) confondre \cle{CONAC} (Commission Nationale Anti-Corruption, organe de prévention et de répression, créé par le décret n°2006/088 du 31 mars 2006) avec la \cle{CONSUPE} (Contrôle Supérieur de l'État, organe de contrôle de l'action administrative) — leurs missions se complètent mais ne se recouvrent pas ; (2) croire que la corruption ne concerne que l'argent : un service rendu, une note gonflée, un piston sont aussi des \cle{avantages indus} ; (3) penser que le signalement est une « délation » : la délation vise à nuire par intérêt personnel, le signalement (\cle{whistleblowing}) protège l'intérêt général et est protégé par la loi (Loi n°2011/002, art. 139). À l'examen, définir précisément les termes avant de répondre.
\end{attention}

\courssub{Restitution des chartes de classe : du papier au mur}

Les chartes de l'intégrité construites en TP (section 4) ne doivent pas dormir dans les cahiers. La restitution suit trois étapes :

\begin{enumerate}
\item \textbf{Présentation orale} : chaque groupe présente sa charte en 3 minutes (contexte, règles retenues, sanctions proposées).
\item \textbf{Vote de la classe} : la meilleure charte est élue selon quatre critères explicites — \cle{clarté} (compréhensible par tous), \cle{exhaustivité} (couvre examens, ateliers, stages, réseaux sociaux), \cle{applicabilité} (règles réalistes et vérifiables), \cle{courage} (ose nommer les vrais problèmes : fraude, racket, notes achetées).
\item \textbf{Affichage} : la charte lauréate est affichée dans la salle, et des copies sont proposées à la salle des professeurs, au CDI et au hall de l'établissement. L'affichage public transforme la charte en \emph{engagement visible} : on ne transgresse pas impunément une règle que l'on a soi-même écrite et affichée.
\end{enumerate}

\courssub{La lettre d'engagement d'intégrité : rédaction individuelle finale}

C'est l'activité centrale de cette section. Chaque élève rédige, sur une page, sa \emph{« Lettre d'engagement d'intégrité pour mon Bac Pro et ma vie active »}. L'intuition est simple : on tient mieux une promesse écrite, datée et signée qu'une bonne intention vague. La psychologie de l'engagement le confirme : l'écrit engage celui qui signe devant lui-même et devant les autres.

\begin{methode}[Rédiger un engagement personnel (la lettre)]
\begin{enumerate}
\item \textbf{Adopter un ton solennel} : commencer par « Je soussigné(e), [Nom, Prénom], élève de Première [spécialité]... ».
\item \textbf{Déclarer ses valeurs de référence} : honnêteté, travail, respect des normes, sens du service public (2 à 3 valeurs maximum, réellement choisies).
\item \textbf{Identifier ses zones de risque} selon sa spécialité : triche aux examens, falsification de mesures en atelier, plagiat de rapports de stage, soudure ou câblage bâclés, logiciels piratés...
\item \textbf{Formuler 3 engagements concrets}, positifs (actions : « Je signalerai toute tentative de corruption ») et négatifs (refus : « Je ne tricherai jamais à un examen »).
\item \textbf{Prévoir une stratégie en cas de pression} : qui alerter (professeur principal, chef d'établissement, CONAC via le 1517), comment dire non, quelles preuves conserver (captures d'écran, témoins, écrits).
\item \textbf{Dater, signer} et conserver la lettre dans son \cle{portfolio} de formation, pour la relire avant chaque examen et chaque stage.
\end{enumerate}
\end{methode}

\begin{attention}[L'engagement n'est pas un vœu pieux]
Une lettre qui se limite à « je serai honnête » ne vaut rien : elle ne dit ni quoi faire, ni quand, ni à quel prix. Un vrai engagement contient des \cle{points de non-retour}, c'est-à-dire des lignes rouges formulées à l'avance : « Si on me demande de falsifier un rapport, je refuse et j'alerte », « Si un examinateur me propose une fuite de sujet, je décline et je le signale ». Le point de non-retour protège : la décision est déjà prise \emph{avant} la tentation, donc la pression du moment ne peut plus la faire plier.
\end{attention}

\begin{savaistu}[Le Serment du Bachelier]
Au Cameroun, la cérémonie officielle de remise des diplômes du Baccalauréat comporte le \emph{Serment du Bachelier} : le lauréat s'engage notamment à \emph{servir la patrie avec honneur et probité}. Ce n'est pas du folklore : c'est un acte à la fois juridique et moral, qui fait de l'intégrité une condition de l'exercice citoyen et professionnel. En rédigeant ta lettre d'engagement en Première, tu t'entraînes à ce que tu prononceras publiquement le jour de ton diplôme.
\end{savaistu}

\courssub{Ouverture : l'intégrité, levier de l'émergence (Vision 2035)}

Pourquoi une leçon d'ECM parle-t-elle d'économie ? Parce que l'intégrité a un \emph{prix} — ou plutôt une \emph{valeur}. Le Cameroun s'est fixé, dans sa Vision 2035, l'ambition de devenir un \cle{pays émergent}. Or l'émergence repose sur trois piliers d'intégrité : une \textbf{administration intègre} (les services publics rendent des services, ils ne les vendent pas), des \textbf{entreprises intègres} (les marchés sont gagnés par la qualité, non par les dessous-de-table), des \textbf{citoyens intègres} (chacun refuse la petite corruption quotidienne). Sans ces trois piliers, les infrastructures se dégradent, les investisseurs fuient, les emplois ne se créent pas.

\begin{exemplebox}[L'intégrité attire l'investissement]
Les rapports de la Banque mondiale (indicateurs de gouvernance mondiale, \emph{Worldwide Governance Indicators}) placent le \emph{contrôle de la corruption} au premier rang des critères d'attractivité d'un pays. Les études comparatives montrent que les pays réputés intègres attirent jusqu'à \textbf{trois fois plus d'Investissements Directs Étrangers (IDE)} que les pays réputés corrompus, à ressources comparables. Concrètement : ton intégrité de futur technicien — un devis honnête, un chantier aux normes, une facture réelle — participe de la réputation du pays, donc des usines qui s'y installeront, donc de \emph{tes emplois futurs}. L'intégrité n'est pas un luxe moral : c'est une infrastructure économique.
\end{exemplebox}

La jeunesse technique a ici un rôle décisif : ce sont les techniciens et ingénieurs qui construiront les routes, les ponts, les réseaux électriques et numériques de 2035. Un pays émergent ne peut pas être bâti sur du béton trafiqué et des rapports falsifiés.

\courssub{Évaluation formative : s'auto-positionner et évaluer ses pairs}

L'évaluation formative ne sert pas à noter, mais à \emph{savoir où l'on en est} pour progresser. Chaque élève remplit le tableau de bord ci-dessous en cochant honnêtement sa position pour chaque savoir et chaque savoir-faire de la leçon. Puis, dans un second temps, chaque groupe évalue la charte d'un autre groupe (évaluation par les pairs) selon les quatre critères du vote : clarté, exhaustivité, applicabilité, courage.

\begin{popfigure}
\begin{center}
\begin{tikzpicture}[scale=0.95, every node/.style={font=\small}, transform shape]
% Titre
\node[font=\bfseries\large, text=popdark, anchor=north] at (0,4.0) {Tableau de bord compétences -- Leçon Intégrité Morale};
% En-têtes de colonnes
\foreach \x/\t in {-3.5/{Je maîtrise}, 0/{Je hésite}, 3.5/{Je ne sais pas}}{
  \node[fill=popblue, text=white, rounded corners=2pt, minimum width=3.0cm, minimum height=0.7cm] at (\x,3.2) {\t};
}
% Lignes savoirs
\node[fill=popblueL, rounded corners=2pt, minimum width=4.0cm, minimum height=0.7cm, anchor=east, text=popdark] at (-5.2,2.3) {\textbf{Savoirs}};
\foreach \y/\t in {1.5/{Concept d'intégrité}, 0.7/{Institutions (CONAC...)}, -0.1/{Intégrité au lycée}, -0.9/{Éthique professionnelle}}{
  \node[anchor=east, text=popdark] at (-5.2,\y) {\t};
  \foreach \x in {-3.5,0,3.5}{
    \draw[popline, thick] (\x-0.45,\y-0.2) rectangle (\x+0.45,\y+0.25);
  }
}
% Lignes savoir-faire
\node[fill=popgreenL, rounded corners=2pt, minimum width=4.0cm, minimum height=0.7cm, anchor=east, text=popdark] at (-5.2,-1.9) {\textbf{Savoir-faire}};
\foreach \y/\t in {-2.7/{Appliquer à des cas concrets}, -3.5/{Rédiger et justifier}, -4.3/{Construire des outils (charte)}}{
  \node[anchor=east, text=popdark] at (-5.2,\y) {\t};
  \foreach \x in {-3.5,0,3.5}{
    \draw[popline, thick] (\x-0.45,\y-0.2) rectangle (\x+0.45,\y+0.25);
  }
}
% Exemple de coches
\node[text=popgreen, font=\bfseries\large] at (-3.5,1.525) {$\times$};
\node[text=poporange, font=\bfseries\large] at (0,0.725) {$\times$};
% Consigne
\node[anchor=west, text=popmuted, font=\itshape\small, text width=10cm] at (-5.2,-5.2) {Consigne : cocher une seule case par ligne, en toute honnêteté (exemple de coches ci-dessus). Les cases « Je hésite » et « Je ne sais pas » désignent les notions à retravailler avant l'évaluation sommative.};
\end{tikzpicture}
\end{center}
\legende{Tableau de bord d'auto-évaluation : l'élève coche une case par ligne (savoirs en bleu, savoir-faire en vert).}
\end{popfigure}

\begin{exercice}[Pour aller plus loin]
\begin{enumerate}
\item Réponds au QCM de révision (10 questions) sans consulter ton cours, puis corrige-toi et reporte tes erreurs dans le tableau de bord.
\item Rédige ta lettre d'engagement d'intégrité (1 page) en suivant la méthode en 6 étapes ; fais-la relire par un pair avant de la signer.
\item En évaluant la charte d'un autre groupe, rédige un court paragraphe justifiant ta note sur le critère « courage » : la charte ose-t-elle nommer les vraies pratiques de fraude de l'établissement ?
\item Explique en dix lignes, avec l'exemple des IDE, pourquoi l'intégrité des techniciens est une condition de l'émergence du Cameroun en 2035.
\end{enumerate}
\end{exercice}

\begin{corrige}[Exercice 4 (indications)]
Le raisonnement attendu suit une chaîne causale en quatre maillons : (1) les investisseurs étrangers choisissent les pays où les règles sont respectées, car la corruption augmente les coûts et les risques ; (2) les pays intègres attirent ainsi jusqu'à trois fois plus d'IDE (indicateurs de gouvernance mondiale) ; (3) ces investissements créent des usines, des chantiers et donc des emplois techniques ; (4) or la réputation d'intégrité d'un pays se construit par les actes quotidiens de ses agents, ses entreprises et ses techniciens — devis honnêtes, chantiers aux normes, rapports véridiques. Conclusion : l'intégrité du futur technicien camerounais n'est pas seulement une vertu personnelle, c'est une contribution directe à l'attractivité du pays et à la création de ses propres emplois futurs, condition de l'émergence visée pour 2035.
\end{corrige}

\begin{aretenir}[L'essentiel de la section]
La leçon se clôt par trois gestes : \textbf{réviser} (QCM et tableau de bord pour situer ses acquis), \textbf{rendre public} (affichage et vote des chartes), \textbf{s'engager} (lettre personnelle datée, signée, avec des points de non-retour). L'intégrité morale n'est pas un savoir parmi d'autres : elle est la condition de la confiance, donc de l'investissement, de l'emploi et de l'émergence du Cameroun à l'horizon 2035. Le technicien intègre est un bâtisseur de la Nation.
\end{aretenir}

\newpage

\coursec{Activité d'intégration}

\courssub{Objectifs de l'activité}
Cette activité d'intégration vise à mobiliser l'ensemble des savoirs, savoir-faire et savoir-être développés au cours de la leçon sur l'intégrité morale, à travers une mise en situation professionnelle authentique. Elle permet à l'élève de :
\begin{itemize}
    \item \textbf{Identifier et qualifier} des faits de corruption et de fraude technique (faux en écriture publique, mise en danger de la vie d'autrui, concussion) dans le cadre d'un marché public.
    \item \textbf{Mobiliser le cadre juridique et institutionnel} camerounais (Constitution, Code pénal, Loi n°2011/022, CONAC, INSPECTION DES TRAVAUX PUBLICS) pour argumenter une défense ou une accusation.
    \item \textbf{Appliquer une démarche éthique} (courage moral, loyauté envers la profession, protection de l'intérêt général) face à un dilemme loyauté/hiérarchie vs intégrité/sécurité.
    \item \textbf{Rédiger des actes juridiques et techniques} (PV d'audience, jugement motivé, rapport de stage, fiche de leçon apprise) conformes aux standards de la communication professionnelle technique.
    \item \textbf{Construire une posture de lanceur d'alerte responsable} : procédure de signalement, protection du témoin, gestion des représailles.
\end{itemize}

\courssub{Situation : L'Affaire du Bâtiment Électrique — « Fondations au rabais »}
\label{sit:affaire_batiment}

\textbf{Contexte général.}
Le Lycée Technique Régional de \textbf{Bafoussam} (Région de l'Ouest) lance un appel d'offres pour la construction d'un nouveau bloc de 12 salles de classe et laboratoires (Marché Public n° 045/MINEDUB/CIPM/2025). L'entreprise \textbf{« BTP Excellence SARL »}, adjudicataire provisoire, est en phase d'exécution des fondations. Le chantier est suivi par le Bureau de Contrôle \textbf{« Ingénierie \& Contrôle Cameroun (ICC) »} et la Mission de Contrôle du Ministère des Travaux Publics (MINTP).

\textbf{Acteurs principaux.}
\begin{itemize}
    \item \textbf{M. TCHINDA Kevin}, élève de \textbf{1re Tech F3 (Bâtiment)}, stagiaire « Conducteur de Travaux » chez BTP Excellence SARL (convention de stage signée avec le Lycée Technique de Bafoussam).
    \item \textbf{M. FOPA Armand}, Chef de Chantier, ingénieur des travaux publics, 15 ans d'expérience, tuteur de stage de Kevin.
    \item \textbf{Mme MEKONTSO Sylvie}, Représentante de la \textbf{CONAC (Commission Nationale Anti-Corruption)}, antenne Région Ouest, missionnée sur alerte anonyme.
    \item \textbf{M. LE JUGE / PROCUREUR}, Président du Tribunal de Grande Instance de Bafoussam (simulé).
    \item \textbf{Le Jury (Observateurs)}, composé d'élèves de 1re Tech F2 (Génie Civil), F4 (Topographie) et F3.
\end{itemize}

\textbf{Les faits techniques et chronologiques.}
\begin{enumerate}
    \item \textbf{Spécifications du CCTP (Cahier des Clauses Techniques Particulières) :} Béton de fondation : \textbf{Béton Armé BA 25 (C25/30)}, dosage minimal en ciment \textbf{300 kg/m$^3$} (classe d'exposition XC2), résistance caractéristique à 28 jours $f_{ck} = 25$ MPa. Écart type admis $\sigma = 4$ MPa. Formule validée par le laboratoire \textbf{LABO-GC} (agréé MINTP).
    \item \textbf{Réalité du chantier (Semaine du 13 au 17 janvier 2025) :} Trois camions toupies (centrale « Béton Moderne » à Bafoussam) livrent le béton pour les semelles filantes du bloc A.
    \item \textbf{Constat de Kevin (Stagiaire) :} En vérifiant les \textbf{bons de livraison (tickets de centrale)} n° BM-2025-0114, BM-2025-0115, BM-2025-0116 (photos à l'appui, horodatées 14/01/2025 09:12, 10:45, 12:30), il relève :
    \begin{itemize}
        \item Dosage ciment indiqué : \textbf{200 kg/m$^3$} (au lieu de 300 kg/m$^3$).
        \item Adjuvant réducteur d'eau : \textbf{0,8 \%} (au lieu de 1,2 \% prescrit).
        \item Affaissement (Slump) mesuré à l'arrivée : \textbf{180 mm} (classe S4, cohérent avec adjuvant), mais \textbf{temps de transit} : 1h45 (supérieur à 1h max préconisée par forte chaleur).
    \end{itemize}
    \item \textbf{Réaction du Chef de Chantier (M. FOPA) :} Alerté par Kevin (« Chef, le dosage n'est pas bon, on risque le refus du contrôleur »), M. FOPA répond : \emph{« Kevin, on a un retard de 3 semaines sur le planning. Le patron a négocié avec le fournisseur pour baisser le coût. Le béton 200 kg tient la route pour des semelles, on a déjà coulé 50 m$^3$. Le contrôleur passe vendredi, on a fait « bonne figure » en surface. Signe le PV de réception des fondations (PV n° REC-003) et on en parle plus. Ton rapport de stage en dépend. »}
    \item \textbf{Action de Kevin :} Refus de signer le PV. Photographie des 3 tickets + du PV vierge. Envoi anonyme (via boîte mail CONAC « Alerte Citoyen ») le 15/01/2025 à 20h00. Copie à son professeur tuteur (M. DJOUMESSI).
    \item \textbf{Intervention CONAC / MINTP :} Le 16/01, Mme MEKONTSO (CONAC) et un Ingénieur du MINTP (Inspection Régionale) descendent sur chantier. Prélèvement de 6 éprouvettes (carottes) sur les semelles coulées les 13 et 14. Saisie des tickets de centrale originaux (vérification numéro de série vs duplicata).
    \item \textbf{Résultats laboratoire (J+7, 23/01) :} Résistance moyenne à la compression sur carottes (Ø 100 mm) : \textbf{14,2 MPa} (Écart type 2,1 MPa). Soit \textbf{$< f_{ck} - 4$ MPa} (non conforme grave). Teneur en chlorures : normale. Carbonatation superficielle accélérée.
\end{enumerate}

\textbf{Enjeux.}
\begin{itemize}
    \item \textbf{Sécurité :} Risque d'effondrement différé du bâtiment (école accueillant 600 élèves), corrosion prématurée des armatures (dosage ciment faible = porosité + carbonatation rapide).
    \item \textbf{Financier :} Économie ciment estimée : $(300-200) \times 50 m^3 \times 55\,000$ FCFA/tonne $\approx$ \textbf{275\,000 FCFA} économisés illégalement sur ce tronçon. Marché total : 450 millions FCFA.
    \item \textbf{Disciplinaire / Pénal :} Faux en écriture publique (PV de réception, tickets), Mise en danger de la vie d'autrui (Art. 282 Code Pénal), Corruption passive/active (Art. 134, 135 Code Pénal ; Loi 2011/022), Concussion.
    \item \textbf{Humain :} Pression sur stagiaire (chantage note de stage), protection du lanceur d'alerte (Art. 38 Loi 2011/022).
\end{itemize}

\courssub{Consignes}
\label{consignes}

\textbf{Durée totale : 2 séances (3h). Travail en groupes de 5-6 élèves.}

\emph{\textbf{Séance 1 : Préparation (30 min) + Audience simulée (45 min)}}

\noindent \textbf{Groupe A — Défense}\ : préparer les arguments de la défense de M. TCHINDA Kevin (lanceur d'alerte), en s'appuyant sur le cadre juridique cité plus haut.

\noindent \textbf{Groupe B — Accusation}\ : préparer l'argumentation de l'entreprise BTP Excellence SARL et de la hiérarchie.

\noindent \textbf{Groupe C — Tribunal}\ : conduire l'audience, rédiger le PV et le jugement motivé.


\newpage

\coursec{Méthodes et exercices résolus}

\coursec{L'intégrité morale}

\courssub{L'intégrité morale : concept, piliers et enjeux}

\begin{definition}[Intégrité morale]
L'\cle{intégrité morale} est la \textbf{conformité constante entre les paroles, les actes et les valeurs éthiques}, même en l'absence de témoin. Elle implique la \cle{probité} (honnêteté dans la gestion des biens publics ou privés), la \cle{transparence} (clarté dans les décisions et reddition de comptes) et la \cle{cohérence} éthique (alignement valeurs/pratiques). C'est une vertu cardinale du citoyen et du professionnel technique.
\end{definition}

\begin{propriete}[Les quatre piliers de l'intégrité (Mnémotechnique \cle{HPTR})]
\begin{itemize}
    \item \textbf{H} -- \textbf{Honnêteté} : Dire la vérité, refuser la tricherie, le mensonge et la fraude sous toutes ses formes.
    \item \textbf{P} -- \textbf{Probité} : Gestion désintéressée des affaires publiques ou privées ; refus catégorique de la corruption, du détournement, de la concussion et du népotisme.
    \item \textbf{T} -- \textbf{Transparence} : Clarté dans les procédures de décision, accès à l'information, obligation de reddition de comptes (\cle{redevabilité}).
    \item \textbf{R} -- \textbf{Responsabilité} : Assumer pleinement les conséquences de ses choix, de ses actes et de ses omissions, tant sur le plan disciplinaire que pénal.
\end{itemize}
\end{propriete}

\begin{aretenir}[Enjeux de l'intégrité pour le développement]
\begin{itemize}
    \item \textbf{Confiance sociale :} Condition \emph{sine qua non} du contrat social (citoyen $\leftrightarrow$ État, client $\leftrightarrow$ entreprise).
    \item \textbf{Efficacité administrative et économique :} La probité préserve les deniers publics (écoles, hôpitaux, routes) ; la transparence attire les investissements (IDE, notation pays).
    \item \textbf{Qualité technique et sécurité :} L'honnêteté du technicien (respect des normes NF, ISO) garantit la solidité des ouvrages et la protection des vies humaines (Art. 289 Code Pénal : mise en danger d'autrui).
    \item \textbf{Mérite et justice :} L'intégrité dans l'évaluation (concours, examens, marchés publics) assure que les compétences réelles sont récompensées, fondement de l'innovation et de la compétitivité.
\end{itemize}
\end{aretenir}

\begin{methode}[Analyser le concept d'intégrité morale : définition, piliers et enjeux]
    \textbf{Objectif :} Maîtriser la définition rigoureuse, identifier les piliers (honnêteté, probité, transparence, responsabilité) et expliquer les enjeux (confiance sociale, efficacité administrative, développement).
    \begin{enumerate}
        \item \textbf{Définir} : L'intégrité morale est la \cle{conformité constante entre les paroles, les actes et les valeurs éthiques}, même en l'absence de témoin. Elle implique la \cle{probité} (honnêteté dans la gestion des biens publics/privés) et la \cle{cohérence} éthique.
        \item \textbf{Lister les 4 piliers} (mnémotechnique \cle{HPTR}) :
            \begin{itemize}
                \item \textbf{H}onnêteté : Dire la vérité, refuser la tricherie.
                \item \textbf{P}robité : Gestion désintéressée, refus de la corruption/détournement.
                \item \textbf{T}ransparence : Clarté dans les décisions, reddition de comptes.
                \item \textbf{R}esponsabilité : Assumer les conséquences de ses choix.
            \end{itemize}
        \item \textbf{Expliquer les enjeux} : Relier chaque pilier à une conséquence concrète (ex: Probité $\rightarrow$ Préservation deniers publics $\rightarrow$ Équipements écoles/hôpitaux).
        \item \textbf{Structurer la réponse} : Définition $\rightarrow$ Piliers (exemples) $\rightarrow$ Enjeux (collectif/individuel).
    \end{enumerate}
\end{methode}

\begin{exoresolu}[Analyse conceptuelle : L'intégrité face à la tentation]
    \textbf{Énoncé :} M. KAMGA, chef de chantier dans une entreprise de BTP, reçoit un appel d'offres pour la réfection d'un pont. Le fournisseur de ciment lui propose discrètement un « pourcentage » (5\% du marché) s'il choisit sa société, dont le ciment est légèrement moins résistant que la norme. M. KAMGA refuse et signale la tentative à sa hiérarchie.
    \textbf{Questions :}
    \begin{enumerate}
        \item Qualifiez l'attitude du fournisseur et celle de M. KAMGA en utilisant le vocabulaire précis de l'ECM.
        \item Identifiez deux piliers de l'intégrité morale illustrés par le refus de M. KAMGA. Justifiez.
        \item Expliquez en deux arguments pourquoi l'intégrité de M. KAMGA protège l'intérêt général.
    \end{enumerate}
    \tcblower
    \textbf{Solution détaillée :}
    \begin{enumerate}
        \item \textbf{Attitude du fournisseur :} C'est une tentative de \cle{corruption active} (offre d'un avantage indu pour obtenir un avantage illégitime). C'est aussi un manquement à la \cle{probité} et à la \cle{transparence} commerciale.
            \textbf{Attitude de M. KAMGA :} C'est un acte d'\cle{intégrité morale} (conformité à l'éthique professionnelle malgré la tentation financière). C'est spécifiquement un acte de \cle{probité} (refus de gain illicite) et de \cle{responsabilité} (signalement pour protéger l'institution).
        \item \textbf{Pilier 1 : La Probité.} Justification : M. KAMGA refuse un avantage personnel illicite (5\%) qui aurait détourné la commande publique de son objet (qualité du pont). Il place l'intérêt de l'ouvrage au-dessus de son intérêt pécuniaire immédiat.
            \textbf{Pilier 2 : La Responsabilité.} Justification : Il ne se contente pas de refuser ; il \emph{signale} la tentative. Il assume son devoir de protéger l'entreprise et la collectivité contre un risque structurel (ciment non conforme) et juridique.
        \item \textbf{Argument 1 (Sécurité/Qualité) :} En refusant le ciment défectueux, M. KAMGA garantit la \cle{solidité de l'ouvrage d'art}. Un pont effondré coûterait des vies humaines et des milliards de FCFA à la collectivité, bien supérieurs au gain personnel refusé.
        \item \textbf{Argument 2 (Confiance/Exemplarité) :} Son signalement renforce la \cle{confiance} dans l'institution entreprise et l'État. Il crée un \cle{précédent éthique} : la corruption n'est pas une fatalité, elle se combat par des actes concrets. Cela préserve l'image de la profession de technicien.
    \end{enumerate}
    \textbf{Conclusion :} L'intégrité morale n'est pas une abstraction ; c'est une compétence technique vitale. Le technicien intègre est le premier garant de la qualité et de la durabilité des infrastructures nationales.
\end{exoresolu}

\courssub{Le dispositif institutionnel camerounais de promotion et de préservation}

\begin{definition}[Cadre juridique de référence]
\begin{itemize}
    \item \textbf{Loi n°2006/014 du 29 décembre 2006} : Création, organisation et fonctionnement de la \cle{CONAC}.
    \item \textbf{Loi n°2011/002 du 6 mai 2011} modifiée par la \textbf{Loi n°2018/011 du 11 juillet 2018} : Organisation et fonctionnement du \cle{Tribunal Criminel Spécial (TCS)}.
    \item \textbf{Décret n°2018/191 du 2 mars 2018} : Organisation et fonctionnement de la \cle{CONSUPE}.
    \item \textbf{Code Pénal (Loi n°2016/007 du 12 juillet 2016)} : Articles 134 (Corruption), 134-1 (Concussion), 134-2 (Trafic d'influence), 289 (Mise en danger d'autrui), 184 (Détournement de deniers publics).
    \item \textbf{Loi n°2019/010 du 24 décembre 2019} : Protection des lanceurs d'alerte.
\end{itemize}
\end{definition}

\begin{propriete}[Triptyque institutionnel majeur (Programme Officiel)]
\begin{center}
\begin{tabularx}{\linewidth}{|l|X|X|X|}\hline
\textbf{Institution} & \textbf{Mission principale} & \textbf{Acteurs clés / Particularités} \\ \hline
\cle{CONAC}\newline(Commission Nationale Anti-Corruption) & \textbf{Prévention, Éducation, Veille, Enquête administrative} & Rapporte au Président de la République. Pas de pouvoir de juger. Publie le rapport annuel sur l'état de la corruption. \\ \hline
\cle{CONSUPE}\newline(Commission Nationale de Surveillance des Marchés Publics) & \textbf{Régulation, Contrôle de la commande publique} & Veille à la transparence, l'équité, la concurrence. Peut suspendre/annuler une procédure. Compétence exclusive : marchés publics. \\ \hline
\cle{TCS}\newline(Tribunal Criminel Spécial) & \textbf{Répression judiciaire (Fort montants)} & Juge les crimes de corruption, détournement, concussion $\ge$ \textbf{50 millions FCFA}. Peines : Prison (10 ans à vie), amendes, confiscation. \\ \hline
\end{tabularx}
\end{center}
\end{propriete}

\begin{aretenir}[Acteurs de proximité et contrôle complémentaire]
\begin{itemize}
    \item \textbf{Contrôle Supérieur de l'État} / \textbf{Inspection Générale des Services (IGS)} : Contrôle a posteriori de la gestion des services publics.
    \item \textbf{Médiateur de la République} : Défense des droits des usagers du service public (médiation, recommandations).
    \item \textbf{Pôle Judiciaire Financier} / \textbf{TGI} : Répression des infractions financières < 50 millions FCFA.
    \item \textbf{Société civile / ONG} : Transparency International Cameroun (rapports, plaidoyer), Unions de consommateurs.
    \item \textbf{École, Famille, Médias} : Éducation de base aux valeurs (premier rempart).
\end{itemize}
\end{aretenir}

\begin{methode}[Identifier les institutions camerounaises de promotion et de préservation de l'intégrité]
    \textbf{Objectif :} Citer les institutions, distinguer leurs missions (prévention, répression, éducation) et rattacher un scénario à l'institution compétente.
    \begin{enumerate}
        \item \textbf{Mémoriser le triptyque institutionnel majeur (Programme Officiel)} :
            \begin{itemize}
                \item \cle{CONAC} (Commission Nationale Anti-Corruption) : \textbf{Prévention}, éducation, veille, enquêtes administratives, rapport annuel au Président.
                \item \cle{CONSUPE} (Commission Nationale de Surveillance des Marchés Publics) : \textbf{Régulation/Contrôle} de la commande publique (transparence, équité).
                \item \cle{Tribunal Criminel Spécial (TCS)} : \textbf{Répression} judiciaire des détournements $\ge$ 50 millions FCFA (forts montants).
            \end{itemize}
        \item \textbf{Compléter avec les acteurs de proximité} : Inspection Générale des Services (IGS), Contrôle Supérieur de l'État, Médiateur de la République, ONG (Transparency International Cameroun), Écoles/Familles (éducation de base).
        \item \textbf{Classer par fonction} :
            \begin{itemize}
                \item \textit{Prévention/Éducation} : CONAC, École, Famille, Médias.
                \item \textit{Contrôle/Enquête administrative} : CONAC, CONSUPE, IGS, Contrôle Supérieur.
                \item \textit{Répression/Jugement} : TCS, Tribunaux ordinaires, Pôle judiciaire financier.
            \end{itemize}
        \item \textbf{Réflexe Bac/Probatoire} : Ne jamais confondre CONAC (enquête administrative, pas de pouvoir de juger) et TCS (jugement). CONSUPE = Marchés publics uniquement. Seuil TCS = 50 millions FCFA.
    \end{enumerate}
\end{methode}

\begin{exoresolu}[Mise en situation institutionnelle : Qui fait quoi ?]
    \textbf{Énoncé :} Trois situations sont observées dans la ville de Douala :
    \begin{enumerate}
        \item \textbf{Situation A :} Un chef de service exige 10 000 FCFA pour signer un certificat de scolarité. Un parent d'élève porte l'affaire à une structure qui organise une sensibilisation dans l'école et ouvre une enquête administrative.
        \item \textbf{Situation B :} Une mairie lance un appel d'offres pour la construction d'un marché. Une entreprise évincée soupçonne un favoritisme dans l'attribution. Elle saisit l'organe de régulation de la commande publique.
        \item \textbf{Situation C :} Un Directeur Général est soupçonné d'avoir détourné 200 millions FCFA sur le budget d'un projet hydraulique. Le dossier est transmis à une juridiction spéciale pour jugement.
    \end{enumerate}
    \textbf{Questions :} Pour chaque situation, citez l'institution camerounaise compétente (sigle + nom complet) et précisez sa mission principale illustrée ici (Prévention, Régulation, ou Répression).
    \tcblower
    \textbf{Solution détaillée :}
    \begin{enumerate}
        \item \textbf{Situation A :}
            \begin{itemize}
                \item \textbf{Institution :} \cle{CONAC} — Commission Nationale Anti-Corruption.
                \item \textbf{Mission illustrée :} \textbf{Prévention et Enquête administrative}. La sensibilisation relève de la prévention/éducation ; l'enquête sur le fonctionnaire relève du pouvoir d'investigation administrative de la CONAC (Art. 2 Loi 2006/014).
            \end{itemize}
        \item \textbf{Situation B :}
            \begin{itemize}
                \item \textbf{Institution :} \cle{CONSUPE} — Commission Nationale de Surveillance des Marchés Publics.
                \item \textbf{Mission illustrée :} \textbf{Régulation et Contrôle de la commande publique}. La CONSUPE veille à la transparence et à l'équité des procédures d'attribution des marchés publics (Décret 2018/191). Elle peut annuler une procédure irrégulière.
            \end{itemize}
        \item \textbf{Situation C :}
            \begin{itemize}
                \item \textbf{Institution :} \cle{TCS} — Tribunal Criminel Spécial.
                \item \textbf{Mission illustrée :} \textbf{Répression judiciaire}. Le montant (200 millions FCFA) dépasse largement le seuil de 50 millions FCFA. Le TCS est la juridiction compétente pour juger les crimes de corruption, détournement de deniers publics et concussion pour les forts montants (Loi 2011/002 modifiée).
            \end{itemize}
    \end{enumerate}
    \textbf{Point de vigilance :} Si le montant était 10 millions (Situation C), ce serait le \textbf{Tribunal de Grande Instance (TGI)} ou le \textbf{Pôle Judiciaire Financier} (selon organisation), mais \emph{pas} le TCS. La distinction des seuils est cruciale au Probatoire/Bac.
\end{exoresolu}

\courssub{TD2 : L'intégrité morale en milieu scolaire - Diagnostic et Analyse}

\begin{experience}[TD2 : Diagnostic d'intégrité en milieu scolaire - Protocole]
\textbf{Objectif :} Appliquer la grille d'analyse « Acteur / Comportement / Valeur violée / Conséquence / Remède » à des situations scolaires concrètes.
\textbf{Matériel :} Études de cas (tricherie, racket, corruption de notes, vol, harcèlement, dégradation), grille d'analyse vierge, Règlement intérieur type, Code Pénal (extraits).
\textbf{Déroulement (3 séances) :}
\begin{enumerate}
    \item \textbf{Séance 1 (Apports) :} Rappel conceptuel (intégrité, piliers HPTR) + Présentation de la grille + Lecture du Règlement intérieur (sanctions disciplinaires) + Rappel sanctions pénales (Art 134 CP).
    \item \textbf{Séance 2 (Travail de groupe) :} Groupes de 4-5. Chaque groupe analyse un cas différent (Cas A : Tricherie au Bac/Probatoire ; Cas B : Corruption notes/surveillant ; Cas C : Racket/Extorsion ; Cas D : Vol matériel scolaire ; Cas E : Harcèlement/Cyberharcèlement). Remplissage grille.
    \item \textbf{Séance 3 (Restitution/Débat) :} Affichage grilles. Débat contradictoire : « La sanction éducative suffit-elle ? » « Comment protéger le lanceur d'alerte élève ? » Synthèse prof : Charte de classe engageante.
\end{enumerate}
\textbf{Compétences évaluées :} Esprit critique, vocabulaire juridique précis (corruption active/passive, faux, abus de fonction), proposition de solutions réalistes (numérisation, traçabilité, protection témoins).
\end{experience}

\begin{methode}[Réaliser un diagnostic d'intégrité en milieu scolaire (TD2) : Grille d'analyse]
    \textbf{Objectif :} Appliquer la grille « Acteur / Comportement / Valeur violée / Conséquence / Remède » à une situation scolaire concrète (tricherie, racket, corruption de notes, vol, harcèlement).
    \begin{enumerate}
        \item \textbf{Identifier l'acte} : Le décrire factuellement (qui ? quoi ? où ? quand ?) sans juger.
        \item \textbf{Qualifier la faute} : Rattacher à une violation de l'intégrité : \cle{Tricherie} (fraude examen), \cle{Corruption} (argent/notes), \cle{Racket} (extorsion), \cle{Vol}, \cle{Harcèlement}.
        \item \textbf{Nommer la valeur bafouée} : Honnêteté, Mérite, Respect, Justice, Solidarité.
        \item \textbf{Analyser les conséquences} : Niveau individuel (sanction, échec, culpabilité, casier) + Niveau collectif (climat de méfiance, dévalorisation du diplôme, injustice pour les travailleurs).
        \item \textbf{Proposer des remèdes cohérents} : Sanction éducative (réparation), Prévention (charte, sensibilisation), Institutionnel (règlement intérieur, implication vie scolaire, numérisation sécurisée).
    \end{enumerate}
\end{methode}

\begin{exoresolu}[Diagnostic scolaire : L'affaire des « notes moyennées »]
    \textbf{Énoncé :} Dans un lycée technique, une enquête de la vie scolaire révèle qu'un groupe d'élèves de Terminale verse 5 000 FCFA par trimestre à un surveillant général adjoint pour que ce dernier « arrange » les bulletins : absences effacées, moyennes majorées de 2 à 3 points. Le surveillant modifie les fichiers Excel avant impression.
    \textbf{Questions :} En utilisant la grille d'analyse, produisez un diagnostic complet de cette situation.
    \tcblower
    \textbf{Solution détaillée (Grille de diagnostic) :}

    \begin{center}
    \begin{tabularx}{\linewidth}{|l|X|}\hline
    \textbf{Critère} & \textbf{Analyse} \\ \hline
    \textbf{Acteurs} & Élèves (corrupteurs actifs), Surveillant général adjoint (corrompu passif, faussaire), Administration (responsabilité de contrôle défaillante). \\ \hline
    \textbf{Comportement} & \cle{Corruption active} (élèves : Art 134 CP) / \cle{Corruption passive} + \cle{Faux en écriture privée/public} (surveillant : modification frauduleuse de documents officiels scolaires, Art 121-123 CP). \\ \hline
    \textbf{Valeurs violées} & \begin{itemize}\item \textbf{Mérite :} La note ne reflète plus le travail réel.
    \item \textbf{Justice/Équité :} Injustice flagrante vis-à-vis des élèves travailleurs.
    \item \textbf{Honnêteté/Probité :} Trahison de la mission de service public éducatif.
    \item \textbf{Responsabilité :} Refus d'assumer son niveau réel ; complicité adulte/enfant.
    \end{itemize} \\ \hline
    \textbf{Conséquences} & \textbf{Individuel :} Risques pénaux (Art 134 CP : corruption = 5 à 10 ans prison + amende), exclusion définitive, casier judiciaire, perte confiance familiale.
    \textbf{Collectif :} \cle{Dévalorisation du diplôme} (Bac/Probatoire suspect aux yeux employeurs/universités), \cle{Climat de cynisme} (« tout s'achète »), \cle{Incompétence technique} future (élèves non formés, danger sur chantiers). \\ \hline
    \textbf{Remèdes proposés} & \textbf{Sanction :} Conseil de discipline $\rightarrow$ Exclusion élèves / Révocation surveillant + Saisine Procureur (TCS si montant global fort, sinon TGI/Pôle Financier).
    \textbf{Prévention :} Numérisation sécurisée des notes (accès traçable, horodatage, double validation), Charte d'intégrité signée à la rentrée, Formation continue surveillants (éthique, droit), Boîte à idées/alertes anonyme (physique/numérique). \\ \hline
    \end{tabularx}
    \end{center}

    \textbf{Synthèse :} Ce n'est pas une « aide entre camarades », c'est un \cle{système mafieux} qui gangrène la certification. L'intégrité du diplôme technique exige une \cle{traçabilité numérique} et une \cle{tolérance zéro} administrative.
\end{exoresolu}

\courssub{TP : Construction d'outils de promotion de l'intégrité (Cartes, Croquis, Chartes)}

\begin{experience}[TP : Élaboration d'une Charte de classe « Zéro Corruption / Excellence »]
\textbf{Objectif :} Produire un outil normatif, pédagogique et juridiquement solide, applicable immédiatement dans la classe.
\textbf{Durée :} 2 séances (1h rédaction collaborative + 1h adoption/signature).
\textbf{Consignes :}
\begin{enumerate}
    \item \textbf{Diagnostic partagé (15 min) :} Brainstorming : « Quels sont les 5 comportements qui tuent le mérite dans NOTRE classe ? » (Tricherie, plagiat, absentéisme complice, racket, dégradation matériel).
    \item \textbf{Rédaction par îlots (30 min) :} Chaque îlot rédige 2 articles (Forme : « Je m'engage à... » ou « Il est interdit de... »). Contraintes : Citer la loi (Code Pénal, Règlement intérieur), prévoir sanction graduée + réparation.
    \item \textbf{Mise en commun / Vote (15 min) :} Lecture, amendements, vote à main levée (majorité 2/3).
    \item \textbf{Formalisation (Séance 2) :} Mise en page (Logo établissement, Titre, Préambule, Articles, Sanctions, Signatures : Élèves, Parents, Professeur principal, Proviseur). Affichage salle de classe + Version numérique (QR Code sur ENT).
\end{enumerate}
\textbf{Livrable :} Charte signée, affichée, diffusée aux parents.
\end{experience}

\begin{methode}[Construire un outil de promotion de l'intégrité : Charte / Affiche / Règlement]
    \textbf{Objectif :} Rédiger un texte normatif court, clair, positif et engageant (Charte de classe, Affiche « Zéro Corruption », Règlement intérieur type).
    \begin{enumerate}
        \item \textbf{Titre accrocheur} : Ex: « Charte de l'Excellence et de l'Honnêteté - Promotion 2025 ».
        \item \textbf{Préambule (Pourquoi)} : 2-3 lignes rappelant le droit à une évaluation juste et le devoir de mérite.
        \item \textbf{Engagements (Articles) :} Rédiger 5 à 7 articles à la forme \textbf{« Je m'engage à... »} (positif) ou \textbf{« Il est interdit de... »} (négatif, pour les interdictions graves).
            \begin{itemize}
                \item Cibler : Examens (pas de fraude), Devoirs (pas de plagiat), Vie scolaire (pas de racket), Biens publics (pas de dégradation), Signalement (devoir d'alerte).
            \end{itemize}
        \item \textbf{Sanctions/Réparations :} Préciser l'échelle (Avertissement $\rightarrow$ Conseil de discipline $\rightarrow$ Poursuites pénales) et la dimension \textbf{réparatrice} (excuses publiques, travail d'intérêt général).
        \item \textbf{Signatures :} Élèves, Parents, Professeur principal, Chef d'établissement. Date.
    \end{enumerate}
\end{methode}

\begin{exoresolu}[Rédaction : Article de Charte « Zéro Tricherie »]
    \textbf{Énoncé :} Vous êtes délégué de classe de 1ère TEE (Électrotechnique). Le professeur principal vous demande de rédiger l'\textbf{Article 3} de la Charte de classe consacré aux \textbf{évaluations et examens}, ainsi que la \textbf{clause de sanction associée}. L'article doit être juridiquement solide (rappel loi), pédagogique et applicable immédiatement.
    \tcblower
    \textbf{Solution détaillée :}

    \textbf{Article 3 — Intégrité des Évaluations (Extrait Charte 1ère TEE)}

    \textbf{Article 3.1 (Devoir de sincérité) :} Tout élève s'engage à composer ses devoirs, interrogations et examens \textbf{en toute autonomie intellectuelle}, sans recourir à la fraude (copie sur voisin, téléphone, « chuettes », communication, substitution de copie).

    \textbf{Article 3.2 (Interdiction de la corruption) :} Toute tentative d'obtenir une note par l'argent, un cadeau, un service ou une menace envers un personnel de l'établissement constitue un acte de \textbf{corruption} passible des sanctions disciplinaires \textbf{et pénales} (Art. 134 et suivants du Code Pénal Camerounais).

    \textbf{Article 3.3 (Protection du mérite / Devoir d'alerte) :} Tout élève témoin de fraude ou de tentative de corruption a le \textbf{devoir moral} de le signaler au professeur ou à la vie scolaire (anonymat garanti), pour protéger la valeur de son propre diplôme.

    \textbf{Clause de Sanction Associée - Article 12 (Échelle graduée)}

    \begin{enumerate}
        \item \textbf{Niveau 1 (Tentative/Flagrance mineure) :} Note de 0/20 à l'épreuve + Avertissement écrit à la famille + Travail d'intérêt général (TIG) 2h (ex: rangement laboratoire, nettoyage atelier).
        \item \textbf{Niveau 2 (Fraude avérée / Récidive) :} 0/20 à l'épreuve + Exclusion temporaire (3 à 8 jours) + Conseil de discipline + Inscription au bulletin « Fraude aux examens ».
        \item \textbf{Niveau 3 (Corruption / Faux / Bande organisée) :} Saisine immédiate du \textbf{Procureur de la République} (infraction pénale) + Exclusion définitive de l'établissement (Décision Ministre/Proviseur) + Invalidation de l'année scolaire.
    \end{enumerate}
    \textbf{Disposition réparatrice :} Toute sanction s'accompagne d'un \textbf{entretien réflexif} avec le CPE/Proviseur sur « Pourquoi j'ai triché ? Comment réussir honnêtement ? ».

    \textbf{Justification pédagogique :} L'article cite la \cle{loi pénale} (effet dissuasif fort), distingue \cle{tentative} et \cle{fait accompli}, intègre le \cle{devoir d'alerte} (citoyenneté active) et prévoit une \cle{dimension réparatrice} (éducative) au-delà du punitif.
\end{exoresolu}

\courssub{Intégrité et avenir professionnel technique : Éthique de l'ingénieur/technicien}

\begin{definition}[Éthique professionnelle du technicien supérieur]
L'éthique du technicien (BTP, Électrotechnique, Mécanique, Informatique, etc.) repose sur la \cle{primauté de la sécurité et de la qualité} sur les impératifs de délai, de coût ou de pression hiérarchique. Elle s'appuie sur :
\begin{itemize}
    \item Les \textbf{Normes techniques} (NF C 15-100, Eurocodes, ISO 9001, ISO 45001) : Loies techniques impératives.
    \item Le \textbf{Code du Travail} (Droit d'alerte Art. 119, Droit de retrait danger grave et imminent).
    \item Le \textbf{Code Pénal} (Responsabilité personnelle Art. 289 : mise en danger d'autrui ; Complicité Art. 41).
    \item La \cle{déontologie} propre aux Ordres professionnels (lorsqu'ils existent) : Indépendance, compétence, secret professionnel, devoir de conseil.
\end{itemize}
\end{definition}

\begin{methode}[Résoudre un dilemme éthique professionnel (Technicien/Ingénieur) : Grille de décision]
    \textbf{Objectif :} Appliquer une démarche structurée face à un conflit de valeurs (Pression hiérarchique vs Norme technique / Sécurité vs Délai / Loyauté vs Intérêt général).
    \begin{enumerate}
        \item \textbf{Énoncer le dilemme} : « Dois-je signer ce PV de réception non conforme pour respecter le délai / faire plaisir au chef ? »
        \item \textbf{Identifier les parties prenantes} : Moi (technicien, responsabilité pénale), Employeur, Client/Usagers, Collectivité, Ordre professionnel (si existant), Assureur.
        \item \textbf{Référencer les normes} : Code de déontologie (Primaire : \textbf{Sécurité personnes/biens} > Secondaire : Intérêt employeur), Normes techniques (NF, ISO), Code du travail (Droit d'alerte), Code pénal (Mise en danger d'autrui).
        \item \textbf{Évaluer les options} (Matrice) :
        \begin{center}
        \begin{tabular}{|l|c|c|}\hline
        \textbf{Option} & \textbf{Avantages courts} & \textbf{Risques \textbackslash Conséquences majeures} \\ \hline
        Signer (Céder) & Paix sociale, prime, rapidité & \textbf{Pénal (mise en danger)}, Perte licence/habilitation, Accident mortel, Remords \\ \hline
        Refuser/Signaler & Respect loi, éthique, sécurité & Pression, mutation, licenciement abusif (protégé par loi lanceur alerte) \\ \hline
        Négocier/Alternative & Solution technique conforme + délai & Exige compétence technique, courage, appui hiérarchique sain \\ \hline
        \end{tabular}
        \end{center}
        \item \textbf{Décider et Argumenter} : Choisir l'option protégeant la \cle{valeur suprême (Vie/Sécurité)}. Rédiger un \textbf{rapport écrit, daté, signé} (traçabilité) exposant le risque technique et le refus de valider.
    \end{enumerate}
\end{methode}

\begin{exoresolu}[Dilemme : Le câblage « économique »]
    \textbf{Énoncé :} Vous êtes technicien supérieur en électrotechnique (1ère TEE) en stage dans une entreprise. Le chef de chantier vous ordonne de câbler un tableau électrique tertiaire avec du câble \textbf{section 1,5 mm²} pour des circuits prises 16A (norme exige 2,5 mm² minimum), arguant : « Le client a payé pour du 1,5, le délai est demain, et ça tient le coup pour l'instant. Signe le PV de conformité. »
    \textbf{Questions :} En utilisant la grille de décision, expliquez votre conduite à tenir. Rédigez la note de service formelle que vous adressez au chef de chantier.
    \tcblower
    \textbf{Solution détaillée :}

    \textbf{1. Analyse du dilemme (Grille) :}
    \begin{itemize}
        \item \textbf{Dilemme :} Obéissance hiérarchique / Respect délai \textbf{vs} Sécurité personnes/biens / Norme NF C 15-100 / Responsabilité pénale personnelle.
        \item \textbf{Parties prenantes :} Moi (responsable exécution, stagiaire), Client (risque incendie), Entreprise (image, assurance, responsabilité civile), Usagers (danger mortel).
        \item \textbf{Normes de référence :} \textbf{NF C 15-100 (Art. 433.1.1 / Tableau 52-C)} : Section minimale cuivre pour circuit 16A en pose encastrée (type E) = \textbf{2,5 mm²}. \textbf{Code Pénal Art 289} (Mise en danger d'autrui). \textbf{Code du Travail} (Droit d'alerte / Retrait).
        \item \textbf{Choix :} \textbf{REFUS CATÉGORIQUE de signer.} La sécurité (valeur suprême) prime sur le délai et l'ordre illégal. Proposition d'alternative : Arrêt chantier le temps de livrer le 2,5 mm².
    \end{itemize}

    \textbf{2. Note de service formelle (Modèle professionnel) :}

    \begin{center}
    \textbf{ENTREPRISE [Nom] — CHANTIER [Nom/Ref]} \\
    \textbf{NOTE DE SERVICE N° 045/2025/TEE} \\
    \textbf{Objet : Refus de validation non-conformité câblage — Mise en demeure technique}
    \end{center}

    \textbf{Monsieur le Chef de Chantier,}

    \textbf{Par la présente, je vous informe formellement de mon refus de signer le Procès-Verbal de conformité du tableau électrique TGBT-01, lot Prises 16A, pour le motif suivant :}

    \textbf{Constat technique :} Le câblage réalisé utilise du câble U-1000 R2V 3G1,5 mm² pour des circuits protégés par disjoncteur 16A (Courbe C).
    \textbf{Référence normative :} Selon la \textbf{NF C 15-100 (Art. 433.1.1 / Tableau 52-C)}, la section minimale de cuivre pour un circuit 16A en pose encastrée (type E) est de \textbf{2,5 mm²}.
    \textbf{Conséquence technique :} Le 1,5 mm² entraîne une \textbf{chute de tension excessive} et un \textbf{risque d'échauffement/d'incendie} sous charge nominale (16A $\times$ 230V = 3680W), incompatible avec la sécurité des personnes et des biens (Art. 289 Code Pénal).

    \textbf{En ma qualité de technicien responsable de l'exécution (Stagiaire 1ère TEE, encadrant : [Nom Tuteur]), j'engage ma responsabilité pénale et civile en validant une installation non conforme.}

    \textbf{Demande :} Je vous demande de bien vouloir commander immédiatement le câble 2,5 mm² et de suspendre la mise sous tension de ce circuit jusqu'à mise en conformité. Je reste à disposition pour toute expertise contradictoire.

    \textbf{Fait à [Lieu], le [Date]} \\
    \textbf{Signature : [Votre Nom] — Technicien Stagiaire 1ère TEE} \\
    \textbf{Copie :} Direction Technique / Bureau de Contrôle / Tuteur Stage / Dossier Qualité.

    \textbf{Conclusion :} L'intégrité du technicien, c'est le \cle{courage technique} de dire « NON » à l'illégal, en s'appuyant sur la \cle{norme écrite} et en laissant une \cle{traçabilité écrite} qui protège l'entreprise... et la vie des usagers.
\end{exoresolu}

\courssub{Synthèse évaluative et Engagement citoyen}

\begin{methode}[Réussir une dissertation / synthèse évaluative type Probatoire/Bac (Sujet d'actualité)]
    \textbf{Objectif :} Structurer une réponse argumentée en 3 parties (Thèse / Antithèse / Synthèse) ou (Constat / Causes / Solutions) selon le libellé, en mobilisant le cours (concepts, institutions, exemples camerounais).
    \begin{enumerate}
        \item \textbf{Analyser le sujet} : Mots-clés (Intégrité, Développement, École, Citoyen), Verbe d'action (Discutez, Montrez, Comment), Définir les termes dans l'intro.
        \item \textbf{Construire le plan détaillé} (sur brouillon) :
            \begin{itemize}
                \item \textbf{Intro :} Accroche (actualité/citation/chiffre CONAC) $\rightarrow$ Définition $\rightarrow$ Problématique $\rightarrow$ Annonce du plan.
                \item \textbf{Partie I (Thèse/Constat) :} L'intégrité est un \cle{fondement} (confiance, mérite, État de droit). Ex: CONAC, TCS, Charte école.
                \item \textbf{Partie II (Antithèse/Obstacles) :} Mais elle est \cle{menacée} (corruption systémique, culture de l'impunité, pauvreté, éducation défaillante). Ex: « Arrangement », « Connaissance », détournements.
                \item \textbf{Partie III (Synthèse/Solutions) :} Elle se \cle{construit} par l'éducation (ECM), le renforcement institutionnel (indépendance justice, moyens CONAC), la transparence (numérisation), l'exemplarité des élites.
                \item \textbf{Conclusion :} Réponse à la problématique $\rightarrow$ Ouverture (Jeune technicien, acteur du changement).
            \end{itemize}
        \item \textbf{Rédiger} : Paragraphes aérés, connecteurs logiques (\emph{En effet, Cependant, Par conséquent}), exemples précis (Chiffres CONAC, Loi 2006, TD2, TP), vocabulaire ECM (\cle{probité, redevabilité, citoyenneté active}).
        \item \textbf{Relire} : Orthographe, respect consignes (nombre de pages), lisibilité.
    \end{enumerate}
\end{methode}

\begin{exoresolu}[Dissertation type Probatoire/Bac : « L'intégrité morale, condition sine qua non du développement du Cameroun »]
    \textbf{Énoncé :} « Sans intégrité morale, point de développement durable. » Discutez cette affirmation en vous appuyant sur les acquis de l'unité et la réalité camerounaise.
    \tcblower
    \textbf{Solution détaillée (Plan rédigé type copie Probatoire/Bac - Structure) :}

    \textbf{INTRODUCTION}
    \begin{itemize}
        \item \textbf{Accroche :} Le rapport CONAC 2023 estime le coût de la corruption à plus de \textbf{1000 milliards FCFA/an} (près de 10\% du PIB), freinant l'émergence à l'horizon 2035 (PND 2030-2035).
        \item \textbf{Définition :} L'intégrité morale = adhésion constante aux valeurs (probité, transparence, responsabilité) dans l'exercice de sa fonction, publique ou privée.
        \item \textbf{Problématique :} En quoi l'intégrité morale constitue-t-elle le socle indispensable (condition \emph{sine qua non}) de tout développement authentique au Cameroun ?
        \item \textbf{Annonce :} Nous montrerons d'abord qu'elle est le ciment de la confiance et de l'efficacité (I), puis que son absence gangrène l'appareil productif (II), pour enfin proposer les leviers de sa reconstruction effective (III).
    \end{itemize}

    \textbf{DÉVELOPPEMENT}

    \textbf{I. L'intégrité morale : Fondement de la confiance et de l'efficacité économique (Thèse)}
    \begin{itemize}
        \item \textbf{A. Confiance institutionnelle et attractivité.} La probité des agents publics (Douanes, Impôts, Marchés publics via CONSUPE) rassure investisseurs (IDE) et bailleurs (FMI, Banque Mondiale). \emph{Ex: Notation pays améliorée si indice corruption baisse.}
        \item \textbf{B. Mérite et capital humain.} L'intégrité à l'école (TD2 : non-tricherie) et au concours (Fonction Publique, Grandes Écoles) garantit que les postes techniques (ingénieurs, médecins, techniciens) reviennent aux \textbf{plus compétents}, pas aux plus riches. Condition de l'innovation technique nationale.
        \item \textbf{C. Préservation des ressources.} La transparence (Budget citoyen, Reddition comptes) assure que l'argent du contribuable finance \textbf{routes, hôpitaux, écoles, eau, électricité} et non des villas. \emph{Ex: Projet « Eau pour tous » réalisé vs détourné.}
    \end{itemize}

    \textbf{II. L'absence d'intégrité : Un frein structurel au développement (Antithèse / Réalité camerounaise)}
    \begin{itemize}
        \item \textbf{A. Coûts directs massifs :} Détournements (TCS : dossiers > 50M FCFA), surfacturations marchés (CONSUPE), fuites fiscales. Ressources détournées = \textbf{Investissements publics non réalisés} (ponts, lycées techniques, centres de santé).
        \item \textbf{B. Coûts indirects délétères :} \textbf{Insécurité juridique} (contrats non respectés), \textbf{Démotivation des méritants} (fuite des cerveaux), \textbf{Mauvaise qualité infrastructures} (ponts qui s'effondrent, câblage 1,5mm² vu en TP, bâtiments non parasismiques). Le « \emph{Article 15} » (débrouille-toi) tue l'esprit d'entreprise et l'innovation.
        \item \textbf{C. Érosion du contrat social.} Perte de légitimité de l'État, montée de l'incivilité, violence. Le citoyen ne paie plus l'impôt $\rightarrow$ cercle vicieux appauvrissant l'État.
    \end{itemize}

    \textbf{III. Reconstruire l'intégrité : Éducation, Institutions, Exemplarité (Synthèse / Solutions concrètes)}
    \begin{itemize}
        \item \textbf{A. Éducation (Base) :} Renforcer l'ECM (Programme 1ère Tech : TD2 diagnostic, Charte TP, Dilemmes techniques). Former le \textbf{technicien citoyen} (compétence technique + éthique professionnelle). « L'intégrité s'apprend comme un métier ».
        \item \textbf{B. Institutions (Gardiennes) :} Indépendance réelle du \textbf{TCS} et \textbf{CONAC} (budget propre, nomination au mérite, protection membres). Protection effective des \textbf{lanceurs d'alerte} (Loi 2019). Numérisation totale (E-procurement marchés publics, Paie biométrique, Télédéclaration fiscale) pour réduire la discrétion humaine.
        \item \textbf{C. Exemplarité (Socle) :} Sanction médiatique et effective des « gros poissons » (Opérations \emph{Epervier} / \emph{Sparrowhawk}). Code de conduite signé par tout agent public/élu. Rémunération décente et paiement régulier pour réduire la « petite corruption » de survie.
    \end{itemize}

    \textbf{CONCLUSION}
    \begin{itemize}
        \item \textbf{Bilan :} L'intégrité n'est pas un luxe moral, c'est une \textbf{infrastructure immatérielle} aussi vitale que le bitume ou l'électricité. Sans elle, le PND 2030-2035 reste une coquille vide.
        \item \textbf{Ouverture :} En tant que futurs techniciens supérieurs (BTP, Élec, Mécano, Info, etc.), nous sommes les \textbf{premiers garants} de la qualité technique. Refuser le câble 1,5mm² pour du 16A, refuser de signer un PV faux, signaler un détournement sur chantier, c'est déjà construire le Cameroun émergent. \textbf{L'intégrité se vit au quotidien, sur le terrain, dans le respect des normes.}
    \end{itemize}
\end{exoresolu}

\begin{methode}[Argumenter à l'oral / Soutenance de projet technique (Grand Oral / Évaluation continue)]
    \textbf{Objectif :} Présenter une prise de position éthique structurée en 3 minutes (Exposé) + 5 min (Échange), sans notes, en lien avec la spécialité technique.
    \begin{enumerate}
        \item \textbf{Accroche (30s) :} Situation vécue en stage / Actualité technique / Citation courte. « Lors de mon stage en maintenance industrielle, j'ai découvert... »
        \item \textbf{Problématique claire (1 phrase) :} « Comment concilier pression de production et impératif de sécurité normative ? »
        \item \textbf{Argumentaire structuré (2 min) :} 2 arguments forts + 1 nuance.
            \begin{itemize}
                \item \textbf{Arg 1 (Technique/Légal) :} Norme NF/ISO, Code du travail, Responsabilité pénale (Art 289 CP). « La norme n'est pas une option. »
                \item \textbf{Arg 2 (Économique/Humain) :} Coût accident $\gg$ Coût conformité. Image entreprise. Devoir de protection collègues.
                \item \textbf{Nuance :} Difficulté réelle (délais, fournisseurs). Proposition de \textbf{processus d'alerte précoce} (Fiche incident, Droit d'alerte CSE/CHSCT).
            \end{itemize}
        \item \textbf{Conclusion engageante (30s) :} « En tant que technicien, mon serment est la sécurité. Je signerai le PV conforme, ou j'écrirai le rapport de non-conformité. C'est ça, l'intégrité technique. »
        \item \textbf{Préparer les questions jury :} « Et si on vous licencie ? » $\rightarrow$ « Protection lanceur alerte (Loi 2019) / Prud'hommes / Devoir de conscience. »
    \end{enumerate}
\end{methode}

\begin{exoresolu}[Simulation Grand Oral : « L'intégrité, frein ou levier pour l'entreprise ? »]
    \textbf{Énoncé :} Vous passez l'oral de l'unité ECM (Coefficient 2). Le jury vous demande : « En entreprise, l'intégrité morale est souvent perçue comme un frein à la rentabilité à court terme. En tant que futur technicien supérieur, comment répondez-vous à cette affirmation ? » Vous avez 3 minutes.
    \tcblower
    \textbf{Solution détaillée (Script oral structuré - 3 min) :}

    \textbf{[Accroche - 20s]} « Monsieur le Président, Mesdames, Messieurs. « L'honnêteté est la meilleure des politiques, mais c'est aussi la plus rentable » (Henry Ford). En stage maintenance dans une agro-industrie, j'ai vu un chef d'équipe vouloir masquer une fuite d'ammoniac pour ne pas arrêter la chaîne. J'ai déclenché l'alarme. L'arrêt a coûté 2h de production. La fuite réparée a évité une explosion potentielle. \textbf{Cet arrêt a sauvé l'entreprise.} »

    \textbf{[Problématique - 10s]} « L'intégrité est-elle un coût ou un investissement ? Je soutiens qu'elle est le \textbf{seul levier de rentabilité durable} pour l'entreprise technique. »

    \textbf{[Argument 1 : La conformité technique EST la rentabilité - 1min]} « Premier argument, technique et juridique. Un technicien qui valide un câblage 1,5mm² pour du 16A (norme 2,5mm²) pour « aller vite » ne gagne pas du temps, il \textbf{crée une bombe à retardement}. L'incendie qui s'ensuivra coûtera : 1) Vies humaines (infini), 2) Pénal (prison dirigeant), 3) Assurance (refus indemnisation = faillite). \textbf{Respecter la norme NF C 15-100, c'est garantir la continuité d'exploitation.} L'intégrité technique = Gestion des risques = Rentabilité. »

    \textbf{[Argument 2 : La confiance, capital immatériel - 1min]} « Deuxième argument, relationnel. Nos clients (État, industriels) exigent la certification ISO 9001 / 45001. Ces normes imposent la \textbf{traçabilité, la non-conformité assumée, l'amélioration continue}. Une entreprise où l'on « arrange » les PV de réception perd sa certification, perd les marchés publics (CONSUPE), perd ses clients. \textbf{L'intégrité est un actif commercial.} Regardez les entreprises camerounaises qui exportent : ce sont celles qui ont banni la corruption interne. »

    \textbf{[Nuance / Réalisme - 30s]} « Certes, à court terme, dire « non » au chef, refuser le pot-de-vin du fournisseur, arrêter la chaîne, ça coûte du courage, du temps, parfois une prime. Mais le \textbf{Code du Travail (Art 119)} protège le salarié qui alerte sur un danger grave. Et la \textbf{Loi 2019 sur les lanceurs d'alerte} renforce ce bouclier. La solution n'est pas la soumission, c'est l'\textbf{alerte précoce écrite} (Fiche HSSE, Note de service) qui protège l'entreprise et le technicien. »

    \textbf{[Conclusion - 30s]} « Pour conclure, l'intégrité n'est pas un frein, c'est le \textbf{frein de secours} qui évite le crash. En 1ère TEE, nous apprenons à calculer des sections de câbles, à dimensionner des poutres. \textbf{Appliquer ces calculs avec honnêteté, même sous la pression, c'est notre honneur professionnel et la condition de survie de nos entreprises.} Merci. Je suis prêt pour vos questions. »

    \textbf{Analyse de la prestation :} Vocabulaire technique (NF C 15-100, ISO, Ammoniac, PV), Juridique (Code Travail, Loi 2019), Personnel (Stage), Structuré (Thèse/Antithèse/Synthèse), Temporalité respectée. \textbf{Note attendue : 18-20/20.}
\end{exoresolu}

\begin{attention}[Les 5 erreurs qui coûtent des points au Probatoire/Bac (ECM - Intégrité Morale)]
    \begin{enumerate}
        \item \textbf{Confondre les institutions (Le classique) :} Dire que la \cle{CONAC} « juge » les corrompus ou que le \cle{TCS} « sensibilise » dans les écoles. \textbf{\rightarrow Règle :} CONAC = Prévention/Enquête administrative (Rapport au PR). TCS = Jugement (Peine de prison) pour montants $\ge$ 50 millions. CONSUPE = Marchés publics uniquement. \textbf{Perte : 2 à 4 pts sur une question d'identification.}
        \item \textbf{Définition floue ou circulaire :} « L'intégrité, c'est être intègre » ou « C'est ne pas voler ». \textbf{\rightarrow Exigence :} Définition opératoire : « Conformité constante actes/paroles/valeurs, même sans témoin, impliquant probité, transparence, responsabilité. » Citer les \textbf{4 piliers (HPTR)}. \textbf{Perte : 3 pts sur la définition.}
        \item \textbf{Oublier la dimension \emph{technique/professionnelle} (Spécificité 1ère Tech) :} Traiter l'intégrité seulement comme « ne pas tricher au bac » ou « ne pas voler au marché ». \textbf{\rightarrow Exigence :} Lier à \textbf{votre spécialité} : Normes (NF, ISO, Eurocodes), Sécurité (Code Pénal Art 289), Qualité (ISO 9001), Responsabilité civile/pénale de l'ingénieur/technicien, Traçabilité écrite. \textbf{Perte : 4-5 pts sur la dissertation/oral (hors-sujet partiel).}
        \item \textbf{Sanctions « scolaires » uniquement dans les cas de corruption :} Répondre « Exclusion définitive » pour un détournement de 100 millions ou une corruption de juge. \textbf{\rightarrow Exigence :} Distinction nette \textbf{Sanction Disciplinaire (École/Entreprise)} vs \textbf{Sanction Pénale (Prison, Amende, Déchéance droits civiques, Interdiction de gérer)}. Citer les articles du Code Pénal (134 corruption, 134-1 concussion, 289 mise en danger, 184 détournement). \textbf{Perte : 3 pts sur l'analyse juridique.}
        \item \textbf{Dissertation sans exemples camerounais concrets et chiffrés :} « La corruption est mauvaise pour le développement » (Généralité). \textbf{\rightarrow Exigence :} \textbf{Chiffres CONAC (Rapports annuels)}, \textbf{Opérations Épervier/Sparrowhawk}, \textbf{Lois 2006/014, 2011/002, 2019}, \textbf{PND 2035}, \textbf{Cas concrets vus en classe (TD2, TP, Stage)}. Une dissertation sans ancrage national = \textbf{Note max 8/20.}
    \end{enumerate}
\end{attention}

\newpage

\coursec{Exercices}

\courssub{Je vérifie mes connaissances}

\begin{exercice}[QCM : les notions de base]
Pour chaque question, une seule réponse est exacte. Recopie la lettre correspondante.
\begin{enumerate}
\item L'\cle{intégrité morale} désigne :
\begin{enumerate}
\item la capacité à intégrer une entreprise publique ;
\item la qualité d'une personne qui agit avec honnêteté et droiture, même sans témoin ;
\item le respect des seules lois pénales du pays ;
\item l'ensemble des règles de politesse en société.
\end{enumerate}
\item La \cle{probité} se distingue de l'intégrité en ce qu'elle :
\begin{enumerate}
\item concerne uniquement les magistrats ;
\item désigne l'honnêteté scrupuleuse dans l'exercice d'une fonction, notamment dans la gestion des biens publics ;
\item est une vertu purement religieuse ;
\item n'existe pas dans le droit camerounais.
\end{enumerate}
\item La \cle{CONAC} (Commission Nationale Anti-Corruption) a pour mission principale :
\begin{enumerate}
\item de juger les criminels de sang ;
\item de contribuer à la lutte contre la corruption au Cameroun ;
\item de contrôler le budget de l'État en amont des dépenses ;
\item de recruter les fonctionnaires.
\end{enumerate}
\item Un \cle{conflit d'intérêts} se produit lorsque :
\begin{enumerate}
\item deux collègues se disputent un poste ;
\item un agent public prend une décision dans laquelle il a un intérêt personnel, financier ou familial ;
\item un élève refuse de faire ses devoirs ;
\item deux entreprises se font concurrence.
\end{enumerate}
\end{enumerate}
\end{exercice}

\begin{exercice}[Vrai ou faux : justifie ta réponse]
Réponds par \textbf{Vrai} ou \textbf{Faux}, puis justifie chaque réponse en une ou deux phrases.
\begin{enumerate}
\item La Chambre des comptes de la Cour suprême (qui a succédé à la CONSUPE) contrôle l'emploi des deniers publics et peut sanctionner les gestionnaires fautifs.
\item Tant que personne ne me voit, falsifier un procès-verbal de chantier n'est pas une atteinte à l'intégrité morale.
\item La tricherie à l'examen du Probatoire technique est passible de sanctions disciplinaires et peut entraîner l'exclusion de l'examen.
\item Le signalement d'un acte de corruption à la CONAC doit obligatoirement révéler l'identité du dénonciateur.
\item L'intégrité morale ne concerne que les hommes politiques et les hauts fonctionnaires.
\end{enumerate}
\end{exercice}

\begin{exercice}[Définitions express]
Définis en deux ou trois phrases chacun des termes suivants, puis illustre chacun par un exemple tiré de la vie scolaire ou professionnelle :
\begin{enumerate}
\item l'intégrité morale ;
\item la corruption ;
\item le détournement de fonds publics ;
\item le népotisme ;
\item le serment professionnel (d'un ingénieur ou d'un technicien).
\end{enumerate}
\end{exercice}

\begin{exercice}[Institutions : qui fait quoi ?]
Recopie et complète le tableau ci-dessous en reliant chaque institution à sa mission principale et en citant un exemple concret d'intervention.

\begin{center}
\begin{tabularx}{\linewidth}{|l|X|X|}\hline
\rowcolor{popblueL} \textbf{Institution} & \textbf{Mission principale} & \textbf{Exemple d'intervention} \\ \hline
CONAC & & \\ \hline
Chambre des comptes de la Cour suprême & & \\ \hline
Agence de Régulation des Marchés Publics (ARMP) & & \\ \hline
Tribunaux répressifs (juridictions pénales) & & \\ \hline
\end{tabularx}
\end{center}
\end{exercice}

\courssub{Je m'entraîne}

\begin{exercice}[Analyse de situations : intégrité ou non ?]
Pour chacune des situations suivantes, indique s'il y a atteinte à l'intégrité morale, nomme la faute commise (corruption, conflit d'intérêts, négligence, fraude, favoritisme...) et propose le comportement intègre à adopter.
\begin{enumerate}
\item Un surveillant d'examen ferme les yeux sur la tricherie du fils de son voisin.
\item Un contrôleur de chantier accepte une enveloppe de \num{50000} FCFA pour valider une livraison de ciment de mauvaise qualité.
\item Un élève refuse de recopier le devoir de son camarade et lui propose plutôt de lui réexpliquer la leçon.
\item Un maire attribue le marché de réfection de la salle des fêtes à l'entreprise de son beau-frère, sans appel d'offres.
\item Une technicienne de laboratoire constate une erreur dans ses mesures et recommence l'analyse, bien qu'elle soit pressée.
\end{enumerate}
\end{exercice}

\begin{exercice}[Le courrier de refus argumenté]
Tu es élève en classe de Première technique, délégué de classe. Le président de la coopérative scolaire te propose de gonfler les factures d'achat de fournitures pour « arrondir les fins de mois » du bureau. Rédige un court texte argumenté (10 à 15 lignes) dans lequel tu refuses, en mobilisant : la définition de l'intégrité morale, les risques disciplinaires et pénaux, et les conséquences pour la confiance au sein de la coopérative. Termine par une proposition honnête pour améliorer les finances de la coopérative.
\end{exercice}

\begin{exercice}[Check-list intégrité en atelier]
En atelier de génie civil, ton groupe doit réceptionner une livraison de matériaux (ciment, ferraillage, granulats). Élabore une check-list de réception comportant :
\begin{enumerate}
\item cinq points de contrôle techniques (quantité, qualité, conformité des documents de livraison...) ;
\item trois points de contrôle éthiques (absence de conflit d'intérêts, refus de tout avantage, signature à deux personnes...).
\end{enumerate}
Pour chaque point, justifie en une phrase le risque évité (négligence, corruption, accident).
\end{exercice}

\begin{exercice}[Le dilemme du logiciel piraté]
En stage dans une PME, ton maître de stage te demande d'installer un logiciel de conception « cracké » pour économiser le coût des licences.
\begin{enumerate}
\item Identifie les valeurs morales et les règles juridiques mises en cause (propriété intellectuelle, honnêteté, responsabilité).
\item Énumère trois risques concrets pour l'entreprise (sécurité informatique, sanctions, image).
\item Propose deux alternatives légales et rédige en cinq lignes la réponse respectueuse mais ferme que tu adresses à ton maître de stage.
\end{enumerate}
\end{exercice}

\courssub{Je me prépare au Bac}

\begin{exercice}[Sujet type Bac : corruption et marchés publics]
\textbf{Dossier documentaire.}

\textbf{Document 1} — Extrait adapté d'un rapport de la CONAC : « Dans plusieurs marchés publics de travaux, des surfacturations et des attributions irrégulières ont été constatées, privant l'État de ressources destinées aux infrastructures scolaires et sanitaires. »

\textbf{Document 2} — Article de presse : « Effondrement d'un immeuble à Douala : l'enquête révèle l'utilisation de ciment frelaté et un contrôle technique défaillant. Plusieurs victimes sont à déplorer. »

\textbf{Document 3} — Organigramme simplifié : CONAC (lutte contre la corruption) ; Chambre des comptes de la Cour suprême (contrôle des deniers publics) ; ARMP (régulation des marchés publics) ; Tribunaux (sanctions pénales).

\textbf{Questions :}
\begin{enumerate}
\item À partir des documents 1 et 2, identifie et nomme trois atteintes à l'intégrité morale présentées. (3 pts)
\item Définis l'intégrité morale et montre en quoi les faits du document 2 la violent doublement (fraude et négligence). (4 pts)
\item Explique le rôle de deux institutions du document 3 dans la prévention ou la sanction de ces faits. (4 pts)
\item En tant que futur technicien supérieur, propose trois mesures concrètes pour garantir l'intégrité sur un chantier de construction. (3 pts)
\item Rédige une lettre de signalement à la CONAC dénonçant les faits du document 1 (en-tête, objet, faits, demande, formule de politesse). (6 pts)
\end{enumerate}
\end{exercice}

\begin{exercice}[Sujet type Bac : l'intégrité en milieu scolaire]
\textbf{Situation.} Dans un lycée technique de Yaoundé, à la veille du Probatoire, des rumeurs font état de la vente de sujets d'épreuves. Par ailleurs, certains élèves copient systématiquement leurs devoirs sur Internet sans citer leurs sources, et un comité d'examen a constaté des falsifications de relevés de notes.

\textbf{Questions :}
\begin{enumerate}
\item Définis l'intégrité morale et montre pourquoi elle est un pilier de la vie scolaire. (3 pts)
\item Identifie dans la situation trois comportements contraires à l'intégrité et classe-les selon leur gravité, en justifiant ton classement. (4 pts)
\item Explique deux conséquences de la fraude aux examens : une pour l'élève, une pour la société (valeur du diplôme, sécurité des futurs ouvrages). (4 pts)
\item Cite deux institutions ou instances (nationales ou scolaires) susceptibles de sanctionner ces faits et précise leur rôle. (3 pts)
\item Rédige le préambule (8 à 12 lignes) d'une charte d'intégrité morale de ta classe, comportant trois engagements précis et les sanctions prévues. (6 pts)
\end{enumerate}
\end{exercice}

\begin{exercice}[Sujet type Bac : intégrité et avenir professionnel du technicien]
\textbf{Situation.} Recruté comme technicien de chantier, tu constates que l'entreprise utilise des barres d'acier de diamètre inférieur à celui prévu par les plans. Le chef de chantier te répond : « Signe le procès-verbal de conformité, c'est une pratique courante, et tu recevras une prime. »

\textbf{Questions :}
\begin{enumerate}
\item Identifie les pressions morales et matérielles exercées sur le technicien et nomme la faute qu'on lui demande de commettre. (3 pts)
\item Explique pourquoi la signature d'un faux procès-verbal engage la responsabilité personnelle du technicien (responsabilité professionnelle et pénale). (4 pts)
\item Décris la démarche intègre à suivre, étape par étape (refus, constat écrit, information de la hiérarchie, signalement si nécessaire). (4 pts)
\item Cite deux institutions camerounaises auxquelles ces faits pourraient être signalés et précise leur mission. (3 pts)
\item Rédige un texte argumentatif de 15 à 20 lignes sur le sujet : « L'intégrité morale du technicien, garante de la sécurité des personnes et de la qualité des ouvrages », en mobilisant au moins trois notions de la leçon. (6 pts)
\end{enumerate}
\end{exercice}

\newpage

\coursec{Corrigés des exercices}

\courssub{Corrigés détaillés des exercices et sujets type Bac}

\begin{corrige}[Exercice 1 — QCM : les notions de base]
\begin{enumerate}
\item \textbf{b} : l'intégrité morale est la qualité de la personne qui agit avec honnêteté et droiture en toutes circonstances, même sans témoin. Les réponses a, c et d la réduisent à tort à un statut, à la seule loi pénale ou à la politesse.
\item \textbf{b} : la probité est l'honnêteté scrupuleuse dans l'exercice d'une fonction, en particulier dans la gestion des biens et deniers publics.
\item \textbf{b} : la CONAC contribue à la lutte contre la corruption ; elle ne juge pas (rôle des tribunaux) et ne contrôle pas le budget en amont.
\item \textbf{b} : il y a conflit d'intérêts lorsqu'un agent prend une décision publique dans laquelle il a un intérêt personnel, financier ou familial.
\end{enumerate}
\end{corrige}

\begin{corrige}[Exercice 2 — Vrai ou faux : justifie ta réponse]
\begin{enumerate}
\item \textbf{Vrai} : la Chambre des comptes de la Cour suprême contrôle l'emploi des deniers publics, juge les comptes des gestionnaires publics et peut les sanctionner en cas de faute de gestion.
\item \textbf{Faux} : l'intégrité morale commande d'agir avec droiture même sans témoin ; falsifier un procès-verbal est un faux, une faute professionnelle et pénale, que l'on soit vu ou non.
\item \textbf{Vrai} : la fraude aux examens officiels est sanctionnée disciplinairement (annulation de l'épreuve, exclusion de l'examen) et peut donner lieu à des poursuites.
\item \textbf{Faux} : le signalement peut se faire de manière confidentielle ; la protection du dénonciateur est recherchée et son identité n'a pas à être révélée publiquement.
\item \textbf{Faux} : l'intégrité morale concerne tout citoyen : élève, enseignant, commerçant, technicien, fonctionnaire, et non seulement les responsables politiques.
\end{enumerate}
\textbf{Point de vigilance} : une réponse « Vrai » ou « Faux » non justifiée ne rapporte aucun point ; la justification doit mobiliser une notion du cours (intégrité, faux, fraude, confidentialité).
\end{corrige}

\begin{corrige}[Exercice 3 — Définitions express]
\begin{enumerate}
\item \textbf{L'intégrité morale} : qualité d'une personne qui agit avec honnêteté, droiture et loyauté en toutes circonstances, même lorsque personne ne la regarde. Elle repose sur la cohérence entre les valeurs professées et les actes posés. \emph{Exemple} : un élève rend à la vie scolaire l'argent trouvé dans la cour au lieu de le garder.
\item \textbf{La corruption} : fait d'offrir, de promettre, de solliciter ou d'accepter un avantage (argent, cadeau, service) pour qu'une personne accomplisse ou s'abstienne d'un acte de sa fonction. \emph{Exemple} : verser de l'argent à un surveillant pour obtenir un sujet d'examen à l'avance.
\item \textbf{Le détournement de fonds publics} : fait, pour un gestionnaire ou un agent, d'utiliser à des fins personnelles des biens ou des deniers publics confiés à sa garde. \emph{Exemple} : le trésorier d'une coopérative scolaire utilise la caisse pour payer ses dépenses personnelles.
\item \textbf{Le népotisme} : favoritisme consistant à accorder un avantage (poste, marché, note) à un parent ou à un proche, au mépris des règles et du mérite. \emph{Exemple} : attribuer un marché de travaux à l'entreprise de son beau-frère sans appel d'offres.
\item \textbf{Le serment professionnel} : engagement solennel pris par un ingénieur ou un technicien d'exercer son métier avec compétence, honnêteté et dans le respect de la sécurité des personnes et de l'intérêt général. \emph{Exemple} : un technicien de chantier refuse de certifier conforme un ouvrage non conforme aux plans.
\end{enumerate}
\end{corrige}

\begin{corrige}[Exercice 4 — Institutions : qui fait quoi ?]
\begin{center}
\begin{tabularx}{\linewidth}{|l|X|X|}\hline
\rowcolor{popblueL} \textbf{Institution} & \textbf{Mission principale} & \textbf{Exemple d'intervention} \\ \hline
CONAC & Contribuer à la lutte contre la corruption : enquêtes, rapports annuels, sensibilisation. & Enquête sur des surfacturations dans un marché de travaux et transmission du dossier à la justice. \\ \hline
Chambre des comptes de la Cour suprême & Contrôler l'emploi des deniers publics et juger les comptes des gestionnaires publics. & Vérification de la gestion d'un lycée technique et sanction d'un gestionnaire fautif. \\ \hline
ARMP & Réguler la passation et l'exécution des marchés publics (transparence, concurrence). & Contrôle d'un appel d'offres irrégulier et recommandation d'annulation. \\ \hline
Tribunaux répressifs & Sanctionner pénalement les infractions (corruption, détournement, faux). & Condamnation d'un agent corrompu à une peine de prison et à une amende. \\ \hline
\end{tabularx}
\end{center}
\textbf{Point de vigilance} : ne pas confondre les rôles — la CONAC enquête et sensibilise mais ne juge pas ; seules les juridictions pénales condamnent.
\end{corrige}

\begin{corrige}[Exercice 5 — Analyse de situations : intégrité ou non ?]
\begin{enumerate}
\item \textbf{Atteinte} : favoritisme et complicité de fraude. Comportement intègre : appliquer le règlement à tous les candidats sans distinction et dresser un procès-verbal de fraude.
\item \textbf{Atteinte} : corruption (acceptation d'un avantage contre un acte de sa fonction). Comportement intègre : refuser l'enveloppe, refuser la livraison non conforme et rédiger un rapport de non-conformité.
\item \textbf{Pas d'atteinte} : comportement intègre et solidaire ; l'élève refuse la fraude tout en aidant honnêtement son camarade.
\item \textbf{Atteinte} : conflit d'intérêts et favoritisme (attribution irrégulière de marché). Comportement intègre : organiser un appel d'offres transparent et se déporter de la décision.
\item \textbf{Pas d'atteinte} : comportement intègre ; la technicienne privilégie l'exactitude des résultats, garante de la qualité et de la sécurité.
\end{enumerate}
\textbf{Point de vigilance} : la réponse complète exige trois éléments — le verdict (atteinte ou non), le nom précis de la faute, et le comportement intègre alternatif.
\end{corrige}

\begin{corrige}[Exercice 6 — Le courrier de refus argumenté]
\textbf{Éléments de correction (texte attendu en 10 à 15 lignes).}

Monsieur le Président, j'ai bien réfléchi à votre proposition de gonfler les factures d'achat de fournitures, et je dois vous signifier mon refus. L'intégrité morale est la qualité de la personne qui agit avec honnêteté et droiture, même lorsque personne ne la regarde : gonfler une facture, c'est établir un faux document et détourner l'argent de la coopérative, qui appartient à tous les élèves. Sur le plan disciplinaire, une telle pratique expose le bureau à des sanctions du conseil de discipline, pouvant aller jusqu'à l'exclusion ; sur le plan pénal, l'établissement de fausses factures et le détournement de fonds sont des infractions punies par la loi. Enfin, si cette pratique était découverte, la confiance des élèves, des enseignants et des partenaires envers la coopérative serait détruite, et aucun projet ne pourrait plus être financé. Je vous propose plutôt d'augmenter honnêtement nos ressources : organiser une kermesse, négocier des remises auprès des fournisseurs, ou lancer une campagne de cotisations volontaires. Comptant sur votre compréhension, je vous prie d'agréer, Monsieur le Président, mes salutations distinguées.

\textbf{Points de vigilance} : la définition de l'intégrité doit être explicite ; les risques disciplinaires ET pénaux doivent être distingués ; la proposition finale doit être concrète et honnête.
\end{corrige}

\begin{corrige}[Exercice 7 — Check-list intégrité en atelier]
\textbf{Cinq points de contrôle techniques :}
\begin{enumerate}
\item Vérifier les quantités livrées par rapport au bon de commande (comptage, pesée) : évite la surfacturation et les pénuries.
\item Contrôler la qualité du ciment (date de fabrication, absence de grumeaux, marquage normatif) : évite les ouvrages fragiles et les accidents.
\item Vérifier la conformité du ferraillage (diamètres, nuances d'acier prévus par les plans) : évite les défauts de résistance des structures.
\item Exiger et archiver les documents de livraison (bon de livraison, certificats de conformité) : évite la fraude documentaire et facilite les recours.
\item Réaliser des prélèvements d'échantillons de granulats pour essais : évite la négligence sur la qualité des matériaux.
\end{enumerate}
\textbf{Trois points de contrôle éthiques :}
\begin{enumerate}
\item Déclarer tout lien personnel avec le fournisseur et se déporter le cas échéant : évite le conflit d'intérêts.
\item Refuser tout avantage (argent, cadeau, invitation) offert par le livreur : évite la corruption.
\item Faire signer le procès-verbal de réception par deux personnes au moins : évite la falsification et répartit la responsabilité.
\end{enumerate}
\textbf{Point de vigilance} : chaque point doit être accompagné de sa justification (risque évité) ; une liste sans justification est incomplète.
\end{corrige}

\begin{corrige}[Exercice 8 — Le dilemme du logiciel piraté]
\begin{enumerate}
\item \textbf{Valeurs et règles mises en cause} : l'honnêteté et la loyauté (valeurs morales) ; le respect de la propriété intellectuelle et du droit d'auteur, car utiliser un logiciel « cracké » est une contrefaçon, une infraction pénale ; la responsabilité professionnelle du technicien, qui ne doit pas participer à un acte illégal.
\item \textbf{Trois risques concrets pour l'entreprise} : un risque de sécurité informatique (les logiciels crackés contiennent souvent des virus ou des programmes espions qui menacent les données) ; un risque de sanctions (amendes et poursuites en cas de contrôle) ; un risque d'image (perte de confiance des clients et des partenaires si la fraude est connue).
\item \textbf{Alternatives légales} : utiliser un logiciel libre et gratuit équivalent ; souscrire une licence éducation ou une licence adaptée au budget de la PME.

\textbf{Réponse au maître de stage (5 lignes)} : « Monsieur, je comprends le souci d'économie de l'entreprise, mais je ne peux pas installer ce logiciel cracké : c'est une contrefaçon qui expose la société à des virus, à des amendes et à une atteinte à son image. Je vous propose d'installer un logiciel libre équivalent ou d'étudier une licence à tarif réduit. Je reste à votre disposition pour mettre en place cette solution légale. »
\end{enumerate}
\textbf{Point de vigilance} : le ton doit rester respectueux mais ferme ; le refus doit toujours être accompagné d'une alternative concrète.
\end{corrige}

\begin{corrige}[Exercice 9 — Sujet type Bac : corruption et marchés publics]
\textbf{Barème indicatif : 20 points.}

\begin{enumerate}
\item \textbf{Trois atteintes (3 pts, 1 pt par atteinte nommée)} : la surfacturation (facturation excessive de travaux) ; l'attribution irrégulière de marchés (favoritisme, absence d'appel d'offres) ; la fraude sur la qualité des matériaux (ciment frelaté) assortie de négligence dans le contrôle technique.
\item \textbf{Définition et double violation (4 pts)} : l'intégrité morale est la qualité de la personne qui agit avec honnêteté et droiture en toutes circonstances, même sans témoin (2 pts). Le document 2 la viole doublement : d'une part par la \emph{fraude} (livraison délibérée de ciment frelaté, contraire à l'honnêteté), d'autre part par la \emph{négligence} (contrôle technique défaillant, manquement au devoir de diligence), avec des conséquences mortelles (2 pts).
\item \textbf{Rôle de deux institutions (4 pts, 2 pts chacune)} : la CONAC enquête sur les faits de corruption, établit des rapports et saisit la justice ; l'ARMP veille à la transparence et à la concurrence dans la passation des marchés publics et peut recommander l'annulation d'une attribution irrégulière. (Autres réponses possibles : Chambre des comptes, qui contrôle l'emploi des deniers publics ; tribunaux, qui sanctionnent pénalement.)
\item \textbf{Trois mesures concrètes (3 pts, 1 pt chacune)} : exiger des certificats de conformité et faire des essais sur les matériaux avant usage ; tenir un cahier de chantier exact et signer les procès-verbaux à deux personnes ; refuser tout avantage et signaler toute pression ou irrégularité à la hiérarchie ou à la CONAC.
\item \textbf{Lettre de signalement (6 pts)} : en-tête (nom, adresse du signataire ; destinataire : Président de la CONAC, Yaoundé ; date et lieu) : 1,5 pt ; objet précis (« Signalement de surfacturations et d'attributions irrégulières de marchés publics ») : 1 pt ; exposé factuel et daté des faits : 1,5 pt ; demande d'enquête et mention de la confidentialité souhaitée : 1 pt ; formule de politesse et signature : 1 pt.
\end{enumerate}
\textbf{Points de vigilance} : distinguer fraude (acte délibéré) et négligence (manquement au devoir de vigilance) ; rappeler que la CONAC enquête mais ne juge pas ; la lettre doit rester factuelle, sans insultes ni accusations non étayées.
\end{corrige}

\begin{corrige}[Exercice 10 — Sujet type Bac : l'intégrité en milieu scolaire]
\textbf{Barème indicatif : 20 points.}

\begin{enumerate}
\item \textbf{Définition et pilier de la vie scolaire (3 pts)} : l'intégrité morale est la qualité de la personne qui agit avec honnêteté et droiture en toutes circonstances, même sans témoin (1,5 pt). Elle est un pilier de la vie scolaire car elle garantit la valeur des évaluations, l'équité entre les élèves et la confiance entre élèves, enseignants et administration (1,5 pt).
\item \textbf{Trois comportements et classement (4 pts)} : la vente de sujets d'examen (fraude aggravée, la plus grave, car elle ruine l'équité d'un examen national et engage la responsabilité pénale) ; la falsification de relevés de notes (faux, très grave, car elle trompe l'institution) ; le plagiat de devoirs sur Internet sans citation des sources (atteinte à l'honnêteté intellectuelle, grave mais de moindre portée). Le classement doit être justifié par l'étendue des conséquences (1 pt par comportement identifié, 1 pt pour la justification du classement).
\item \textbf{Deux conséquences (4 pts)} : pour l'élève, annulation de l'épreuve, exclusion de l'examen, voire poursuites, et formation non acquise (2 pts) ; pour la société, dévalorisation du diplôme et risque de techniciens incompétents construisant des ouvrages dangereux (2 pts).
\item \textbf{Deux institutions (3 pts)} : l'OBC (Office du Baccalauréat du Cameroun) et les jurys d'examen, qui constatent et sanctionnent les fraudes ; le conseil de discipline du lycée, qui sanctionne les fautes scolaires ; en cas de faux, les tribunaux peuvent être saisis (1,5 pt par institution avec son rôle précisé).
\item \textbf{Préambule de charte (6 pts)} : rappel des valeurs (honnêteté, équité, respect) : 2 pts ; trois engagements précis (ne pas tricher, ne pas plagier, signaler toute fraude) : 2 pts ; sanctions prévues (avertissement, rapport au conseil de discipline) : 2 pts.
\end{enumerate}
\textbf{Points de vigilance} : le classement par gravité doit être argumenté, pas seulement énoncé ; les engagements de la charte doivent être précis et mesurables, non de simples formules vagues.
\end{corrige}

\begin{corrige}[Exercice 11 — Sujet type Bac : intégrité et avenir professionnel du technicien]
\textbf{Barème indicatif : 20 points.}

\begin{enumerate}
\item \textbf{Pressions et faute (3 pts)} : pression morale (« c'est une pratique courante », appel à la complicité et à la banalisation du mal) et pression matérielle (la prime promise, forme de corruption) (1,5 pt). La faute demandée est l'établissement et la signature d'un faux procès-verbal de conformité, c'est-à-dire un faux, assorti d'une acceptation d'avantage (corruption passive) (1,5 pt).
\item \textbf{Responsabilité personnelle (4 pts)} : en signant, le technicien atteste personnellement d'un fait qu'il sait inexact ; il commet un faux, infraction pénale, et engage sa responsabilité professionnelle (sanctions disciplinaires, perte d'emploi, interdiction d'exercer) et pénale (poursuites, amende, emprisonnement). En cas d'effondrement de l'ouvrage, sa responsabilité peut être recherchée pour les dommages causés aux personnes.
\item \textbf{Démarche intègre étape par étape (4 pts)} : (1) refuser clairement de signer et refuser la prime ; (2) dresser un constat écrit et mesuré (diamètres réels relevés, photos, témoins) ; (3) informer par écrit la hiérarchie ou le maître d'ouvrage ; (4) si rien ne change, signaler les faits aux institutions compétentes en demandant la confidentialité.
\item \textbf{Deux institutions (3 pts)} : la CONAC, qui reçoit les signalements de corruption et enquête ; les tribunaux répressifs, qui sanctionnent le faux et la corruption ; on peut aussi citer l'ARMP si le chantier relève d'un marché public (1,5 pt par institution avec sa mission).
\item \textbf{Texte argumentatif (6 pts)} — \textbf{Éléments attendus} : introduction posant le rôle central du technicien dans la chaîne de sécurité ; développement mobilisant au moins trois notions de la leçon, par exemple l'\cle{intégrité morale} (agir avec droiture même sous pression), la \cle{probité} (refus de tout avantage), le \cle{faux} (le procès-verbal signé à contrecœur n'efface pas la faute, il l'aggrave), la \cle{responsabilité} professionnelle et pénale ; exemples concrets (effondrement d'ouvrages dus à des matériaux non conformes) ; conclusion montrant que l'intégrité protège à la fois les usagers, la réputation du technicien et la qualité des ouvrages. Barème : respect de la longueur et de la structure argumentative : 2 pts ; mobilisation correcte d'au moins trois notions : 2 pts ; exemples et cohérence : 2 pts.
\end{enumerate}
\textbf{Points de vigilance} : ne pas écrire que « l'entreprise seule est responsable » : le signataire d'un procès-verbal engage sa responsabilité personnelle ; la démarche intègre doit laisser des traces écrites à chaque étape ; le signalement à la CONAC peut demander la confidentialité.
\end{corrige}

\newpage

\coursec{Fiche bilan}

\begin{popfigure}
\begin{tikzpicture}[mindmap, grow cyclic, every node/.style={concept}, root concept/.append style={concept color=popblue, font=\bfseries, minimum size=3.2cm, text=white}, level 1/.append style={level distance=4.2cm, sibling angle=51.4}, level 1 concept/.append style={font=\small\bfseries, minimum size=2.6cm, text=white}]
\node[root concept]{L'intégrité morale}
  child[concept color=popgreen]{ node[align=center]{Concept\\ honnêteté, droiture, cohérence} }
  child[concept color=poporange]{ node[align=center]{Piliers\\ vérité, justice, respect, responsabilité} }
  child[concept color=poppurple]{ node[align=center]{Enjeux\\ confiance, paix, développement} }
  child[concept color=popgold]{ node[align=center]{Institutions\\ CONAC, justice, Églises, école, famille} }
  child[concept color=poppink]{ node[align=center]{Milieu scolaire\\ triche, vol, corruption : diagnostic} }
  child[concept color=popblue!70]{ node[align=center]{Outils de promotion\\ chartes, croquis, sensibilisation} }
  child[concept color=popgreen!70!black]{ node[align=center]{Éthique professionnelle\\ technicien intègre} };
\end{tikzpicture}
\legende{Carte mentale de la leçon : les sept idées forces autour de l'intégrité morale (concept, piliers, enjeux, institutions, milieu scolaire, outils de promotion, éthique professionnelle).}
\end{popfigure}

\begin{aretenir}[L'essentiel de la leçon]
\textbf{Définitions clés}
\begin{itemize}
\item \cle{Intégrité morale} : qualité d'une personne qui agit avec honnêteté, droiture et cohérence entre ses paroles, ses valeurs et ses actes, même lorsque personne ne la surveille.
\item \cle{Corruption} : fait d'offrir, de promettre, de solliciter ou d'accepter un avantage indu pour accomplir ou s'abstenir d'accomplir un acte de sa fonction.
\item \cle{Probité} : scrupuleuse honnêteté dans l'exercice d'une charge ou d'une fonction.
\item \cle{Triche} : violation délibérée d'une règle (examen, concours, jeu) pour obtenir un avantage indu.
\end{itemize}

\textbf{Repères à connaître par cœur}
\begin{center}
\begin{tabularx}{\linewidth}{|l|X|}\hline
\rowcolor{popblueL} \textbf{Notion} & \textbf{À retenir} \\ \hline
Piliers de l'intégrité & honnêteté, vérité, justice, respect, responsabilité, courage moral \\ \hline
Enjeux & confiance sociale, paix, bonne gouvernance, développement économique \\ \hline
Institutions de promotion & famille, école, Églises et chefferies traditionnelles, médias \\ \hline
Institutions de préservation & CONAC, justice (tribunaux), ministère de la Fonction publique, forces de maintien de l'ordre \\ \hline
CONAC & Commission Nationale Anti-Corruption, créée par la loi n° 2006/004 du 14 juillet 2006 ; sa mission : contribuer à la lutte contre la corruption (dénonciation, prévention, répression) \\ \hline
Textes de référence & Constitution du 18 février 1996 (préambule : lutte contre la corruption) et Code pénal camerounais (répression de la corruption et des infractions assimilées) \\ \hline
\end{tabularx}
\end{center}

\textbf{Méthodes express}
\begin{itemize}
\item \emph{Étude de cas} : identifier les faits $\rightarrow$ nommer la valeur bafouée $\rightarrow$ citer l'institution compétente $\rightarrow$ proposer une solution concrète.
\item \emph{Diagnostic scolaire} : observer les pratiques (triche, vol, favoritisme), mesurer leurs causes, proposer des remèdes réalistes.
\item \emph{Construction d'une charte} : règles courtes, formulées positivement, signées par tous, avec sanctions équitables.
\item \emph{Rédaction au Probatoire} : introduction (définition + problématique), développement en 2 ou 3 arguments illustrés d'exemples camerounais, conclusion engagée.
\end{itemize}
\end{aretenir}

\begin{aretenir}[Checklist avant le Probatoire / Baccalauréat technique]
\begin{itemize}
\item[$\square$] Je sais définir l'intégrité morale en une phrase complète.
\item[$\square$] Je distingue intégrité morale, probité et honnêteté.
\item[$\square$] Je cite au moins quatre piliers de l'intégrité morale.
\item[$\square$] J'explique deux enjeux de l'intégrité pour la société camerounaise.
\item[$\square$] Je cite la CONAC et sa mission principale.
\item[$\square$] Je distingue les institutions de promotion (famille, école, Églises) des institutions de préservation (justice, CONAC).
\item[$\square$] Je donne trois exemples d'atteintes à l'intégrité en milieu scolaire et leurs conséquences.
\item[$\square$] Je sais proposer des solutions concrètes contre la triche à l'école.
\item[$\square$] Je relie l'intégrité morale à l'éthique du technicien (qualité du travail, sécurité, refus des pots-de-vin).
\item[$\square$] Je sais construire une charte ou un outil de sensibilisation simple.
\item[$\square$] Je rédige une conclusion engagée avec un appel à l'action citoyenne.
\item[$\square$] Je cite un texte officiel (Constitution de 1996 ou Code pénal) pour appuyer mon propos.
\end{itemize}
\end{aretenir}

\findecours

\end{document}
